\newcommand{\langue}{french}
\newcommand{\nomchapitre}{Chapitre}
\newcommand{\titrecourt}{MANUEL DE RÉFÉRENCE}
\input{version.tex}
% ================================================================
%  Préambule du manuel de Coupole (FR et EN), d'après le format
%  maison : Latin Modern, titres sans-serif colorés, encadrés
%  tcolorbox, fontawesome5.  \langue{french} ou \langue{english}
%  est défini par le fichier principal avant \input.
% ================================================================
\documentclass[11pt,a4paper]{article}

\usepackage[T1]{fontenc}
\usepackage[utf8]{inputenc}
\usepackage[\langue]{babel}
\usepackage{lmodern}
\usepackage{textcomp}
\usepackage{amsmath,amssymb}
\usepackage{mathtools}
\usepackage{siunitx}
\sisetup{group-minimum-digits=4,separate-uncertainty=true,per-mode=symbol}
\usepackage{cancel}
\usepackage{microtype}
\DeclareUnicodeCharacter{2248}{\ensuremath{\approx}}
\DeclareUnicodeCharacter{03B4}{\ensuremath{\delta}}
\DeclareUnicodeCharacter{03C6}{\ensuremath{\varphi}}
\DeclareUnicodeCharacter{03B8}{\ensuremath{\theta}}
\DeclareUnicodeCharacter{2295}{\ensuremath{\oplus}}
\DeclareUnicodeCharacter{2081}{\ensuremath{_1}}
\DeclareUnicodeCharacter{2082}{\ensuremath{_2}}
\DeclareUnicodeCharacter{03C0}{\ensuremath{\pi}}
\DeclareUnicodeCharacter{2212}{\ensuremath{-}}
\DeclareUnicodeCharacter{00D7}{\ensuremath{\times}}
\DeclareUnicodeCharacter{2192}{\ensuremath{\rightarrow}}
\DeclareUnicodeCharacter{00A0}{~}
\DeclareUnicodeCharacter{202F}{\,}
\DeclareUnicodeCharacter{2033}{\ensuremath{''}}
\DeclareUnicodeCharacter{2032}{\ensuremath{'}}
\DeclareUnicodeCharacter{2264}{\ensuremath{\le}}
\DeclareUnicodeCharacter{2265}{\ensuremath{\ge}}
\DeclareUnicodeCharacter{207B}{\ensuremath{^{-}}}
\DeclareUnicodeCharacter{2070}{\ensuremath{^{0}}}
\DeclareUnicodeCharacter{2713}{\ensuremath{\checkmark}}
\DeclareUnicodeCharacter{2717}{\ensuremath{\times}}
\DeclareUnicodeCharacter{2194}{\ensuremath{\leftrightarrow}}
\DeclareUnicodeCharacter{00B2}{\textsuperscript{2}}
\DeclareUnicodeCharacter{2026}{\ldots}
\DeclareUnicodeCharacter{00B0}{\textdegree}
\DeclareUnicodeCharacter{2080}{\ensuremath{_0}}
\DeclareUnicodeCharacter{03A9}{\ensuremath{\Omega}}
\DeclareUnicodeCharacter{039B}{\ensuremath{\Lambda}}
\DeclareUnicodeCharacter{03C3}{\ensuremath{\sigma}}
\DeclareUnicodeCharacter{00B3}{\textsuperscript{3}}
\DeclareUnicodeCharacter{03B1}{\ensuremath{\alpha}}
\DeclareUnicodeCharacter{03B2}{\ensuremath{\beta}}
\DeclareUnicodeCharacter{03BB}{\ensuremath{\lambda}}
\DeclareUnicodeCharacter{0153}{\oe}

\usepackage[a4paper,margin=2cm,headheight=24pt]{geometry}
\usepackage{graphicx}
\graphicspath{{img/}}
\usepackage[table]{xcolor}
\usepackage[most]{tcolorbox}
\usepackage{booktabs}
\usepackage{tabularx}
\usepackage{array}
\usepackage{multicol}
\usepackage{enumitem}
\usepackage[nobottomtitles*]{titlesec}
\usepackage{needspace}
\usepackage{fancyhdr}
\usepackage{fontawesome5}
\usepackage{tikz}
\usetikzlibrary{arrows.meta,calc,decorations.pathreplacing,patterns,angles,quotes}
\usepackage{pgfplots}
\pgfplotsset{compat=1.18}
\usepgfplotslibrary{fillbetween}
\usepackage{pdfpages}
\usepackage[colorlinks=true,linkcolor=sky,urlcolor=sky,bookmarksopen=true,bookmarksopenlevel=1,
  pdftitle={Coupole},
  pdfauthor={ARP273-ROSE}]{hyperref}
\usepackage{bookmark}
% Les notes posées dans des tableaux (annexes) par \footnotemark/\footnotetext cassent les ancres
% de hyperref : les appels de note ne sont donc pas des liens.
\hypersetup{hyperfootnotes=false}
\usepackage{xurl}
\usepackage{fvextra}
\fvset{breaklines=true,fontsize=\small}

% ---------------------------------------------------------------- Langue des titres d'encadrés
\newcommand{\tr}[2]{\ifdefstring{\langue}{english}{#2}{#1}}

% ---------------------------------------------------------------- Couleurs
\definecolor{sky}{HTML}{1C6DB4}
\definecolor{skyl}{HTML}{E7F0F9}
\definecolor{alert}{HTML}{B5382B}
\definecolor{alertl}{HTML}{FBEAE7}
\definecolor{night}{HTML}{1B2A4A}
\definecolor{ink}{HTML}{1F2430}
\definecolor{summit}{HTML}{6B7A8F}
\definecolor{gold}{HTML}{B8860B}
\definecolor{goldl}{HTML}{FBF3DC}
\definecolor{pine}{HTML}{2C7A55}
\definecolor{pinel}{HTML}{E2F1E9}
\definecolor{prune}{HTML}{5B3F8C}
\definecolor{prunel}{HTML}{EEE8F6}
\AtBeginDocument{\color{ink}}

% ---------------------------------------------------------------- Titres
% \section = chapitre du cours (« Chapitre 1 : ... »), \subsection = I, II, III..., \subsubsection = A, B, C...
\titleformat{\section}[display]
  {\sffamily\color{night}}
  {\large\bfseries\color{alert}\MakeUppercase{\sectionname}~\thesection}{2pt}{\LARGE\bfseries}
  [{\vspace{2pt}\color{night!35}\titlerule[1.3pt]}]
\titlespacing*{\section}{0pt}{6pt}{10pt}
\newcommand{\sectionname}{\nomchapitre}
\renewcommand{\thesubsection}{\arabic{section}.\arabic{subsection}}
\renewcommand{\thesubsubsection}{}
\titleformat{\subsection}
  {\large\bfseries\sffamily\color{sky}}{\thesubsection}{0.6em}{}
\titlespacing*{\subsection}{0pt}{12pt}{5pt}
\titleformat{\subsubsection}
  {\normalsize\bfseries\sffamily\color{alert}}{}{0em}{}
\titlespacing*{\subsubsection}{0pt}{9pt}{3pt}
\setcounter{secnumdepth}{3}
\setcounter{tocdepth}{2}

\usepackage{etoolbox}
\pretocmd{\section}{\clearpage}{}{}
\pretocmd{\subsection}{\needspace{8\baselineskip}}{}{}
\pretocmd{\subsubsection}{\needspace{5\baselineskip}}{}{}
\renewcommand{\bottomtitlespace}{0.15\textheight}
\widowpenalty=10000
\clubpenalty=10000
\raggedbottom

% ---------------------------------------------------------------- En-têtes
\pagestyle{fancy}
\fancyhf{}
\renewcommand{\headrulewidth}{0.4pt}
\renewcommand{\headrule}{\hbox to\headwidth{\color{night!40}\leaders\hrule height \headrulewidth\hfill}}
\fancyhead[L]{\sffamily\small\color{night}\textbf{COUPOLE · \titrecourt}}
\fancyhead[R]{\sffamily\small\color{summit}\partiecourante}
\fancyfoot[L]{\sffamily\scriptsize\color{summit}Coupole \versioncoupole{} · GPL-3.0 · ARP273-ROSE}
\fancyfoot[R]{\sffamily\small\color{summit}\thepage}
\newcommand{\partiecourante}{}
\newcommand{\partie}[1]{\clearpage\renewcommand{\partiecourante}{#1}}


% ---------------------------------------------------------------- Parties
\titleformat{\part}[display]{\centering\sffamily\color{night}}
  {\Large\bfseries\color{alert}PARTIE \thepart}{8pt}{\Huge\bfseries}
\pretocmd{\part}{\clearpage}{}{}
\setcounter{tocdepth}{2}

% ---------------------------------------------------------------- Encadrés du site (une icône par type de section)
\tcbset{maison/.style={enhanced,breakable,boxrule=0.9pt,arc=2.5pt,left=2.5mm,right=2.5mm,top=1.5mm,bottom=1.5mm,
  fonttitle=\bfseries\sffamily\small,coltitle=white,
  before upper={\widowpenalties 4 10000 10000 10000 0 \clubpenalties 4 10000 10000 10000 0}}}
% \typesite{env}{icône}{libellé}{fond}{cadre}
\newcommand{\typesite}[5]{\newtcolorbox{#1}[1][]{maison,colback=#4,colframe=#5,title={#2~~#3\ifx\relax##1\relax\else~·~##1\fi}}%
  \BeforeBeginEnvironment{#1}{\par\needspace{6\baselineskip}}}
\typesite{cdefinition}{\faBook}{\tr{Définition}{Definition}}{skyl}{sky}
\typesite{crappel}{\faUndo}{\tr{Rappel}{Reminder}}{goldl}{gold!85!black}
\typesite{cremarque}{\faInfoCircle}{\tr{Remarque}{Note}}{skyl}{sky!70!black}
\typesite{cattention}{\faExclamationTriangle}{\tr{Attention}{Warning}}{alertl}{alert}
\typesite{cexemple}{\faLightbulb}{\tr{Exemple}{Example}}{pinel}{pine}
\typesite{cexercice}{\faPencil*}{\tr{Exercice}{Exercise}}{white}{pine}
\typesite{cdemonstration}{\faCalculator}{\tr{Démonstration}{Proof}}{prunel}{prune}
\typesite{cconseil}{\faThumbsUp}{\tr{Conseil}{Tip}}{pinel}{pine!80!black}
\typesite{ccomplement}{\faPlusCircle}{\tr{Complément}{Supplement}}{prunel}{prune!80!black}
\typesite{censavoirplus}{\faSearchPlus}{\tr{En savoir plus}{Further reading}}{prunel}{prune!70!black}
\typesite{cactivite}{\faFlask}{\tr{Activité}{Activity}}{goldl}{gold}
\typesite{cqcm}{\faCheckSquare}{\tr{QCM}{Quiz}}{white}{sky}
\typesite{cobjectifs}{\faBullseye}{\tr{Objectifs}{Goals}}{skyl}{night}
\typesite{cprerequis}{\faListUl}{\tr{Prérequis}{Prerequisites}}{skyl}{night}
\typesite{ccadre}{\faStar}{\tr{À retenir}{Key point}}{goldl}{gold}
\newtcolorbox{ccode}{maison,colback=night!4,colframe=night!50,fontupper=\small\ttfamily,title={\faTerminal~~\tr{Commande}{Command}}}

\newcommand{\intertitre}[2][]{\par\needspace{4\baselineskip}\medskip\noindent{\sffamily\bfseries\color{night}#2}\par\smallskip}
\newcommand{\question}[1]{\par\smallskip\noindent{\sffamily\bfseries\color{pine}#1}~}
\newcommand{\auteurs}[1]{\par\noindent{\sffamily\small\color{summit}\faUserEdit~#1}\par\medskip}
\newcommand{\appli}{\textcolor{summit}{\faMousePointer~\textit{(application interactive en ligne)}}}
% \figsite{titre}{image}{légende}{crédit}
\newcommand{\figsite}[4]{\par\needspace{8\baselineskip}\begin{center}\begin{minipage}{\linewidth}\centering
  \if\relax\detokenize{#1}\relax\else{\sffamily\small\bfseries\color{sky}#1}\par\smallskip\fi
  \if\relax\detokenize{#2}\relax{\sffamily\small\color{summit}\faFilm~\textit{animation ou vidéo en ligne, non reproduite}}\par\else#2\par\fi
  \if\relax\detokenize{#3}\relax\else\smallskip{\small #3}\par\fi
  \if\relax\detokenize{#4}\relax\else{\scriptsize\color{summit}Crédit : #4}\par\fi
  \end{minipage}\end{center}}
\DeclareUnicodeCharacter{2009}{\,}
\DeclareUnicodeCharacter{2013}{--}
\DeclareUnicodeCharacter{2014}{---}
\DeclareUnicodeCharacter{2019}{'}
\DeclareUnicodeCharacter{2018}{`}
\DeclareUnicodeCharacter{201C}{``}
\DeclareUnicodeCharacter{201D}{''}
% \ressource{titre}{niveau · format · volume}{url}{description} — guide du portail
\newcommand{\ressource}[4]{\par\needspace{5\baselineskip}\noindent{\sffamily\bfseries\color{night}#1}\\
  {\sffamily\footnotesize\color{summit}#2 · \href{#3}{\faExternalLink*~lien}}\\
  {\small #4}\par\smallskip}

\renewcommand{\partiecourante}{Coupole \versioncoupole}
\newcommand{\capture}[2]{\par\needspace{10\baselineskip}\begin{center}\includegraphics[width=#2\linewidth]{#1}\end{center}}
\newcommand{\cmd}[1]{\texttt{#1}}

\begin{document}
% ---------------------------------------------------------------- page de garde
\begin{titlepage}
\centering
\vspace*{2.5cm}
\includegraphics[width=4.2cm]{../../logo/coupole_512.png}\par\vspace{1.2cm}
{\sffamily\bfseries\color{night}\fontsize{38}{44}\selectfont Coupole\par}
\vspace{0.4cm}
{\sffamily\Large\color{sky}Manuel de référence\par}
\vspace{0.3cm}
{\sffamily\large\color{summit}Version \versioncoupole{} --- \datecoupole\par}
\vspace{1.4cm}
\begin{minipage}{0.78\linewidth}\centering\sffamily\color{ink}
Boîte à outils libre pour les étudiants du DU « Explorer et Comprendre l'Univers » (Observatoire de Paris).\\[4pt]
Banque d'images de l'OHP · partages du DU · qualité des images · spectres et séries · sites et heures · cosmologie · animations · calculateurs de TD.
\end{minipage}
\vfill
{\sffamily\small\color{summit}Auteur : ARP273-ROSE · Licence GNU GPL version 3 ou ultérieure\\
Données : banque « OHP student observations », Observatoire de Paris / PADC (licence Etalab 2.0)\par}
\end{titlepage}

\tableofcontents

% =====================================================================
\section{Présentation}
Coupole est une application libre (GPL-3.0), écrite en Python 3.10+, qui fonctionne sous Windows, macOS et Linux,
avec une interface graphique (PyQt6) et une ligne de commande complète. Tout l'affichage existe en français et en
anglais ; la langue du système est détectée, et peut être imposée (menu \emph{Langue}, Préférences, ou
\cmd{--lang fr|en}).

\begin{center}\small
\begin{tabularx}{\linewidth}{@{}lX@{}}
\toprule
\textbf{Module} & \textbf{Fonction}\\
\midrule
Banque OHP & Inventaire, catalogue, téléchargement, contrôle astrométrique, en-têtes corrigés, conversion, lots empilables, anomalies, carte du ciel.\\
Partages du DU & Dossiers collaboratifs des défis et TP : structure, aperçus, nouveautés, tableaux de mesures contrôlés, outils d'analyse (Parallaxe Lune et ISS), dépôt de sa mesure sans rien écraser.\\
Archives des observatoires & Images publiques de Hubble, JWST, ESO, Spitzer, WISE, 2MASS, NOIRLab, Keck, SDSS, Voyager, Cassini, Juno : recherche, volume, téléchargement, extraction, alignement, couleur.\\
Qualité des images & Facultatif. FWHM, ellipticité, fond, bruit, gradient, saturation, traînées, échantillonnage.\\
Spectres et séries & Lecture et tracé de données 1D (spectres radio, courbes de lumière, tables).\\
Cosmologie & Du redshift aux distances, âges, volume, module de distance, échelle (Planck 2018 et autres).\\
Sites et heures & Base des sites, carte OpenStreetMap, heure UTC et heure locale.\\
Calculateurs de TD & La démarche des TD et TP pas à pas avec vos nombres, puis le calcul exact : Parallaxe (TP du cours).\\
Ma machine & Diagnostic matériel ; exemple minimal de module.\\
\bottomrule
\end{tabularx}
\end{center}

Principes : tout travail long s'exécute hors du fil graphique (l'interface reste réactive, avec progression et
annulation) ; un traitement interrompu reprend où il s'est arrêté ; rien n'est supprimé sans trace ; aucune
donnée n'est envoyée sans accord ; les composants facultatifs (ASTAP, SEP, CuPy) peuvent manquer.

% =====================================================================
\section{Installation}
\subsection{Paquets tout compris}
Les paquets embarquent Python et toutes les bibliothèques (PyQt6, numpy, scipy, astropy, zstandard, lz4) : rien
d'autre n'est à installer. Ils se mettent à jour d'eux-mêmes. ASTAP reste facultatif et n'est pas inclus.

\textbf{Bibliothèques facultatives embarquées (0.2.1).} Facultatives pour une installation par pip (extras
\texttt{qualite}, \texttt{alignement}), elles sont dans \emph{tous} les paquets (installeur Windows, .dmg, tar.gz,
.deb) : \texttt{reproject} (alignement exact et adaptatif du module Archives, avec \texttt{astropy-healpix},
\texttt{dask} et \texttt{zarr} qu'il demande ; jusqu'en 0.2.0 il manquait aux paquets, et l'alignement retombait sur
un rééchantillonnage bilinéaire de scipy), \texttt{sep} (module Qualité des images), \texttt{psutil} (mesures de la
machine) et \texttt{lxml} (validation XISF). Seul CuPy (carte NVIDIA, inutilisé pour l'instant) reste à part. Chaque
paquet construit les importe et les fait tourner avant publication ; poids ajouté : environ 20~Mo à l'archive Linux
(101 → 121~Mo), 12~Mo au .deb (109 → 121~Mo). Le module \emph{Ma machine} et \emph{Aide > À propos} donnent les
bibliothèques facultatives présentes et absentes, avec leur version.

\begin{center}\small
\begin{tabularx}{\linewidth}{@{}l>{\ttfamily\small}lX@{}}
\toprule
\textbf{Système} & \textbf{\normalfont Fichier} & \textbf{Installation}\\
\midrule
Windows 10/11 64 bits & Coupole-Setup-x.y.z.exe & Double-clic. Avertissement SmartScreen au premier lancement (installeur non signé) : « Informations complémentaires », « Exécuter quand même ». Installation dans le profil, sans droit administrateur.\\
macOS Apple Silicon & Coupole-x.y.z-macos-arm64.dmg & Glisser Coupole dans Applications. Premier lancement : clic droit, « Ouvrir ».\\
macOS Intel & Coupole-x.y.z-macos-x86\_64.dmg & Idem.\\
Ubuntu / Debian x86\_64, ARM64 & coupole\_x.y.z\_amd64.deb & \texttt{sudo apt install ./coupole\_x.y.z\_amd64.deb} (\texttt{\_arm64.deb} sur ARM) : apt installe les bibliothèques système de Qt ; commande \texttt{coupole}, entrée de menu. Ubuntu 22.04 ou plus récent, Debian 12 ou plus récent. Désinstaller : \texttt{sudo apt remove coupole}.\\
Autre Linux x86\_64 / ARM64 & Coupole-x.y.z-linux-*.tar.gz & Décompresser, lancer \texttt{installer.sh} (copie dans \texttt{\textasciitilde/.local/share/coupole}, entrée de menu, commande \texttt{coupole} par un lien \texttt{\textasciitilde/.local/bin/coupole} ; si ce dossier n'est pas dans le \texttt{PATH}, l'installeur le dit et donne la ligne à ajouter). Le lanceur \texttt{Coupole.sh} suit les liens symboliques (lien de lien compris).\\
\bottomrule
\end{tabularx}
\end{center}

\subsection{Avec le Python de l'ordinateur}
Les scripts du dossier des sources cherchent un Python 3.10 ou plus récent, créent un environnement virtuel,
installent les dépendances et créent les lanceurs. Une seule question : Entrée pour confirmer (\cmd{--oui} /
\cmd{-Oui} pour l'omettre).
\begin{ccode}
sh install.sh                \# Linux, macOS\\
install.bat                  \# Windows (double-clic ; lance install.ps1)
\end{ccode}
Sans Python, ou sans le module \texttt{venv} (Debian : paquet \texttt{python3-venv}), le script le dit et indique
comment l'obtenir. Emplacements : Linux \texttt{\textasciitilde/.local/share/coupole/venv} et lanceurs dans
\texttt{\textasciitilde/.local/bin} ; macOS \texttt{\textasciitilde/Library/Application Support/Coupole/venv} et
\texttt{\textasciitilde/Applications/Coupole.app} ; Windows \texttt{\%LOCALAPPDATA\%\textbackslash
Programs\textbackslash CoupolePython} et raccourcis Bureau et menu Démarrer.

\subsection{pip et pipx}
\begin{ccode}
pipx install git+https://github.com/ARP273-ROSE/coupole\\
pip install .                 \# depuis les sources ; pip install ".[qualite]" ajoute SEP\\
python -m coupole             \# interface ; coupole --help : ligne de commande
\end{ccode}
Extras : \texttt{qualite} (SEP), \texttt{validation} (lxml, xisf : validation XSD et second lecteur),
\texttt{gpu} (CuPy), \texttt{test} (pytest).

\subsection{Mise à jour, désinstallation, données}\label{sec:donnees}
Les paquets vérifient au démarrage (désactivable) la dernière version publiée sur GitHub et proposent de
l'installer : seule l'archive du code (\texttt{coupole-app-x.y.z.zip}, environ 1 Mo) est téléchargée, extraite par
l'application, l'ancienne version étant restaurée en cas d'échec. Une installation par pip/pipx affiche la commande
à lancer. Le paquet \texttt{.deb} ne se met pas à jour par archive : ses fichiers (\texttt{/opt/coupole}) appartiennent à
dpkg ; Coupole signale la nouvelle version, donne l'adresse du nouveau \texttt{.deb} et la commande \texttt{apt} à
lancer (bouton \emph{Télécharger le paquet}). Les réglages, le cache et les rapports vivent hors du dossier
d'installation et survivent aux mises à jour comme à la désinstallation :

\begin{center}\small
\begin{tabular}{@{}ll@{}}
\toprule
Windows & \texttt{\%LOCALAPPDATA\%\textbackslash Coupole}\\
macOS & \texttt{\textasciitilde/Library/Application Support/Coupole} (cache : \texttt{\textasciitilde/Library/Caches/Coupole})\\
Linux & \texttt{\textasciitilde/.config/coupole} (cache : \texttt{\textasciitilde/.cache/coupole})\\
Tous & variable \texttt{COUPOLE\_HOME} : remplace les deux (installations portables, tests)\\
\bottomrule
\end{tabular}
\end{center}

% =====================================================================
\section{Prise en main}
La fenêtre comporte une barre de modules à gauche (Ctrl+1 à Ctrl+9) et le module courant à droite. Chaque élément
interactif porte une info-bulle. Menu \emph{Aide} : aide de l'écran affiché (F1), ce manuel (Maj+F1), raccourcis
clavier, signaler un problème, rapports d'incident, recherche de mise à jour, À propos (version, licence,
configuration détectée).

\begin{center}\small
\begin{tabular}{@{}ll@{}}
\toprule
F1 & aide de l'écran\\ Maj+F1 & manuel PDF\\ Ctrl+1 … Ctrl+9 & module 1 à 9\\ Ctrl+, & Préférences\\
Ctrl+Maj+A & assistant ASTAP\\ Ctrl+Maj+D & thème sombre / clair\\ Ctrl+R & Banque OHP : rafraîchir l'inventaire\\ Ctrl+Q & quitter\\
\bottomrule
\end{tabular}
\end{center}

\capture{img/reglages_fr.png}{0.62}
Préférences (Ctrl+,) : langue de l'interface ; langue des noms de dossiers, d'objets et des commentaires FITS ;
apparence (section~\ref{sec:apparence}) ; boîtes de dialogue de fichiers (section~\ref{sec:dialogues}) ; PixInsight déjà ouvert (section~\ref{sec:ouvrir}) ; dossier de sortie ; format ; parallélisme ; débit de téléchargement (case \emph{Limiter le débit}, décochée par défaut : aucun plafond ; cochée : plafond de 0,1 à 1~000~Mo/s, pour la Banque OHP comme pour les Archives) ; mode économe ; recherche de mise à jour ; vérification des nouveautés de la banque au démarrage et sa fréquence maximale (section~\ref{sec:nouveautes}). L'onglet
\emph{Sources} donne les adresses des services (chapitre~\ref{chap:evolutif}). Un fichier de réglages illisible est
mis de côté (suffixe \texttt{.corrompu-…}) et les valeurs par défaut reprennent : l'application démarre toujours.

\subsection{Choisir un fichier ou un dossier, dossier sur un NAS}\label{sec:dialogues}
Les boutons \emph{Parcourir} ouvrent l'explorateur de fichiers du système, avec ses emplacements, favoris et partages
réseau : Explorateur sous Windows, Finder (NSOpenPanel) sous macOS, et sous Linux celui du bureau par le
\textbf{portail XDG} (\texttt{xdg-desktop-portal}) : Dolphin sous KDE Plasma, Fichiers (Nautilus) sous GNOME. Le paquet
autonome embarque sa propre bibliothèque Qt, qui ne peut pas charger le greffon du bureau installé sur le système :
jusqu'à la 0.1.7, il retombait sur le dialogue générique de Qt, en anglais et sans les emplacements du système. Au
démarrage, si l'utilisateur n'a pas déjà posé \texttt{QT\_QPA\_PLATFORMTHEME}, Coupole regarde (sans attendre plus
d'une seconde) si le portail est installé ou en service avec un moteur qui sait choisir des fichiers, et le prend ;
sinon, sous un bureau GTK (GNOME, Cinnamon, Xfce, MATE…), le dialogue GTK ; sinon, le dialogue de Qt, \textbf{en
français} et complété dans sa barre latérale par les dossiers personnels, \texttt{/media}, \texttt{/mnt},
\texttt{/run/media/\$USER}, les partages ouverts dans le gestionnaire de fichiers (\texttt{/run/user/<uid>/gvfs/…}) et
les montages cifs, nfs et sshfs trouvés dans \texttt{/proc/mounts}.

Préférences > \emph{Boîtes de dialogue de fichiers} : \emph{Système} (défaut) ou \emph{Qt}, le second seulement si le
dialogue du système ne s'ouvre pas ou fonctionne mal (portail défaillant). Sous Linux, revenir à \emph{Système} prend
effet au prochain lancement. Tous les textes fournis par Qt (boutons standard, menus contextuels des champs, dialogue
de repli) suivent la langue de l'interface et changent avec elle.

\paragraph{Dossier de sortie sur un NAS.}
\begin{itemize}[nosep]
\item \emph{Linux} : ouvrir d'abord le partage dans Dolphin ou Fichiers (\texttt{smb://serveur/partage}) ; il apparaît
  sous \texttt{/run/user/<uid>/gvfs/smb-share:server=serveur,share=partage} (GNOME ; KDE avec le paquet
  \texttt{kio-fuse}, sans lequel un partage \texttt{smb://} non monté ne peut pas être choisi par une application),
  ou monter le partage en cifs ou nfs (\texttt{/mnt/partage}, conseillé pour de gros volumes ; options par défaut, \texttt{nobrl} inutile : section~\ref{sec:partage}) ;
\item \emph{macOS} : Finder > Aller > Se connecter au serveur (Cmd+K), \texttt{smb://serveur/partage}, puis choisir le
  volume sous \texttt{/Volumes} ;
\item \emph{Windows} : lecteur réseau (\texttt{Z:}) ou chemin UNC \texttt{\textbackslash\textbackslash
  serveur\textbackslash partage}, qu'on peut aussi taper dans le champ du dossier.
\end{itemize}
Ces chemins sont acceptés partout : base d'état lue en lecture seule (URI SQLite d'un chemin UNC écrite avec une
autorité vide), place libre, test d'écriture, chemins longs Windows (\texttt{\textbackslash\textbackslash?\textbackslash
UNC\textbackslash…}), affichage coupé à la ligne ; un partage est reconnu comme tel (UNC, lecteur réseau, cifs, nfs,
sshfs, gvfs, kio-fuse) et le module Qualité y limite ses lectures parallèles.

\subsection{Apparence}\label{sec:apparence}
Coupole a son propre thème, indépendant de celui de l'ordinateur (un système en mode sombre ne mélange plus ses
couleurs à celles de l'application : style Fusion, palette complète, taille de police fixe, contraste WCAG $\geq$ 4,5
pour le texte, les champs, les boutons, les info-bulles et la sélection). Le thème \textbf{sombre} est celui par défaut ;
le thème clair reste au choix. Il se change par le menu \emph{Affichage} > \emph{Apparence} (entrées cochables
\emph{Clair} / \emph{Sombre}, Ctrl+Maj+D pour basculer), ou dans les Préférences : c'est le même réglage, appliqué
immédiatement et enregistré ; les deux restent synchronisés. Un utilisateur qui avait choisi le thème clair le garde :
le réglage enregistré prime sur le défaut. Les captures de ce manuel sont prises dans le thème par défaut.

\subsection{Réglages conservés}\label{sec:conserves}
Coupole retrouve à chaque lancement l'écran tel qu'on l'a laissé. Sont conservés :
\begin{itemize}
\item \textbf{fenêtre} : taille, position, écran, état maximisé ; taille de chaque dialogue (Préférences et son
  onglet, assistant ASTAP, À propos, signalement, aide, correction d'un objet, ajout d'un site) ;
\item \textbf{disposition} : module affiché, onglet de la Banque OHP, position des séparateurs, largeur et ordre des
  colonnes, colonne et sens du tri de chaque tableau ;
\item \textbf{Banque OHP} : recherche, type, télescope, cases \emph{Nouveaux}, \emph{À vérifier}, \emph{À télécharger
  seulement}, \emph{Sans dates douteuses}, nuit et filtre, objets choisis (rétablis dès que l'inventaire est chargé),
  filtres des anomalies et de la carte du ciel ; onglet Téléchargement : dossier de sortie \emph{dès qu'il est modifié},
  format, langue des noms, \emph{Garder les doublons}, \emph{Garder les FITS}, mode ASTAP, \emph{Vérifier la qualité} ;
\item \textbf{Qualité des images} : dossier analysé, mode échantillon et nombre d'images par lot ;
\item \textbf{Spectres et séries} : dossier du dialogue d'ouverture et menu \emph{Récents} (10 derniers fichiers ;
  un fichier disparu y est grisé) ; axe choisi ;
\item \textbf{Cosmologie} : jeu de paramètres, paramètres personnalisés ($H_0$, $\Omega_m$, gardés même en repassant
  par Planck), $\Omega_k$, $z$ courant, option SH0ES, échelle des courbes (logarithmique ou linéaire) ;
\item \textbf{Sites et heures} : site choisi, zoom et centre de la carte du monde, \emph{Carte en ligne} ;
\item \textbf{dialogues de fichiers} : chacun rouvre dans le dernier dossier utilisé (export CSV, réorganisation,
  ASTAP, signalement…).
\end{itemize}
La disposition vit dans \texttt{interface.json}, à côté de \texttt{reglages.json} (section~\ref{sec:donnees}) ; les
choix de traitement restent dans \texttt{reglages.json}. L'écriture est \emph{différée et groupée} : au plus une toutes
les 2~s et une à la fermeture, atomique (fichier temporaire puis remplacement) ; une frappe dans un champ n'écrit
jamais le disque. Lecture tolérante : une valeur absente, d'un mauvais type ou hors bornes laisse le défaut ; un
fichier illisible est mis de côté (\texttt{.corrompu-…}). Garde-fous : un écran débranché ou une position hors des
écrans recentre la fenêtre, une taille plus grande que l'écran est ramenée à l'écran (comme à 150 ou 200~\%). Un
dossier gardé sur un partage réseau absent ne bloque jamais le démarrage : son existence est vérifiée en fond, et le
dialogue s'ouvre au dossier par défaut tant qu'elle n'est pas confirmée. Rien de ce qui serait dangereux à rejouer
n'est gardé : un traitement interrompu n'est jamais relancé seul (il reprend quand on le relance).

Pour revenir à l'origine : Préférences > \emph{Réinitialiser la disposition} (immédiat ; la langue, le thème, le
dossier de sortie et les autres réglages ne changent pas), ou \cmd{coupole --reinitialiser-interface}
(\cmd{--reset-interface}).

% =====================================================================
\section{Module Banque OHP}
\subsection{La banque}
« OHP student observations » réunit les images prises par les étudiants du DU à l'Observatoire de Haute-Provence :
7~989 images FITS au 8 octobre 2026, 166 objets sous 194 noms, 60 nuits de 2015 à 2026, 78,0 Go ; télescope de
1,20 m (T120, 7~419 images, 1024 pixels de 0,77″, champ de 13,1′) et télescope robotique IRIS de 50 cm (570 images,
2048 pixels de 0,68″ ou 4096 pixels de 0,46″, champs de 23′ et 31′). Les images sont déjà calibrées (dark, flat) et
résolues astrométriquement : \textbf{elles ne doivent pas être recalibrées}.

\subsection{Télécharger en trois gestes (0.7.2)}\label{sec:trois-gestes}
\capture{img/ohp_guide_fr.png}{1.0}
Au premier lancement, un bandeau en tête du catalogue guide pas à pas ; chaque étape faite est cochée.
\begin{enumerate}[leftmargin=*,itemsep=1pt]
\item \textbf{Choisir une cible} : un clic sur une ligne de la liste de gauche (Ctrl+clic ou Maj+clic pour
  plusieurs ; la recherche et la liste des types aident à la trouver). Tant que rien n'est choisi, la liste des
  images le dit : « Aucune cible choisie : cliquez sur une cible dans la liste de gauche ».
\item \textbf{Voir ses images} : elles s'affichent à droite, filtrables par nuit, filtre et dates douteuses.
\item \textbf{Télécharger} : le bouton bleu sous la liste des images annonce ce qu'il va faire,
  « Télécharger (840 images · 7,15 Go → ≈ 2,64 Go en XISF) ». La première fois, Coupole demande le dossier où
  ranger les images (dialogue de fichiers complet du système ; il y crée \texttt{OHP\_DU\_ECU}) ; ensuite il
  reprend le même dossier et le même format. La progression s'affiche sous le bouton, avec \emph{Arrêter} ; à la
  fin, la pastille de chaque image devient une coche verte et la barre d'état le dit (« Téléchargement terminé :
  … Double-clic sur une image pour l'ouvrir »).
\end{enumerate}
Le bouton est grisé, avec la raison au survol, tant qu'aucune cible n'est choisie ; il affiche « Déjà téléchargé ✓ »
quand tout est là, « Téléchargement en cours… » pendant un téléchargement. \emph{Compris, masquer} ferme le
bandeau ; \emph{Aide > Guide de démarrage (Banque OHP)} le réaffiche. Les réglages (dossier, format XISF ou FITS,
doublons, FITS d'origine, astrométrie, parallélisme) sont dans l'onglet \emph{Téléchargement}, sous
\emph{Options} (bouton \emph{Options du téléchargement →} du catalogue).

\capture{img/ohp_menu_fr.png}{1.0}
\textbf{Clic droit et double-clic.} Le menu d'une image commence par \emph{Télécharger cette image} (ou
\emph{Télécharger ces N images} si plusieurs lignes sont choisies), puis \emph{Ouvrir}, \emph{Ouvrir avec},
\emph{Ouvrir l'emplacement du fichier}, grisés avec leur raison (« pas encore téléchargée ») tant que le fichier
n'existe pas ; celui d'une cible commence par \emph{Télécharger toutes les images de NGC 6888}. Un double-clic sur
une image (ou une cible) pas encore téléchargée propose de la télécharger ; sur une image téléchargée, il l'ouvre.

\paragraph{À donner aux étudiants.} « Dans Coupole, module \emph{Banque OHP} : 1) cliquez sur une cible dans la
liste de gauche ; 2) ses images s'affichent à droite ; 3) cliquez sur le bouton bleu \emph{Télécharger} sous la
liste des images et choisissez votre dossier la première fois. Ensuite, double-clic sur une image pour l'ouvrir. »

\subsection{Catalogue}
\capture{img/catalogue_fr.png}{1.0}
L'inventaire vient du service TAP public de l'Observatoire (\texttt{http://tap-ufe.obspm.fr/tap}, table
\texttt{ivoa.obscore}, une requête ADQL synchrone), gardé en cache ; un instantané est livré pour travailler hors
ligne. \emph{Rafraîchir l'inventaire} (Ctrl+R) réinterroge le service et signale les images et les noms nouveaux
(« nouveau depuis le … »). Les noms de la base sont regroupés en objets canoniques ; la sélection d'un ou plusieurs
objets affiche leurs images à droite, filtrables par nuit, filtre et dates douteuses ; cette liste est la sélection
à traiter. L'estimation donne le volume à télécharger, le volume attendu en sortie et la place libre. Les heures
sont affichées en UTC et à l'heure du site (survol : également l'heure de l'utilisateur). Chaque en-tête de colonne
s'explique au survol ; les colonnes clés (état, date, télescope, filtre, pose ; objet, possédé, volume) viennent en
premier et restent visibles sans défiler en 1366~×~768.

\paragraph{Ce qu'on possède déjà.} D'après le dossier de sortie (\texttt{\_traitement/etat.sqlite}, relu en fond à
l'ouverture, après chaque traitement et à chaque changement de dossier, jamais bloquant), la liste des images
distingue par une pastille et une couleur de texte douce (lisible dans les deux thèmes, contraste $\geq 4{,}5$) :
\textbf{possédée} (convertie et rangée, coche verte), \textbf{doublon écarté} (mêmes pixels qu'une image déjà
convertie, rond gris), \textbf{échec} (à retenter, triangle orange) et \textbf{à télécharger} (texte normal,
flèche) ; l'info-bulle de la pastille donne le statut, la date et le chemin local du fichier. La \textbf{liste des
objets} montre le même état, aux mêmes couleurs : chaque ligne commence par une pastille (coche verte : tout possédé,
doublons écartés compris ; demi-disque vert : en partie ; triangle orange : au moins un échec ; flèche : rien de
téléchargé), le nom de l'objet prend la même couleur, et l'info-bulle détaille (« 12 possédées / 12 · 0 à télécharger
· 0 échec(s) · 3 doublon(s) écarté(s) »). La colonne \emph{possédé} (« 120 / 300 », mini-barre) suit directement le
nom (un ordre de colonnes choisi à la main est respecté) ; la trier range les objets par état (échec, rien, en
partie, complet). La légende, sous les deux listes, vaut pour les deux. La ligne de résumé de l'inventaire ajoute la
possession globale (« 7 989 images (144 doublons), 166 objets, 60 nuits, 78,03 Go — possédées : 7 625 (+ 220
doublons écartés), échecs : 0, à télécharger : 0 »). La case \emph{À télécharger seulement} cache les objets
entièrement possédés et, dans la liste, les images déjà traitées ; une liste vidée l'explique au centre (« Tout est
déjà téléchargé dans … »), comme la liste des images tant qu'aucun objet n'est choisi (« Aucune cible choisie : cliquez sur une cible dans
la liste de gauche » ; au premier lancement, aucune cible n'est choisie : le guide et cette liste disent quoi faire). L'estimation ne compte que ce qui manque
(« 180 images manquantes sur 300 : 1,5 Go »). Changer de dossier de sortie (Préférences, champ ou \emph{Parcourir}
de l'onglet Téléchargement) relit aussitôt sa possession en fond : pastilles, colonne et résumé suivent, et la barre
d'état l'annonce (« Possession recalculée : 7 625 images trouvées dans … »).
Ligne de commande : \cmd{coupole ohp inventaire --manquantes [--dest DOSSIER] [--json]}, et
colonnes \texttt{possedee}, \texttt{statut\_local}, \texttt{fichier\_local} dans le CSV de \cmd{coupole ohp images}.

\subsection{Téléchargement (onglet et options)}
\capture{img/traitement_fr.png}{1.0}
L'onglet \emph{Téléchargement} (anciennement \emph{Traitement}) montre la sélection, le même bouton
\emph{Télécharger}, \emph{Tout télécharger…}, \emph{Pause}, \emph{Arrêter}, la progression et le journal ; il réunit
sous \emph{Options} le dossier et le format, la langue des noms, \emph{Garder les doublons}, \emph{Garder les FITS
d'origine}, \emph{Vérifier la qualité des images}, la vérification astrométrique (ASTAP) et le parallélisme. Les noms
de fichiers et de dossiers internes ne changent pas (\texttt{\_traitement/etat.sqlite}…), ni la ligne de commande
(\cmd{coupole ohp traiter}) : les dossiers déjà remplis restent compatibles.
Pour chaque image, hors doublons :
\begin{enumerate}[leftmargin=*,itemsep=1pt]
\item téléchargement : au plus 4 simultanés (3 par défaut), débit libre par défaut (plafond total réglable dans les Préférences), User-Agent
  identifiant Coupole, reprise par requêtes HTTP Range, contrôle de la taille (serveur et inventaire), du multiple
  de 2~880 octets et de la carte \texttt{SIMPLE} ;
\item lecture des cartes FITS telles qu'écrites et des pixels ; empreinte SHA-1 des pixels (doublons exacts) ;
\item contrôle de la solution astrométrique (section~\ref{sec:astrometrie}) ;
\item en-tête corrigé (section~\ref{sec:entetes}) ;
\item écriture au format choisi, relecture et comparaison pixel à pixel ;
\item à la fin : tri en lots, \texttt{LOT.txt}, \texttt{INDEX\_LOTS.csv}, \texttt{\_traitement/journal.csv}.
\end{enumerate}
L'état est gardé dans \texttt{\_traitement/etat.sqlite} : \emph{Arrêter} puis \emph{Télécharger} (ou relancer la même
commande) reprend sans rien refaire. Un dossier garde la langue de noms et le format de son premier traitement.
\emph{Garder les doublons} traite aussi ce qui serait écarté ; \emph{Garder les FITS d'origine} les conserve dans
\texttt{\_traitement/fits\_origine}.

La chaîne est un \emph{pipeline} : les téléchargements (fils) et les conversions (processus séparés, un par cœur
disponible) se recouvrent, la file des FITS en attente est bornée ($2\times$ conversions + téléchargements ; si la
place libre à destination est juste, elle est réduite automatiquement — jamais sous conversions + 1 — et le journal le
dit, plutôt que de refuser le traitement), et le
débit global est celui du réseau ou du plafond choisi (mesuré : 24 images de 8 Mo à 8 Mo/s en 24,5 s, soit la durée
du seul téléchargement ; sans plafond, 86 Mo/s de FITS convertis sur 6 cœurs). \emph{Pause} suspend
l'alimentation du pipeline (les images en cours se terminent) et \emph{Reprendre} la rétablit ; \emph{Arrêter}
interrompt téléchargements et conversions sans attendre (les processus sont terminés, ASTAP compris : par signal
sous Linux et macOS, par un \emph{job object} sous Windows). À la fin, un
rapport résume images, durée, converties, doublons, échecs (avec leur raison), volume écrit et lots.

Sur un partage réseau (NAS par SMB), où chaque accès est un aller-retour, l'état est validé au plus une fois par
seconde et \texttt{JOURNAL.txt} écrit par paquets (une coupure ne fait refaire que la dernière seconde) ; le rangement
final ne réécrit que les \texttt{LOT.txt} qui ont changé et ne parcourt plus toute l'arborescence : relancer un
traitement pour une image ne coûte plus une minute de rangement.
La base d'état, elle, n'est jamais écrite directement sur un partage : section~\ref{sec:partage}.

\subsection{Tout télécharger, journal, réorganiser}\label{sec:tout}
\emph{Tout télécharger…} prend la banque entière (environ 8~000 images, 78 Go de FITS, 30 Go en XISF). Au premier
usage (ou si aucun dossier de sortie n'est défini), Coupole demande le dossier de destination par un dialogue natif et
propose \texttt{Documents/Coupole/OHP\_DU\_ECU} ; puis il affiche, \textbf{avant toute confirmation}, le volume à
télécharger, le volume attendu en sortie, la durée estimée --- au plafond de débit s'il est posé (78 Go à 8 Mo/s : 2~h~42), sinon au débit mesuré
lors des téléchargements précédents, sinon « selon la connexion » --- et la place
libre ; il refuse clairement s'il en manque. Le traitement est celui de toute sélection : reprise après coupure ou
fermeture (état par image dans \texttt{\_traitement/etat.sqlite}, FITS déjà téléchargés retrouvés), pause, rapport
de fin. Ligne de commande : \cmd{coupole ohp tout --oui}.

\emph{Ouvrir le journal} affiche \texttt{\_traitement/JOURNAL.txt} : une ligne horodatée (UTC) et bilingue par
événement --- début et fin de session (reprise, pause), chaque image téléchargée, convertie ou écartée avec sa raison,
chaque erreur, chaque conflit de nom ; rotation à 5 Mo (\texttt{JOURNAL-1.txt}, \texttt{JOURNAL-2.txt}). Le détail
tabulaire reste dans \texttt{journal.csv}.

\emph{Réorganiser des fichiers…} range dans l'arborescence des lots des images déjà converties par Coupole qui se
trouvent ailleurs ou selon un ancien rangement : chaque fichier est reconnu par son en-tête (adresse d'origine
\texttt{OHP:Source:URL}, ou carte \texttt{HISTORY} « converti de … »), rattaché à l'inventaire et \textbf{déplacé}
(jamais copié) ; un nom déjà pris par un fichier étranger reçoit un suffixe et une ligne de journal ; rien n'est
jamais écrasé. Les en-têtes sont lus par plusieurs processus, la progression s'affiche (« 1 200 / 7 625 en-têtes
lus ») et \emph{Arrêter} interrompt la lecture (ce qui est déjà reconnu est rangé). Ligne de commande :
\cmd{coupole ohp reorganiser DOSSIER --dest DEST}.

\subsection{Dossier de sortie sur un partage réseau}\label{sec:partage}
\textbf{Pourquoi.} La base d'état (\texttt{\_traitement/etat.sqlite}) est une base SQLite, qui pose des verrous sur
son fichier à chaque écriture. Sous Linux, un partage SMB monté par \texttt{cifs} avec les options par défaut
transmet ces verrous au serveur comme verrous de plage SMB, et l'écriture échoue : « database is locked » au bout
d'une minute, base de 0 octet (vérifié sur un vrai serveur Samba ; c'était le cas de toute version jusqu'à 0.1.9 :
télécharger les nouveautés, \emph{Tout télécharger} ou réorganiser vers un dossier du NAS échouait). La lecture
seule et la copie de fichiers, elles, fonctionnent.

\textbf{Ce que fait Coupole (0.1.10).} Quand le dossier de sortie est reconnu comme un partage (cifs, smb3, nfs,
sshfs, gvfs, kio-fuse sous Linux ; smbfs, afpfs, nfs sous macOS ; chemin UNC ou lecteur réseau sous Windows) ou
quand SQLite n'arrive pas à y écrire en 3~s (partage non reconnu, autre programme), il travaille sur une
\textbf{base de travail locale}, dans le dossier de cache de l'utilisateur (\texttt{bases/<clé>/etat.sqlite},
une par dossier de sortie, la clé étant tirée du chemin), et le dit dans le journal et la barre d'état :
« Dossier sur un partage réseau : base de travail locale, recopiée sur le partage toutes les 30~s ».
\begin{itemize}[nosep]
\item \emph{Au début} : la base du partage est copiée (simple copie de fichier) dans la base de travail.
\item \emph{Pendant le traitement}, toutes les 30~s, \emph{à la fin}, à l'\emph{arrêt} et à l'\emph{annulation} :
  la base de travail est recopiée sur le partage --- instantané cohérent, \texttt{PRAGMA integrity\_check},
  copie dans un fichier temporaire du même dossier, relecture et comparaison de l'empreinte SHA-256, puis
  remplacement d'un coup (\texttt{os.replace}). Aucun verrou n'est jamais posé sur le partage : pas de verrou
  orphelin possible après un plantage, et un lecteur (autre ordinateur, NAS) voit toujours une base complète.
\item \emph{Au démarrage suivant} : si la base du partage est plus récente (traitée ailleurs), elle est recopiée ; si
  la base de travail a des écritures qui n'ont pas été recopiées (plantage, coupure), elles le sont aussitôt.
\item \emph{Deux écrivains} (deux ordinateurs, ou le NAS, ont traité vers le même dossier pendant ce temps) : rien
  n'est écrasé. Coupole le dit clairement et propose une \textbf{fusion} : pour chaque image, le statut le plus
  avancé est gardé (convertie > doublon écarté > échec > en cours), à égalité le plus récent ; empreintes de
  pixels réunies. Les deux bases d'origine sont gardées à côté de la base de travail
  (\texttt{etat.sqlite.local-<date>}, \texttt{etat.sqlite.partage-<date>}). Ligne de commande :
  \cmd{coupole ohp fusionner --dest DOSSIER}.
\end{itemize}
La possession, les nouveautés et les anomalies lisent la base de travail quand elle est à jour (plus rapide qu'une
lecture page par page sur le réseau), sinon la base du partage, en lecture seule. \texttt{JOURNAL.txt}, les CSV et
les \texttt{LOT.txt} restent écrits directement sur le partage (fichiers texte, remplacés d'un coup). Le cache du
module Qualité n'est qu'un cache : il reste dans le dossier de cache de l'utilisateur, sans recopie
(section~\ref{sec:qualite}).

\textbf{Selon le système.} \emph{Linux} : cifs par défaut échoue (cause ci-dessus) ; monter avec \texttt{nobrl}
n'est plus nécessaire. \emph{macOS} : le client smbfs gère ses verrous autrement (selon la façon dont SQLite a été
compilé) ; ce n'a pas pu être vérifié ici, Coupole n'en dépend donc pas et utilise la même base de travail.
\emph{Windows} : SQLite sait en principe écrire sur un partage (verrous SMB natifs), mais plusieurs écrivains y
risquent d'abîmer la base ; même base de travail, vérifiée en intégration continue sur
\texttt{\textbackslash\textbackslash localhost\textbackslash C\$}. Si la base du partage est ouverte par un autre
programme au moment du remplacement, Coupole réessaie puis attend le passage suivant.

\textbf{Windows : lecteur réseau ou chemin UNC (0.2.1).} Un dossier de sortie sur une lettre de lecteur réseau
(\texttt{O:\textbackslash OHP\_DU\_ECU}, partage mappé) ou sur un chemin UNC
(\texttt{\textbackslash\textbackslash serveur\textbackslash partage\textbackslash OHP\_DU\_ECU}) est reconnu comme
partage (\texttt{GetDriveTypeW} = lecteur distant). Jusqu'en 0.2.0, la base d'état d'un dossier sur une \emph{lettre}
de lecteur réseau était illisible : le chemin était « résolu » en chemin UNC, et l'adresse SQLite qui en résultait
(\texttt{file://serveur/…}) était refusée ; l'erreur passait inaperçue et tout paraissait à télécharger, alors que le
même dossier vu de Linux était entièrement possédé. Depuis 0.2.1, l'adresse est construite sans résolution
(\texttt{file:///O:/…}, \texttt{file:////serveur/partage/…}) et la base d'un partage est lue par une \textbf{copie
locale} (cache de l'utilisateur, refaite seulement quand la base change) ; une base en journal WAL (ancien script)
est recopiée avec son \texttt{-wal}, que Windows ne sait pas lire sur un partage. Lettre et chemin UNC d'un même
partage ont chacun leur base de travail (la clé vient du chemin tapé) : sans conséquence, chacune suit le partage.
Vérifié en intégration continue sur le partage administratif \texttt{C\$} de la machine, sous son nom (pas
« localhost », que SQLite accepte dans une adresse), en UNC et par une lettre posée par \texttt{net use}.

\textbf{Dossier de suivi.} La ligne de résumé de la Banque OHP et la barre d'état disent le dossier réellement lu et
ce qu'on y a trouvé : « Dossier de suivi : <chemin> --- N images possédées lues ». Si la base de suivi existe mais
ne peut pas être lue, un \textbf{bandeau d'erreur} le dit avec la cause (« La base de suivi de ce dossier n'a pas pu
être lue : … ») au lieu d'afficher « rien de possédé » ; \texttt{coupole.log} note chaque lecture (base de travail,
copie ou lecture directe, nombre d'images lues) et chaque synchronisation avec le partage. Si \texttt{\_traitement/}
manque alors que des images converties sont rangées (copie partielle, sauvegarde restaurée), le bandeau propose
\textbf{Reconnaître les fichiers existants} : l'en-tête de chaque fichier (\texttt{OHP:Source:URL}, sinon la carte
\texttt{HISTORY} « converti de … ») donne l'image d'origine, inscrite comme possédée ; en fond, avec progression ;
rien n'est téléchargé, déplacé ni réécrit (\cmd{coupole ohp reconnaitre --dest DOSSIER}). \emph{Rafraîchir
l'inventaire} relit aussi la possession.

\subsection{Nouveautés de la banque}\label{sec:nouveautes}
À l'ouverture (au plus une fois par jour, réglable, désactivable dans les Préférences), si Internet est disponible et
qu'une copie locale existe dans le dossier de sortie, Coupole compare l'inventaire TAP frais à la copie. S'il y a des
images nouvelles, un bandeau apparaît dans la Banque OHP (et un message dans la barre d'état) : « N nouvelles images
(M objets, X Go) depuis le JJ/MM --- Télécharger maintenant ? / Plus tard / Voir ». \emph{Télécharger} traite les
seules nouveautés dans la même arborescence ; \emph{Voir} filtre le catalogue ; rien n'est jamais téléchargé sans
accord explicite. Le bouton \emph{Vérifier les nouveautés} lance la comparaison à la demande. Ligne de commande :
\cmd{coupole ohp nouveautes} (liste) et \cmd{coupole ohp nouveautes --telecharger}.

\subsection{Solution astrométrique}\label{sec:astrometrie}
Contrôles de cohérence, toujours faits : solution présente et lisible ; échelle à moins de 2~\% de la nominale
(\texttt{XPIXSZ}/\texttt{FOCALLEN}) ou de la médiane de la nuit ; angle à moins de 5° de la médiane de la nuit
(modulo 180°) ; centre à moins de 1,5° du médoïde de l'objet (fixe) ou 2° du médoïde objet et nuit (mobile).
Médianes et médoïdes viennent de l'inventaire entier.

Avec ASTAP (facultatif), chaque image est en outre résolue indépendamment :
\begin{center}\small
\begin{tabularx}{\linewidth}{@{}lX@{}}
\toprule
\textbf{Statut} & \textbf{Signification}\\
\midrule
confirmée & ASTAP retrouve le même centre à moins de 10″ et la même échelle à moins de 1~\%.\\
validée & solution gardée : contrôles bons, ASTAP muet (champ pauvre, trop dense, pose courte) ou absent.\\
refaite & solution absente ou incohérente, remplacée par celle d'ASTAP (avec SIP) ; l'ancienne passe en HISTORY.\\
douteuse & contrôles en échec et pas d'ASTAP pour trancher : image rangée à part.\\
échec & rien d'utilisable : image rangée dans \texttt{\_sans\_solution\_astrometrique}.\\
\bottomrule
\end{tabularx}
\end{center}

\subsection{En-têtes corrigés}\label{sec:entetes}
Toutes les cartes d'origine sont gardées dans l'ordre ; chaque valeur remplacée est recopiée dans une carte
\texttt{HISTORY valeur d'origine : …}.
\begin{center}\small
\begin{tabularx}{\linewidth}{@{}>{\ttfamily}lX@{}}
\toprule
\normalfont\textbf{Mot-clé} & \textbf{Correction}\\
\midrule
OBJECT & nom canonique (alias de la base fusionnés).\\
FILTER, FILTSYS, FILTORIG & filtre normalisé (B, V, R, Ha, OIII, g, r, i, z), système (Johnson, Cousins, SDSS, Gunn, Narrowband), valeur d'origine.\\
LATITUDE, LONGITUD & corrigées \textbf{seulement si l'inversion est démontrée} (valeurs échangées retombant sur le site) ; justes : intactes ; ambiguës : signalées, intactes. Opération idempotente.\\
SITELAT, SITELONG, SITEELEV & ajoutés s'ils manquent (site MPC 511 converti en WGS 84) ; présents et cohérents : gardés ; incohérents : signalés.\\
RA, DEC, OBJCTRA, OBJCTDEC & centre réel de la solution.\\
PIXSCALE, ORIENTAT & échelle mesurée, angle de position de l'axe Y.\\
FOCALLEN, XPIXSZ & recalculés seulement s'ils manquent ou sont incohérents de plus de 2~\%.\\
CTYPE1/2 & « -SIP » sans coefficients (IRIS, PinPoint) ramené à TAN.\\
DATE-OBS, DATEDOUT & ajoutée si absente ; \texttt{DATEDOUT = T} quand la date est partagée avec une autre pose.\\
XISFCONV & description de la conversion.\\
\bottomrule
\end{tabularx}
\end{center}
La colonne \texttt{t\_min} de la base vaut \texttt{DATE-OBS} − pose/2 ; \texttt{DATE-OBS} (début de pose) fait foi.
Une date d'en-tête écrite en heure locale (écart entier d'heures égal au décalage du fuseau du site avec
\texttt{MJD-OBS}/\texttt{JD}, ou pose « nocturne » datée en plein jour) est signalée dans le journal, jamais
corrigée en silence.

\subsection{Lots}
\capture{img/lots_fr.png}{0.9}
Un lot réunit ce qui s'empile ensemble : même objet, même instrument (T120, T120-bin1, T120-sCMOS, IRIS-2048,
IRIS-4096), même filtre, et même champ pour un objet fixe (toutes nuits) ou même nuit pour un objet mobile.
\begin{ccode}
07\_Nebuleuses/NGC\_6888/champ\_1\_T120/V/20250716-214427\_NGC\_6888\_V\_1s.xisf\\
01\_Asteroides/914\_Palisana/nuit\_2023-08-15\_T120/R/LOT.txt
\end{ccode}
\texttt{LOT.txt} (français puis anglais) donne objet, noms de la base, instrument, filtre, poses et pose totale,
nuits, centre moyen, angle, statut des solutions, dates douteuses et conseil d'empilement (StarAlignment ou
CometAlignment, rejet adapté au nombre de poses). Noms de fichiers : date et heure UTC du début de pose, objet,
filtre, pose ; tri par nom = tri dans le temps. Noms sûrs sous Windows. L'onglet \emph{Lots} ajoute une colonne
\emph{complet} : poses converties contre poses de la banque pour la même clé (objet, instrument, filtre, et nuit
pour un objet mobile) ; un objet fixe découpé en plusieurs champs compte ses lots ensemble.

\paragraph{Possédé $\neq$ empilable ensemble.} « 840 / 840 » dans la colonne \emph{possédé} dit que toutes les
images d'un objet sont là, pas qu'elles forment un seul empilement : la colonne \emph{lots} de la liste des objets
et son info-bulle donnent la répartition (« 840 images rangées en 9 lots : champ 1 T120 : B 200 · V 300 · R 300 ;
… — on empile lot par lot »), et le clic droit sur un objet propose \emph{Voir les lots de cet objet} (onglet Lots
filtré, bouton \emph{Tous les lots}) et \emph{Ouvrir le dossier de la cible} (double-clic : de même).

\paragraph{Pour PixInsight / N.I.N.A.} Sous le tableau des lots, un encadré donne, pour le lot choisi, ce qu'il
faut saisir dans PixInsight (ImageSolver, panneau \emph{Astrometric solution} de WBPP) ou N.I.N.A. : instrument,
\textbf{focale} (celle qui donne l'échelle mesurée : $f = 206{,}265 \times \mathrm{XPIXSZ} / \mathrm{PIXSCALE}$,
par exemple 7~234,1~mm pour le T120 au lieu des 7~200 de l'en-tête d'origine), \textbf{pixel effectif} (binning
compris : T120 iKon-L en binning 2 : 27~\textmu m), binning, \textbf{échelle} (\si{\arcsecond}\,px$^{-1}$), champ et
centre, chacun avec \emph{Copier}, plus \emph{Tout copier} (aussi au clic droit d'un lot : \emph{Copier les
paramètres d'astrométrie}). Les valeurs sont lues dans l'en-tête des images du lot, jamais codées en dur ; les
colonnes \emph{focale}, \emph{pixel} et \emph{échelle} du tableau se masquent au clic droit sur l'en-tête.
\texttt{LOT.txt} donne les mêmes valeurs. Si un logiciel affiche le double de l'échelle, il a appliqué le binning
deux fois : saisir le pixel non binné (13,5~\textmu m) avec le binning 2, ou le pixel effectif avec le binning 1.

Vérifié dans le code de PixInsight (PJSR : \texttt{AstrometricMetadata.js}, \texttt{BPP-Solver.js}) : ImageSolver
lit les propriétés \texttt{Instrument:Telescope:FocalLength} (m) et \texttt{Instrument:Sensor:XPixelSize}
(\textmu m), sinon \texttt{FOCALLEN} et \texttt{XPIXSZ}, sans multiplier par \texttt{XBINNING} ; WBPP résout le
\emph{master} intégré avec ses métadonnées si elles survivent à l'intégration, sinon avec les valeurs de son panneau,
qui restent celles du dernier instrument utilisé : y saisir focale et pixel (un fichier de réglages WBPP par
instrument évite de les ressaisir). Coupole écrit donc \texttt{FOCALLEN} accordé à l'échelle mesurée (ancienne valeur
en \texttt{HISTORY}) et les propriétés \texttt{Instrument:Sensor:XPixelSize}/\texttt{YPixelSize},
\texttt{Instrument:Camera:XBinning}/\texttt{YBinning}, \texttt{Instrument:Camera:Name},
\texttt{Instrument:Telescope:FocalLength} et \texttt{Aperture} (T120 : 1,2~m). Pour les fichiers déjà convertis :
\cmd{coupole ohp metadonnees DOSSIER --reecrire} (sans \texttt{--reecrire} : simulation) complète les en-têtes
\emph{sans toucher aux pixels} (bloc recopié octet pour octet, relu et comparé ; écriture atomique ; journal
\texttt{\_traitement/metadonnees.csv}) et note focale et pixel dans la base d'état ; \cmd{coupole ohp ranger} récrit
ensuite les \texttt{LOT.txt}. Pour vérifier un master : l'ouvrir, \emph{FITSHeader} ou \emph{Image > Properties},
chercher \texttt{XPIXSZ}, \texttt{FOCALLEN} et \texttt{Instrument:*}.

\subsection{Ouvrir une image}\label{sec:ouvrir}
Double-clic sur une image possédée (liste des images, tableau du module Qualité) : l'application associée du
système. Clic droit : \emph{Ouvrir}, \emph{Ouvrir avec} et \emph{Ouvrir l'emplacement du fichier} (fichier
sélectionné : Explorateur \texttt{/select}, Finder \texttt{open -R}, gestionnaire freedesktop par D-Bus
\texttt{FileManager1.ShowItems} — Dolphin, Nautilus… —, sinon le dossier). \emph{Ouvrir avec} ne propose que les
logiciels d'astronomie trouvés sur la machine (PixInsight, Siril, ASTAP, Aladin, N.I.N.A. : PATH, emplacements
usuels, flatpak, \texttt{/Applications}, Program Files) \emph{et} capables de lire ce fichier d'après la table
vérifiée de la section~\ref{sec:compat} ; les autres apparaissent grisés avec la raison (« ASTAP ne lit que le XISF
non compressé : choisir le format FITS dans les options du téléchargement »). Une image pas encore téléchargée : le menu commence par
\emph{Télécharger cette image}, et le double-clic propose de la télécharger (0.7.2). Une entrée désactivée porte
son motif dans son libellé (« Ouvrir avec (pas encore téléchargée) ») et le motif
complet s'affiche dans la barre d'état au survol : plus d'info-bulle par-dessus le menu (sous KDE, elle en captait
la souris).

\paragraph{Comment chaque logiciel est lancé (0.1.12).} Toujours le programme détecté, avec une \emph{liste}
d'arguments (jamais d'interpréteur de commandes) ; chaque lancement est noté dans \texttt{coupole.log} (méthode,
programme, arguments, numéro de processus ; si le programme s'arrête dans les 3~premières secondes : code de retour
et début de ce qu'il a écrit) et son résultat s'affiche dans la barre d'état : « Ouverture dans Siril… », ou
l'erreur (programme introuvable, arrêt immédiat avec son message). Jamais de silence.
\begin{itemize}
  \item \textbf{PixInsight} — Linux : le script \texttt{bin/PixInsight.sh} (l'éditeur précise que le binaire seul
  ne démarre pas sous Linux). Ce script reconstruit ses arguments par un \texttt{eval} qui n'entoure de guillemets
  que ceux qui contiennent une espace : un chemin avec \texttt{\$}, \texttt{`}, \texttt{'}, \texttt{"},
  \texttt{\textbackslash}, \texttt{(}, \texttt{*}… serait interprété ; Coupole lui passe alors un \emph{lien
  symbolique temporaire} au nom sans caractère spécial (pour enregistrer : « Enregistrer sous » dans le dossier
  d'origine). Windows : \texttt{bin\textbackslash PixInsight.exe}. macOS : \texttt{open -a PixInsight.app
  \emph{fichier}}, ou \texttt{open -n -a PixInsight.app --args -n \emph{fichier}} pour une nouvelle instance.
  \item \textbf{PixInsight déjà ouvert.} Sans option, PixInsight confie le fichier à l'instance ouverte (« Yielded
  execution to running application instance \#1 », code~0) ; si celle-ci est occupée (script en cours), \emph{rien
  ne s'ouvre}. D'où le réglage \emph{Préférences > PixInsight déjà ouvert} : « Ouvrir une nouvelle fenêtre (-n) »
  (par défaut, sûr) ou « Envoyer à la fenêtre ouverte (essai) ». Quand une instance tourne (processus détecté), le
  menu propose les deux : « PixInsight (nouvelle fenêtre) » et « PixInsight (fenêtre ouverte, essai) », celle du
  réglage d'abord ; si la sortie dit que PixInsight a cédé la main, la barre d'état le signale.
  \item \textbf{Siril} : \texttt{siril \emph{fichier}} ; en flatpak : \texttt{flatpak run --file-forwarding
  org.free\_astro.siril @@ \emph{fichier} @@} (le portail de documents ne transmet le fichier que si le bac à sable ne
  le voit pas). Le paquet Flathub voit tout le système (\texttt{--filesystem=host}) sauf \texttt{/tmp}, \texttt{/run}
  (hors \texttt{/run/media}), \texttt{/usr}… ; si un réglage \texttt{flatpak override} lui retire un dossier, la
  barre d'état le dit et donne la commande : \cmd{flatpak override --user --filesystem=/mnt/partage
  org.free\_astro.siril}.
  \item \textbf{ASTAP, Aladin} : \texttt{\emph{programme} \emph{fichier}}.
  \item \textbf{N.I.N.A.} n'ouvre pas une image passée en argument (ses seules options : profil, séquence, débogage,
  accélération graphique) : il reste grisé dans \emph{Ouvrir avec}, avec la raison ; ouvrir l'image depuis N.I.N.A.
\end{itemize}

\paragraph{Double-clic qui ouvre PixInsight sans l'image (Linux).} \texttt{xdg-open} lance l'application associée
au type du fichier par son fichier \texttt{.desktop}. Si sa ligne \texttt{Exec=} n'a pas de code \texttt{\%F}
(\texttt{\%f}, \texttt{\%u}, \texttt{\%U}), l'application démarre \emph{sans} le fichier — cas d'un
\texttt{pixinsight.desktop} écrit à la main. Coupole le détecte (\texttt{xdg-mime query filetype}, puis
\texttt{xdg-mime query default}, lecture du \texttt{.desktop}) et lance alors le logiciel directement (barre
d'état : « PixInsight lancé directement… ») ; une autre application dans ce cas n'est pas lancée, un message dit
pourquoi. Pour corriger le bureau tout entier (Coupole ne modifie pas ce fichier), dans
\texttt{\textasciitilde/.local/share/applications/pixinsight.desktop} :
\begin{verbatim}
Exec=/opt/PixInsight/bin/PixInsight.sh %F
MimeType=image/x-xisf;image/fits;
\end{verbatim}
puis \cmd{update-desktop-database \textasciitilde/.local/share/applications} et, au besoin,
\cmd{xdg-mime default pixinsight.desktop image/x-xisf}.

\paragraph{Copie d'un ancien traitement.} Une sortie produite avant Coupole (script \texttt{ohp\_xisf.py}) ou sur
une autre machine note dans sa base d'état le chemin absolu d'alors (« /srv/ancien/OHP\_DU\_ECU/… »).
Coupole retrouve chaque fichier dans le dossier de sortie courant, dans cet ordre : chemin relatif noté dans la base,
\texttt{\_traitement/journal.csv} (colonne \texttt{destination}), ancienne racine ou dossier de type, puis le nom
attendu dans le dossier du lot ; les chemins trouvés par le journal sont notés une fois pour toutes dans la base
(en fond, au premier chargement ; par la base de travail locale sur un partage). Ouvrir, \emph{Ouvrir avec},
l'emplacement, le dossier de la cible et les lots fonctionnent donc sur la copie de la banque entière (7~625 images
retrouvées en 0,34~s, notées dans la base en 0,8~s). \cmd{coupole ohp metadonnees DOSSIER --reecrire} note aussi
ces emplacements.

\subsection{Anomalies}
\capture{img/anomalies_fr.png}{1.0}
\begin{center}\small
\begin{tabularx}{\linewidth}{@{}>{\ttfamily}llX@{}}
\toprule
\normalfont\textbf{Genre} & \textbf{Action} & \textbf{Signification}\\
\midrule
meme\_fichier & écartée & même pose sous un autre chemin ou une autre orthographe\\
copie\_diurne & écartée & copie datée en plein jour d'une image bien datée\\
pixels\_identiques & écartée & mêmes pixels (SHA-1) qu'une autre image, au traitement\\
memes\_meta & signalée & mêmes date, filtre, pose, taille, nom différent\\
date\_partagee & signalée & DATE-OBS identique à celle d'une autre pose\\
date\_diurne & signalée & Soleil au-dessus de l'horizon au site (jamais pour le Soleil ni la radio)\\
champ\_incoherent & signalée & centre loin des autres poses de l'objet\\
classement & signalée & nom classé automatiquement ou non classé\\
instrument & signalée & instrument absent de la table\\
\bottomrule
\end{tabularx}
\end{center}
Au 8 octobre 2026 : 133 + 11 = 144 doublons écartés, 1~654 dates partagées, 3 poses diurnes, 92 champs
incohérents. Export CSV aux colonnes bilingues ; rien n'est supprimé.

\subsection{Carte du ciel}
\capture{img/ciel_fr.png}{0.95}
Projection d'Aitoff, ascension droite croissante vers la gauche, 12~h au centre ; un point par objet fixe (taille :
nombre d'images) et par nuit pour les objets mobiles ; écliptique en tirets, plan galactique en pointillés.

\subsection{Fiche en ligne (facultative)}\label{sec:fiche}
\capture{img/fiche_fr.png}{0.95}
Si Internet est disponible, l'onglet \emph{Fiche en ligne} (ou le bouton du même nom sous le catalogue) affiche des
informations publiques sur l'objet choisi :
\begin{itemize}
\item ciel profond et étoiles : \textbf{SIMBAD} (CDS, Strasbourg) --- type, coordonnées ICRS, magnitudes, parallaxe
(et la distance qu'elle donne quand elle est assez précise), distance mesurée préférée avec sa méthode et sa
référence, vitesse radiale et redshift, taille angulaire, types spectral et morphologique, autres noms ; liens
vers la page SIMBAD, \textbf{Aladin Lite} (image du ciel centrée sur l'objet) et, pour les galaxies, \textbf{NED} ;
si SIMBAD ne connaît pas le nom, \textbf{Sesame} (CDS) interroge aussi NED et VizieR ;
\item astéroïdes, comètes, transneptuniens : base des petits corps du \textbf{JPL} (SBDB) --- classe orbitale,
géocroiseur ou non, éléments (a, e, i, q, Q, $\Omega$, $\omega$, période, époque), diamètre, albédo, période de
rotation, magnitude absolue ;
\item planètes et désignations NEOCP : pas de fiche (message explicatif).
\end{itemize}
\textbf{L'identifiant retenu est toujours affiché} : SIMBAD rattache par exemple « NGC 6888 » à son étoile
Wolf-Rayet HD 192163 ; la fiche le signale. Un redshift positif s'envoie au module Cosmologie (bouton \emph{Envoyer
ce redshift au module Cosmologie}).

\paragraph{Politesse et hors ligne.} Rien ne part tant que l'onglet n'est pas affiché ; une seule requête par
objet choisi (après 0,6~s sans autre changement de sélection), au plus une par seconde et par service, délai de
8~s, User-Agent de Coupole. Les réponses sont gardées 30~jours dans le cache de Coupole (\texttt{enligne.json},
avec leur date ; un « introuvable » un jour). Hors ligne, la dernière fiche connue s'affiche avec sa date, sinon un
message clair ; le reste de l'application n'en dépend jamais. La case \emph{Interroger les services en ligne} coupe
tout envoi. Les adresses (\texttt{simbad.base}, \texttt{simbad.page}, \texttt{aladin.page}, \texttt{ned.page},
\texttt{sbdb.page}, \texttt{resolution.sbdb}, \texttt{resolution.sesame}) se règlent comme les autres sources.

% =====================================================================
\section{Partages du DU}\label{sec:partages}
\capture{img/partages_contenu_fr.png}{1.0}
Les défis et les TP du DU se font souvent dans un \emph{dossier partagé} (Nextcloud de l'Observatoire) : un lien
« \texttt{https://…/s/…} », dans lequel chacun dépose ses images, ses vidéos et parfois une ligne dans un tableau.
Ce module rassemble tout au même endroit : plusieurs partages mémorisés (nom libre), leur structure, les aperçus, les
nouveautés, le téléchargement, le contrôle des tableaux de mesures, les outils d'analyse et le dépôt de votre propre
mesure. \textbf{Aucun lien n'est livré avec Coupole} : vous collez celui que l'équipe du DU vous a donné ; il reste
sur votre ordinateur (fichier \texttt{partages.json} lisible par vous seul) et n'apparaît dans aucun rapport.

\paragraph{Banque OHP ou partage ?} La Banque OHP est une archive professionnelle (service TAP, métadonnées
normalisées, FITS des télescopes T120 et IRIS). Un partage est un dossier collaboratif d'étudiants : formats et
métadonnées hétérogènes, saisis à la main, à valider. Coupole devine la structure (dossiers datés, observateurs,
séries, tableaux), dit d'où vient chaque information (en-tête, tableau, nom de fichier) et signale les doutes. La
présentation est la même que celle de la Banque : liste à gauche, contenu à droite, bouton bleu
« Télécharger (N · taille) », clic droit qui commence par « Télécharger », guide en trois étapes au premier lancement
(menu \emph{Aide > Guide des partages}), aperçus (FITS, JPEG, PNG, TIFF, première image d'une vidéo SER ou AVI) et
bandeau des nouveautés depuis votre dernière visite.

\paragraph{Ce que Coupole devine.} Le genre de chaque fichier ; l'instant d'après le nom (« \texttt{AAAA-MM-JJ-HH-MM-SS} »,
« \texttt{JJ-MM-AA-…} » du sujet, « \texttt{IMG20260618223000} »…) — un nom qui se lit de deux façons
(17-10-26 : 17 octobre 2026 ou 26 octobre 2017) est tranché par la date de dépôt, et une heure lue dans un nom est
marquée « d'après le nom » (souvent l'heure locale) ; l'observateur (fin du nom, ou dossier
« date – lieu – nom ») ; la cible (ISS, Lune…) ; les prises (variantes FITS/PNG/TIFF d'une même photo) et les séries
(poses espacées de moins de 2 min). Les en-têtes FITS se lisent \emph{sans télécharger l'image} (requête partielle :
les 34 premiers ko) : lieu (SITELAT/SITELONG), heure de début et milieu de pose.

\capture{img/partages_mesures_fr.png}{1.0}
\paragraph{Tableaux de mesures.} Les classeurs (xlsx, csv) sont lus \emph{sans être modifiés}. Formats normalisés
(« 2,45° », « 48°47′05″ N », « +44:13:51:00 », nombres de jours d'Excel, dates américaines) et contrôles ligne par
ligne, avec la preuve quand il y en a une : latitude et longitude échangées (certain si $|$latitude$| > 90^\circ$ ;
prouvé par l'en-tête du fichier, par la position de la Lune — hauteur et azimut notés ne collent qu'avec les
coordonnées échangées — ou par la comparaison aux lieux sûrs) ; heure locale au lieu du TU (l'en-tête ou la Lune
s'accordent avec l'heure moins 1 ou 2~h) ; date ambiguë tranchée par le nom du fichier ; champs manquants, fichier
introuvable. Les corrections proposées (marquées « * ») servent aux calculs si vous le voulez, jamais au tableau
d'origine ; « Exporter en CSV » donne le tableau normalisé.

\capture{img/partages_analyse_fr.png}{1.0}
\paragraph{Outil « Parallaxe (Lune, ISS) ».} Les outils d'analyse se branchent selon le contenu du partage (le
premier est celui-ci ; d'autres s'ajoutent sans toucher au reste). \emph{Chercher les paires} : même cible, deux
observateurs éloignés (base $\ge 100$~km), presque au même instant (ISS $\le 5$~s, Lune $\le 60$~s). \emph{Mesurer}
(clic droit sur une image) : centre de la Lune ajusté sur le limbe (cercle de Kåsa, le terminateur est écarté ; une
Lune saturée ou un disque de taille anormale est signalé et l'incertitude élargie), objet mobile d'une série de poses
courtes ou trait d'une pose longue pour l'ISS ; la direction $(\alpha, \delta)$ vient de la solution astrométrique de
l'image. \emph{Calculer}, puis le rapport HTML ou PDF.
\begin{itemize}
\item \textbf{Lune} : méthode vectorielle topocentrique (positions WGS84 des deux lieux dans le même repère, point le
plus proche des deux visées). Les directions mesurées sur une photo sont \emph{astrométriques} : la Lune où elle était
il y a $\approx 1{,}2$~s, vue d'une Terre qui a bougé depuis ; le temps de lumière n'étant pas le même pour les deux
observateurs, Coupole ramène chaque visée à la direction géométrique avant de trianguler (sans cela, $\approx 20$~km
d'erreur pour 600~km de base). Comparaison à JPL DE440 (fichier de 32~Mo téléchargé une fois) ou à Meeus, budget
d'erreur (direction, lieu, heure).
\item \textbf{ISS} : triangulation des deux traces à l'instant commun, estimation du décalage d'horloge entre les
observateurs, méthode des plans (sans horloge commune), comparaison au TLE (SGP4).
\item Méthodes à un observateur : parallaxe diurne de la Lune (plusieurs heures de mesures), vitesse angulaire ou
orbite circulaire de l'ISS.
\end{itemize}

\paragraph{Vérification.} Les calculs sont comparés à des références indépendantes (dossier
\texttt{outils/verif\_parallaxe/}) : skyfield et JPL DE440 (position de l'observateur à 2~m près, directions de la Lune
à 0,05″ avec DE440, distance de la Lune retrouvée à 0,5~km par la triangulation, altitude de l'ISS à 20~m, prédiction
TLE à 10~m) et SageMath à 200 bits (ellipsoïde WGS84 dans les deux sens, point le plus proche de deux droites,
réfraction, cercle de Kåsa, orbite au zénith).

\capture{img/partages_depot_fr.png}{1.0}
\paragraph{Déposer sa mesure.} Uniquement sur votre action, après une confirmation qui liste exactement ce qui part.
Chaque fichier prend un nom libre (« \texttt{\_2} » si le nom existe) et part avec la condition « le fichier ne doit
pas exister » ; votre ligne s'ajoute dans la première ligne libre du tableau (aucune cellule remplie n'est jamais
remplacée, styles et autres feuilles recopiés octet pour octet) avec la condition « le tableau n'a pas changé depuis
sa lecture » ; une copie locale du tableau est gardée avant l'ajout ; après trois conflits, la ligne part dans un
petit fichier « \texttt{ligne\_<nom>\_<date>.csv} » à côté. Le client WebDAV de Coupole n'a \emph{aucune} méthode de
suppression, de renommage ou de déplacement : même si le lien le permet (Coupole vous le signale), il ne touche jamais
aux fichiers des autres.

\paragraph{Astrométrie des images téléchargées.} La solution de l'en-tête est gardée ; sinon, au choix : ASTAP ou
astrometry.net installés sur votre ordinateur, puis, en dernier recours, le service en ligne nova.astrometry.net
(section~\ref{sec:nova}).

\begin{ccode}
coupole partages ajouter "Défi parallaxe" "https://…/s/…"\\
coupole partages explorer "Défi parallaxe"\\
coupole partages tableau "Défi parallaxe" --json\\
coupole partages paires "Défi parallaxe" --entetes\\
coupole partages mesurer image.fits --cible lune
\end{ccode}

% =====================================================================
\section{Archives des observatoires}\label{sec:archives}
\capture{img/module_archives_fr.png}{1.0}
Les grands observatoires publient leurs images une fois écoulée la période réservée à l'équipe qui les a
demandées (souvent douze mois). Le module cherche ces images publiques, actuelles et passées, les télécharge
poliment, en extrait l'image scientifique et aide à les aligner et à les colorer, comme la Banque OHP le fait pour
les stages. Méthode détaillée et sources : \texttt{docs/archives\_methode.md}.

\begin{center}\small
\begin{tabularx}{\linewidth}{@{}llX@{}}
\toprule
\textbf{Archive} & \textbf{Étape} & \textbf{Ce qui est proposé}\\
\midrule
MAST (STScI) & 1 & Hubble (\texttt{\_drz}, \texttt{\_drc}), JWST (\texttt{\_i2d}), GALEX (\texttt{-int}) ; service TAP officiel, vue \texttt{dbo.obspointing}\\
ESO & 1 & produits réduits « Phase 3 » (VISTA, VST, HAWK-I, FORS, MUSE…) ; TAP ObsCore\\
IRSA (NASA/IPAC) & 1 & Spitzer (mosaïques SEIP), WISE (Atlas AllWISE), 2MASS (Atlas) ; SIA 2\\
NOIRLab & 2 & piles et images rééchantillonnées (DECam, Mosaic…)\\
Keck (KOA) & 2 & poses brutes NIRC2, OSIRIS, MOSFIRE (aucun produit final d'imagerie)\\
SDSS & 2 & images calibrées u, g, r, i, z\\
Gemini & 2 & \textbf{compte exigé} même pour les données publiques : signalé, rien n'est téléchargé\\
SMOKA (Subaru) & 2 & formulaire web seulement : signalé, non pris en charge\\
OPUS (PDS) & 3 & Voyager ISS (\texttt{\_GEOMED}, distorsion corrigée), Cassini ISS (\texttt{\_CALIB})\\
PDS Imaging Atlas & 3 & JunoCam (RDR, bandes de \emph{framelets})\\
\bottomrule
\end{tabularx}
\end{center}

\subsection{Ce que vous téléchargez, par rapport à votre télescope}
Avec votre télescope, vous prenez des poses unitaires (« lights »), vous les calibrez (offset, noir, PLU), puis vous
les alignez et les empilez en un \emph{master}. Les archives livrent les mêmes étapes, déjà faites :
\begin{center}\small
\begin{tabularx}{\linewidth}{@{}lX@{}}
\toprule
\textbf{Nature (colonne)} & \textbf{Équivalent amateur}\\
\midrule
pose unitaire · calibrée & une de vos « lights » calibrées (Cassini, Voyager, JunoCam, SDSS, Keck) ; badge bleu \emph{Pose calibrée}\\
produit combiné · produit final & votre master empilé : poses alignées, rayons cosmiques rejetés, distorsion corrigée (Hubble \texttt{drz}/\texttt{drc}, JWST \texttt{i2d}, ESO Phase 3) ; badge vert \emph{Produit final}\\
mosaïque & plusieurs visites ou champs réunis sur une grille commune (Hubble MVM, JWST \texttt{c…}, Spitzer SEIP, WISE)\\
image de détection & somme de tous les filtres d'une visite, faite pour les catalogues : pas un filtre, écartée par défaut\\
\bottomrule
\end{tabularx}
\end{center}
La colonne \emph{pose totale} donne la somme des poses et leur nombre quand il est connu (« 8 × 650 s = 5~200 s
(1~h~27~min) ») : à comparer à vos nuits. Le niveau de calibration de la fiche est expliqué en clair (2 = calibré
pose par pose, 3 = produit final combiné, 4 = produit d'une équipe scientifique).

\paragraph{Unités.} Hubble donne des électrons par seconde, JWST des MJy/sr : pas des ADU. Les valeurs sont
minuscules (0,01 à 10) ; l'image paraît donc noire tant qu'on ne l'étire pas (STF de PixInsight, auto-ajustement de
Siril). Rien n'est perdu : Coupole garde les valeurs physiques et les constantes photométriques.

\paragraph{Hubble : drz, drc, SVM, MVM et « skycell ».} Pour les mêmes poses, MAST publie jusqu'à quatre produits
(vérifié sur Arp~273, programme 12326 : 17 fichiers pour 3 filtres) :
\begin{itemize}
\item \texttt{ipppssoot\_drz} : pipeline standard, une association de poses d'une visite et d'un filtre ;
  \texttt{\_drc} : le même à partir de poses corrigées de l'efficacité de transfert de charge (CTE, ACS/WFC et
  WFC3/UVIS) — \emph{HST Data Handbook}, « Types of ACS files » ;
\item \texttt{hst\_<prog>\_<visite>\_…\_drc} : \textbf{mosaïque d'une visite} (SVM) ; tous les filtres de la visite
  alignés entre eux et sur Gaia, sur la même grille ; \texttt{…\_total\_…} est l'image de détection ;
\item \texttt{hst\_skycell-pPPPPxXXyYY\_…\_all\_drc} : \textbf{mosaïque multi-visites} (MVM) : toutes les poses
  publiques du filtre dans une \emph{cellule du ciel} (« skycell »), une grille fixe de 0,2° de côté à 0,04″ par
  pixel, remplie seulement là où Hubble a observé. La cellule est rognée au plus petit rectangle qui contient
  \emph{toutes} les poses de la cellule, tous filtres confondus (DrizzlePac, « Multi-visit Mosaic Processing ») : un
  filtre peut n'en occuper qu'un coin. Exemple réel : M~16, ACS F775W, 19~806 × 15~741 pixels (2,49~Go chez MAST)
  pour une bande de données de 762 × 5~217 pixels.
\end{itemize}
\textbf{Un produit logique par filtre} : Coupole retient la MVM quand elle existe (toutes les poses, temps total le
plus long, CTE corrigée, grille commune à tous les filtres de la cellule : couleur directe — MAST, « Multi-Visit
Mosaics from Hubble Now Available », 2022) ; sinon, par visite, la SVM ; sinon le produit du pipeline standard
(WFPC2 par exemple). Pour JWST : la mosaïque de plusieurs observations (\texttt{c…}), sinon l'observation au plus
long temps de pose. Les autres produits des mêmes poses deviennent des \emph{variantes}, cachées par défaut (case
\emph{Montrer les variantes}) ; la fiche dit pourquoi chaque produit est retenu ou mis de côté. Limites documentées
des MVM : décalage résiduel de quelques pixels entre visites, objets mobiles dédoublés, WFPC2 absent.

\subsection{Chercher}
\capture{img/archives_produits_fr.png}{1.0}
Le champ \emph{Objet} accepte :
\begin{itemize}
\item un nom ou un numéro de catalogue, résolu par SIMBAD puis Sesame (qui interroge aussi NED et VizieR) ; les
  écritures courantes sont réécrites (« m16 » → M~16, « ngc6611 », « arp273 », « sh2-155 ») et les noms usuels
  français et anglais reconnus (« Piliers de la Création », « Nébuleuse de la Carène », « Quintette de Stephan »,
  « HUDF ») ;
\item un \textbf{corps du Système solaire}, en français ou en anglais, sans souci de casse ni d'accents (jupiter,
  Saturne, Encelade, Japet, Phœbé…). SIMBAD ne connaît pas les planètes et une planète bouge : le corps est cherché
  \textbf{par nom de cible} dans MAST (Hubble, JWST : \texttt{target\_name} « JUPITER », « JUPITER-MAPS »… avec le
  drapeau de cible mobile), Voyager/Cassini (OPUS) et Juno (PDS), même si ces archives ne sont pas cochées ;
\item des coordonnées en degrés (\texttt{83.82 -5.39}) ou sexagésimales (\texttt{05:35:17 -05:23:28}).
\end{itemize}
Des suggestions s'affichent pendant la frappe (corps, noms usuels). Un objet introuvable donne des propositions et
un conseil. Rayon : 3~minutes d'arc par défaut. Filtres : instrument (morceau du nom : \texttt{NIRCAM},
\texttt{ACS}), filtre (\texttt{F444W}), dates. Cases : \emph{Produits finaux seulement}, \emph{Données publiques
seulement}, \emph{Montrer les variantes}. Les archives répondent en parallèle ; une archive injoignable, qui exige
un compte ou n'a pas d'interface est signalée sans empêcher les autres. Une réponse tronquée par la limite est
signalée ; quand une archive sert deux missions (OPUS : Voyager et Cassini), la limite est partagée entre elles (en
0.2.1, « saturne » ne rendait que Voyager).

\paragraph{Taille, fiche et aperçu.} La taille de chaque fichier s'affiche d'elle-même : annoncée par l'archive, ou
mesurée en fond par une requête d'un octet (60 au plus d'un coup, les autres extrapolées et marquées « ≈ ») ; le total
de la sélection est sous le tableau. La fiche donne mission, instrument, filtre (avec sa longueur d'onde et son
équivalent amateur), nature, pose totale, dates, programme, investigateur, niveau de calibration, fichier, crédit et
conditions d'usage. L'aperçu est la vignette officielle de l'archive, lue en entier, réduite sans être décompressée en
pleine taille et recadrée sur la zone utile (les JPEG de MAST sont en pleine résolution, jusqu'à 19~806 × 15~741
pixels, fond blanc là où il n'y a pas de données) ; à défaut une vignette générée depuis le fichier préparé ;
sinon la raison (« l'archive ne fournit pas de vignette… »).

\subsection{Télécharger et préparer}
Choisir des lignes (Ctrl+clic, Maj+clic), puis le bouton bleu \emph{Télécharger la sélection (N · taille)}, juste sous
le tableau (0.7.2 : grisé avec la raison au survol tant que rien n'est choisi, « Déjà téléchargé ✓ » si tout est là).
Le clic droit commence par \emph{Télécharger cette observation} (ou \emph{ces N observations}), puis \emph{Ouvrir},
\emph{Ouvrir avec}, \emph{Ouvrir l'emplacement du fichier} et la page de l'archive ; un double-clic sur une ligne pas
encore téléchargée propose de la télécharger. Au-delà de 2~Go (réglage \texttt{archives\_seuil\_go}),
confirmation. Débit libre par défaut (plafond : Préférences, ou réglage \texttt{archives\_debit\_mo\_s}), deux fichiers à la fois, reprise
d'un fichier interrompu. La barre avance en \textbf{deux phases} : téléchargement (jusqu'à 70~\%, d'après les octets
reçus), puis préparation (jusqu'à 100~\%) ; dessous : fichiers téléchargés et préparés, octets reçus sur le total,
débit, temps restant de l'ensemble et du fichier en cours. À la fin, l'état est écrit en clair : « Terminé :
N fichiers, X téléchargés, Y gardés, Z erreurs — Ouvrir le dossier ».

Chaque image est préparée puis rangée \textbf{par champ}, comme dans la Banque OHP :
\texttt{<dossier>/Archives/<champ>/<cible>\_<instrument>\_<filtre>\_<date>\_<type>\_<n>x<pose>s.xisf}, par exemple
\texttt{Arp273\_WFC3-UVIS\_F390W\_2010-12-17\_MVM\_8x1200s.xisf} ou
\texttt{HUDF\_WFC3-IR\_F160W\_2009-11-03\_HLSP\_26x1361s.xisf} (type : MVM, SVM-v21 pour la visite 21, drz, drc,
pose, i2d, HLSP…). L'identifiant de l'archive reste dans l'en-tête (\texttt{ARCHID}, \texttt{ORIGFILE}) et dans la
fiche. Préparation :
\begin{itemize}
\item extension scientifique extraite (\texttt{SCI} ; à défaut la première image), un fichier par détecteur pour une
  pose à plusieurs détecteurs de l'ESO (\texttt{\_d1}, \texttt{\_d2}…) ;
\item \textbf{recadrage sur la zone utile} : plus petit rectangle des pixels valides, la WCS suit
  (\texttt{CRPIX} décalés ; \texttt{CROPX0}, \texttt{CROPY0}, \texttt{ORIGNAX1}, \texttt{ORIGNAX2} gardent la trace) ;
  tailles avant et après dans la barre d'état, la fiche et le journal ;
\item en-tête : WCS (SIP comprise), \texttt{BUNIT}, constantes photométriques, télescope, instrument, filtre, dates,
  temps de pose, nombre de poses (\texttt{NPOSES}), programme, investigateur, \texttt{CREDIT} et conditions d'usage ;
\item pixels sans donnée restants (NaN) remplacés par 0 (\texttt{NANCOUNT}, \texttt{NANFILL}) ;
\item XISF Float32 (valeurs physiques gardées) ou FITS Float32.
\end{itemize}
\textbf{Par défaut, rien d'autre que le produit préparé} : le fichier d'origine (case \emph{Garder aussi le fichier
d'origine}, rangé dans \texttt{<champ>/\_originaux}) et le masque des pixels sans donnée (case \emph{Masque…},
entiers 8 bits compressés, écrit après recadrage) sont des options. Sans masque, l'alignement considère les 0 exacts
comme « sans donnée ». L'état est tenu dans \texttt{Archives/\_etat/etat.sqlite} et le journal bilingue dans
\texttt{Archives/\_etat/JOURNAL.txt}.

\paragraph{Réorganiser les téléchargements.} Le bouton \emph{Réorganiser les téléchargements…} (et
\cmd{coupole archives reorganiser}) reprend ce qui a été téléchargé avec la 0.2.0 ou la 0.2.1 : noms lisibles,
rangement par champ, recadrage, suppression des masques pleine taille, variantes et images de détection dans
\texttt{<champ>/\_variantes}. Une confirmation montre ce qui va changer ; rien n'est jamais écrasé (un nom pris
reçoit \texttt{\_2}) ; chaque déplacement est écrit dans le journal et la base d'état suit. Essai réel sur 18 fichiers
d'une base écrite sous Windows : 18 rangés, 11 recadrés, 16 masques supprimés, 1~641 → 910~Mo.

\subsection{Catalogue des cibles}
\capture{img/archives_catalogue_fr.png}{1.0}
L'onglet \emph{Catalogue} liste les cibles d'une archive \textbf{sans critère}, les plus observées d'abord, par pages
de 100, filtrables par type d'objet (galaxie, amas, nébuleuse, étoile, Système solaire, champ profond) et par nom ;
double-clic : la recherche est lancée. Requêtes agrégées (\texttt{GROUP BY target\_name}) mesurées : Hubble 12 à
17~s (limité aux mosaïques multi-visites : la table entière dépasse le délai du service), JWST 34~s, ESO 2~s,
Voyager, Cassini (OPUS) et Juno (PDS) en moins d'une seconde ; les listes sont gardées 30 jours (\emph{Actualiser}
pour les relire). IRSA, NOIRLab, Keck et SDSS n'ont pas de liste utile (relevés du ciel entier, champs numérotés,
poses brutes) : l'onglet l'explique. \emph{Programmes célèbres} : Hubble Deep Field (6337), Hubble Ultra Deep Field
(9978, 10086), HUDF09 (11563), HUDF12 (12498), Piliers de la Création 1995 (5773), 2014 (13926) et JWST 2022 (2739),
Carène (10475 ; JWST 2731), Quintette de Stephan (2732), Anneau austral (2733), Arp~273 (12326) : la recherche ne
garde alors que les produits du programme.

\subsection{Hubble Ultra Deep Field}
\capture{img/archives_hudf_fr.png}{1.0}
L'onglet guide vers le HUDF (centre 3~h~32~min~39~s, −27°~47′~29″) : les poses calibrées une à une (\texttt{flt},
\texttt{flc}) se comptent par centaines par filtre ; les équipes ont publié des \textbf{produits de haut niveau}
(HLSP) déjà empilés et alignés entre filtres : XDF (Illingworth et al. 2013 : toutes les poses ACS et WFC3/IR de 2002
à 2012, environ 1,8 million de secondes), HUDF09, HUDF12. Temps de pose vérifiés dans MAST (somme des produits du
pipeline) : programmes 9978 + 10086, ACS/WFC F435W 134~880~s, F606W 135~320~s, F775W 347~110~s, F850LP 346~620~s ;
HUDF12 (readme de l'équipe) F105W 198~423~s, F140W 82~676~s, F160W 71~652~s. Préréglage « couleur en 3 filtres » :
XDF F435W, F775W et F160W à 60~mas, 3 × 110~Mo, même grille (\cmd{coupole archives hudf --telecharger}).

\subsection{Instruments et filtres}
L'onglet \emph{Instruments et filtres} donne, pour chaque instrument (Hubble ACS/WFC, WFC3 UVIS et IR, WFPC2 ; JWST
NIRCam, MIRI, NIRISS ; ESO HAWK-I, FORS2, MUSE, VIRCAM, OmegaCAM ; Spitzer IRAC et MIPS ; WISE ; 2MASS ; SDSS ;
Keck NIRC2 ; Voyager et Cassini ISS ; JunoCam) son champ, son échelle et son domaine, et pour chaque filtre sa
longueur d'onde, sa largeur, son usage et son \textbf{équivalent amateur}. Survoler un nom de filtre dans le tableau
des résultats en donne le résumé. Le nom code le filtre : F = filtre, nombre = longueur d'onde (nm pour Hubble,
centièmes de µm pour JWST), lettre = largeur (W large, M moyen, N étroit, LP passe-haut, X très large).
\begin{center}\small
\begin{tabularx}{\linewidth}{@{}llllX@{}}
\toprule
\textbf{Filtre} & \textbf{Instrument} & \textbf{λ (nm)} & \textbf{Largeur} & \textbf{Équivalent amateur}\\
\midrule
F435W & ACS/WFC & 430 & 104 & B\\
F555W, F606W & ACS/WFC & 535, 591 & 119, 234 & V ; F606W, très large, s'approche d'une luminance\\
F656N, F657N & WFC3/UVIS & 656 & 2, 12 & Hα (vos filtres de 3 à 7~nm)\\
F502N & WFC3/UVIS & 501 & 7 & OIII\\
F673N & WFC3/UVIS & 677 & 12 & SII\\
F775W, F814W & ACS/WFC & 776, 833 & 153, 251 & proche infrarouge (filtre « IR-pass »)\\
F160W & WFC3/IR & 1~537 & 268 & aucun : au-delà de 1~µm, un capteur CMOS ne voit plus rien\\
F444W, F770W & NIRCam, MIRI & 4~402, 7~639 & 1~024, 1~950 & aucun : infrarouge, en fausses couleurs\\
\bottomrule
\end{tabularx}
\end{center}
Sources : manuels des instruments WFC3 (tableaux 6.2 et 7.2, longueur d'onde pivot) et ACS (tableau 5.1, longueur
d'onde centrale), JDox (NIRCam, MIRI, NIRISS), pages des instruments de l'ESO, manuels Spitzer, suppléments
explicatifs WISE et 2MASS, SDSS, PDS Ring-Moon Systems Node, Hansen et al. 2017 (JunoCam) ; valeurs approchées
marquées « ≈ » dans l'onglet.

\paragraph{Formats planétaires.} Les produits de Voyager et de Cassini sont des fichiers VICAR accompagnés d'un label
PDS3 ; ceux de Juno des fichiers PDS3 ; le PDS4 (label XML) est lu aussi. Coupole lit ces formats lui-même (entiers
8 à 32 bits, réels IEEE, ordre des octets, préfixes de ligne, mise à l'échelle), retourne les lignes pour que l'image
s'affiche comme dans l'archive et écrit FITS ou XISF. Particularité : le label PDS3 des images calibrées de Cassini
place l'image une ligne trop tôt (une ligne binaire VICAR la précède) ; les données sont lues d'après l'en-tête
VICAR. Non gérés (message clair) : fichiers compressés (EDR de Voyager), réels VAX, entrelacement BIL/BIP à
plusieurs bandes, cubes ; une image JunoCam est convertie telle quelle (bandes de 128 lignes par filtre), sans
reconstruction géométrique.

\subsection{Aligner et colorer}
\capture{img/archives_alignement_fr.png}{1.0}
Les images d'archive ont une astrométrie exacte : on ne cherche pas d'étoiles (elles changent d'aspect d'un
instrument et d'une longueur d'onde à l'autre), on \textbf{rééchantillonne} chaque image sur une grille commune :
celle d'une image (souvent le filtre le plus fin) ou une \emph{grille optimale} qui les couvre toutes, au pas le plus
fin, éventuellement agrandi (\emph{Pixels ×}). Méthodes : bilinéaire (défaut, rapide), adaptative (anti-crénelage,
meilleure quand les pas diffèrent beaucoup), exacte (recouvrement exact des pixels, lente ; sous 0,05″ par pixel,
l'adaptative la remplace), plus proche voisin. Adaptative et exacte demandent la bibliothèque facultative
\texttt{reproject} (\cmd{pip install ".[alignement]"}) ; sans elle, la bilinéaire est faite par scipy. Une quantité
\textbf{par pixel} (électrons/s de Hubble, nanomaggies de SDSS) est multipliée par le rapport des surfaces de pixel ;
une brillance de surface (MJy/sr de JWST) est gardée telle quelle. Recadrage facultatif sur le plus grand rectangle
couvert par toutes les images (une colonne morte au milieu d'une mosaïque ne le réduit pas). La précision vaut celle
des astrométries des archives : un écart résiduel de l'ordre du pixel peut subsister entre missions ; l'alignement
par étoiles de PixInsight ou de Siril sur les fichiers rééchantillonnés n'a plus alors qu'une translation à corriger.

Sortie dans \texttt{Archives/\_alignes/<date>/} : un fichier par filtre (\texttt{\_aligne}), le masque commun, un
aperçu couleur (\texttt{apercu\_couleur.png}, étirement asinh) et \texttt{composition.txt}. Les filtres sont rangés
par longueur d'onde (\emph{ordre chromatique} des images de presse de Hubble et de JWST) : trois filtres donnent
bleu, vert et rouge ; davantage, des teintes réparties du bleu au rouge. \texttt{composition.txt} donne le tableau des
couleurs et l'expression PixelMath (identifiants d'image tels que PixInsight les crée à l'ouverture) ; dans Siril,
la composition RVB se fait avec les mêmes fichiers, déjà alignés.

\subsection{Ligne de commande}
\begin{verbatim}
coupole archives liste
coupole archives chercher "M 16" --rayon 6 --mission JWST --instrument MIRI
coupole archives chercher arp273                 # un produit par filtre
coupole archives chercher arp273 --variantes     # toutes les variantes
coupole archives chercher saturne                # MAST, Voyager, Cassini par nom de cible
coupole archives telecharger "M 16" --mission JWST --instrument MIRI --dest ~/Ciel
coupole archives telecharger M51 --archive sdss --filtre g --filtre r --filtre i --dest ~/Ciel --oui
coupole archives catalogue                       # programmes célèbres
coupole archives catalogue jwst --type nebuleuse --page 2
coupole archives chercher HUDF --rayon 2 --programme 9978 --programme 10086
coupole archives hudf --telecharger --dest ~/Ciel
coupole archives instruments F656N
coupole archives reorganiser --dest ~/Ciel --simuler
coupole archives preparer fichier_i2d.fits --format fits
coupole archives aligner A.xisf B.xisf C.xisf --dest ~/Ciel/Piliers --optimale
coupole archives bilan --dest ~/Ciel
\end{verbatim}
\cmd{--id} garde des observations précises (identifiant, avec ou sans le préfixe de l'archive, ou début du nom de
fichier) ; \cmd{--tous-niveaux} montre aussi les poses calibrées et brutes ; \cmd{--garder-original} et
\cmd{--masque} gardent le fichier d'origine et le masque ; \cmd{--json} et \cmd{--csv} pour les scripts. Au-delà du
seuil, \cmd{telecharger} demande \cmd{--oui}.

\paragraph{Crédits.} Toute image publiée à partir de ces données doit citer l'archive et la mission : « NASA/ESA/CSA,
STScI » (JWST), « NASA/ESA, STScI » (Hubble), « ESO » et l'identifiant du programme, « NASA/JPL-Caltech » (Spitzer),
« NASA/JPL-Caltech/UCLA » (WISE), « 2MASS/UMass/IPAC-Caltech/NASA/NSF », NOIRLab, KOA, SDSS, et pour les sondes
« NASA/JPL-Caltech » (Voyager), « NASA/JPL-Caltech/Space Science Institute » (Cassini), « NASA/JPL-Caltech/SwRI/MSSS »
(Juno). Coupole l'écrit dans la fiche et dans l'en-tête de chaque fichier (\texttt{CREDIT}).

% =====================================================================
\section{ASTAP (facultatif)}
Sans ASTAP, Coupole fait tout le reste. ASTAP (Han Kleijn, licence MPL~2.0, \url{https://www.hnsky.org/astap.htm})
ajoute la vérification indépendante de chaque solution. Il faut le programme et un catalogue d'étoiles.

\capture{img/astap_fr.png}{0.72}

\paragraph{Catalogue.} Le site officiel indique que pour un champ de 0,6° ou plus, D05, D20, D50 et D80 conviennent,
et que les catalogues denses servent les petits champs. Les champs de la banque vont de 10,7′ (0,18°) à 31,4′
(0,52°) : \textbf{D80} est conseillé (environ 1,2 Go ; catalogue du traitement de référence, 6~888 solutions
confirmées sur 7~625) ; D50 (environ 0,9 Go) convient aussi pour IRIS.

\begin{center}\small
\begin{tabularx}{\linewidth}{@{}lX@{}}
\toprule
\textbf{Système} & \textbf{Programme (liens de la page officielle)}\\
\midrule
Windows x64 & \texttt{astap\_setup.exe} puis \texttt{d80\_star\_database.exe} (même dossier, \texttt{C:\textbackslash Program Files\textbackslash astap}) ; facultatif : \texttt{astap\_command-line\_version\_win64.zip}\\
Windows ARM64 & \texttt{astap\_command-line\_version\_win11\_aarch64.zip} dans \texttt{C:\textbackslash astap} ; catalogue : \texttt{d80\_star\_database.exe} (installeur x86, émulé) vers ce dossier, ou \texttt{d50\_star\_database.zip} (pas de D80 en zip)\\
macOS Apple Silicon & \texttt{astap\_M1.pkg}, \texttt{d80\_star\_database.pkg} ; une fois : \texttt{codesign --force -s - /Applications/ASTAP.app/Contents/MacOS/astap}\\
macOS Intel & \texttt{astap.pkg}, \texttt{d80\_star\_database.pkg}\\
Linux Debian/Ubuntu & \texttt{astap\_amd64.deb} (ou \texttt{astap\_aarch64.deb}), \texttt{d80\_star\_database.deb} (dans \texttt{/opt/astap})\\
Linux Fedora/openSUSE & \texttt{astap\_amd64.rpm} ; D80 extrait du \texttt{.deb} (ci-dessous) ou D50 en zip dans \texttt{/opt/astap}\\
Linux Arch/Manjaro & \texttt{astap\_amd64\_gtk3.pkg.tar.zst} (ou GTK2) ; D80 : paquet AUR \texttt{d80-star-db-astap} ou extrait du \texttt{.deb} ; D50 : \texttt{d50\_star\_database.pkg.tar.zst}\\
Autres Linux & \texttt{astap\_amd64.tar.gz} ou \texttt{astap\_command-line\_version\_Linux\_amd64.zip} ; pas d'AppImage officielle\\
\bottomrule
\end{tabularx}
\end{center}

\paragraph{Pas de D80 en zip.} Le lien « D80 (zipped) » de la page officielle est mort (404, vérifié le 9/10/2026) :
le D80 n'existe qu'en \texttt{.exe}, \texttt{.pkg} et \texttt{.deb}. Hors Debian, l'extraire du \texttt{.deb} sans
toucher aux droits de \texttt{/} (commande vérifiée sur le vrai paquet : 1~213~392~756 octets, 1~476 tuiles) :
\begin{ccode}
bsdtar -xf d80\_star\_database.deb data.tar.xz\\
sudo mkdir -p /opt/astap\\
sudo tar -xJf data.tar.xz -C /opt/astap --strip-components=3 --no-same-owner\\
\hspace*{2em}--wildcards \textquotesingle./opt/astap/d80\_*\textquotesingle
\end{ccode}
(jamais \texttt{sudo tar -x … -C /} : l'archive contient \texttt{./} et \texttt{./opt/}). \textbf{Arch/Manjaro} : le
paquet officiel met \texttt{astap} et \texttt{astap\_cli} dans \texttt{/opt/astap} seulement (rien dans le PATH :
Coupole les y trouve) ; l'\texttt{astap} graphique GTK2 demande \texttt{libgtk-x11-2.0}, absent des dépôts officiels
(AUR seulement) et non déclaré ; \texttt{astap\_cli} est autonome (lié statiquement) et suffit à Coupole ; pour
l'interface d'ASTAP, le paquet GTK3. Un test de l'intégration continue vérifie chaque lien, pour chaque système.

\begin{sloppypar}\noindent\textbf{Détection.} Ordre : chemin choisi (Préférences, \cmd{coupole astap --definir}), variable
\texttt{COUPOLE\_ASTAP}, \texttt{astap\_cli} partout (PATH puis emplacements par défaut), puis seulement \texttt{astap}, emplacements par défaut (Windows
\texttt{Program Files\textbackslash astap}, \texttt{\%LOCALAPPDATA\%\textbackslash astap} ; macOS
\texttt{/Applications/ASTAP.app}, \texttt{/usr/local/opt/astap} ; Linux \texttt{/opt/astap}, \texttt{/usr/bin},
\texttt{/usr/local/bin}, \texttt{\textasciitilde/astap}). Catalogue : dossier choisi, \texttt{COUPOLE\_ASTAP\_CATALOGUE},
dossier du programme, \texttt{/opt/astap}, \texttt{/usr/share/astap/data}, \texttt{/usr/local/opt/astap}. Un catalogue
est déclaré complet avec ses 1~476 tuiles ; la version est lue sur \texttt{astap\_cli -h}.\end{sloppypar}


\paragraph{Astrométrie en ligne (nova.astrometry.net)}\label{sec:nova}
\capture{img/reglages_astrometrie_fr.png}{0.6}
Quand une image n'a pas de solution astrométrique et qu'aucun outil n'est installé, Coupole peut la confier au service
\emph{libre et gratuit} nova.astrometry.net (Lang, Hogg, Mierle, Blanton \& Roweis 2010). Ordre suivi partout (Partages
du DU, Banque OHP pour les images sans solution quand ASTAP manque) : solution de l'en-tête, ASTAP installé,
astrometry.net installé (\texttt{solve-field}), puis nova. Il faut \emph{votre} clé d'API (compte gratuit sur
\url{https://nova.astrometry.net}, page de votre profil) dans \emph{Préférences > Astrométrie} ; elle est gardée à
part (\texttt{nova.json}, lisible par vous seul) et n'est envoyée qu'à nova. \textbf{Rien ne part sans votre accord},
demandé avant chaque lot. Ce que Coupole envoie, d'après la documentation officielle de l'API : une image allégée
(niveaux de gris 16 bits, réduite au-delà de 4 millions de pixels ; la solution est ensuite ramenée à l'image
d'origine, coefficients SIP compris), les indices connus (échelle en ″/px avec bornes $\pm 20$~\%, centre approché et
rayon), images \emph{non publiques} par défaut, sans licence commerciale ni modification. Une seule soumission à la
fois, sondage espacé (5~s puis jusqu'à 30~s), délai maximal, abandon propre ; une image déjà envoyée (même empreinte)
n'est \emph{jamais} renvoyée — son résultat, ou l'échec, est gardé en cache, et une soumission interrompue reprend le
sondage au lieu de renvoyer l'image.

% =====================================================================
\section{Formats de sortie}
\begin{center}\small
\begin{tabularx}{\linewidth}{@{}lXl@{}}
\toprule
\textbf{Format} & \textbf{Contenu} & \textbf{Lecteurs}\\
\midrule
XISF (défaut) & Float32 en ADU, \texttt{bounds="-1000:65535"} ; UInt16 pour les FITS IRIS entiers ; compression zstd+sh niveau 9 ; propriétés \texttt{Observation:*}, \texttt{Instrument:*} & PixInsight\\
XISF compatible & UInt16 + piédestal de 1~000~ADU (\texttt{PEDESTAL}), zlib+sh (ce qu'écrit N.I.N.A.) & PixInsight, Siril $\geq$ 1.4, N.I.N.A.\\
FITS \texttt{.fits.fz} & compression par tuiles sans perte : GZIP\_2 sans quantification (flottants), RICE\_1 (entiers) & Siril, astropy, DS9\\
FITS & float32 (ou UInt16) non compressé & tous\\
\bottomrule
\end{tabularx}
\end{center}
\textbf{XISF} : \texttt{bounds} fixe la plage que PixInsight ramène à [0,1] à l'ouverture : rien à régler, même
échelle pour toutes les poses ; la borne basse à −1~000 garde les valeurs négatives des poses courtes calibrées.
Taille : 36~\% du FITS d'origine pour le T120 (float64), 62~\% pour IRIS. \textbf{FITS compressé} : mesuré avec
astropy 8.0, RICE\_1 et HCOMPRESS\_1 quantifient les flottants (écart d'un ADU, même avec
\texttt{quantize\_level=0}) ; seuls GZIP\_1 et GZIP\_2 sans quantification restituent les flottants bit à bit.
\textbf{FITS float32} : dans PixInsight, régler la plage de lecture FITS sur 0–65~535 avec \emph{Truncate}.
Pour tous les formats sauf le XISF compatible, l'écart au FITS d'origine reste sous le demi-ULP du float32
(vérifié pixel par pixel).

\paragraph{XISF compatible (N.I.N.A., Siril).} UInt16 = ADU + 1~000 arrondis, écrêtés à [0, 65~535] ; le piédestal
garde les fonds négatifs des poses courtes calibrées et figure dans l'en-tête (\texttt{PEDESTAL}, convention de
PixInsight). Mesuré sur 66 poses de la banque (\texttt{outils/audit2/mesure\_perte\_uint16.py}) : arrondi
$\le 0{,}5$~ADU, moyenne du fond changée de 0,014~ADU au plus (bruit du fond : 10 à 109~ADU) ; pixels écrêtés en
haut : 0,012~\% au plus (étoiles au-delà de 64~535~ADU) ; plus aucun pixel perdu en bas. Sans piédestal, une pose
T120 de 10~s au fond sur-soustrait (−220~ADU) perdait 99,75~\% de ses pixels à 0 et 6,1~\% des pixels des poses de
1–2~s tombaient à 0 : d'où le piédestal. Les IRIS entiers restent sans perte.

\subsection*{Qui lit quoi (vérifié)}\label{sec:compat}
\begin{center}\small
\begin{tabularx}{\linewidth}{@{}>{\raggedright}p{0.25\linewidth}ccccX@{}}
\toprule
\textbf{Fichier} & \textbf{PixInsight} & \textbf{Siril} & \textbf{N.I.N.A.} & \textbf{ASTAP} & \textbf{Remarque}\\
\midrule
XISF Float32 (défaut) & oui & non & non & non & Siril $\geq$ 1.4 le lit mais hors de sa plage [0, 1] ; N.I.N.A. : image blanche\\
XISF compatible (UInt16, zlib) & oui & $\geq$ 1.4.0 & oui & non & ASTAP : XISF non compressé seulement\\
XISF UInt16 zstd (IRIS entiers) & oui & $\geq$ 1.4.0 & non (3.2) & non & zstd lu par N.I.N.A. depuis janvier 2026 (développement)\\
FITS Float32 & oui & oui & non & oui & N.I.N.A. suppose [0, 1] ; Aladin, astropy : oui\\
FITS UInt16 & oui & oui & oui & oui & \\
\texttt{.fits.fz} Float32 (GZIP) & — & oui & — & non & ASTAP : Rice seulement\\
\texttt{.fits.fz} UInt16 (Rice) & — & oui & — & $\geq$ 2026.07.16 & — : non vérifié\\
\bottomrule
\end{tabularx}
\end{center}
Sources : \textbf{PixInsight}, code de PCL (\texttt{XISF.cpp} : codecs zlib, lz4, lz4hc, zstd, avec ou sans
\texttt{+sh} ; \texttt{XISFReader.cpp} : \texttt{bounds} obligatoire pour un flottant, écrêtage puis ramené à
[0, 1]) ; spécification XISF 1.0 révision 1 (septembre 2026, \url{https://pixinsight.com/doc/docs/XISF-1.0-spec/XISF-1.0-spec.html}),
schéma XSD identique à celui qu'emploie Coupole. \textbf{Siril}, essai réel des AppImage officielles (siril-cli,
pixels relus et comparés) : 1.2.0 et 1.2.6 ne lisent pas le XISF (« file type not supported »), 1.4.0 et 1.4.4
lisent tous nos codecs bit à bit ; mais les flottants XISF y restent en ADU (statistiques multipliées par 65~535),
alors que les FITS flottants sont ramenés à [0, 1] ; paquets Debian 13 et Ubuntu 24.04 (Siril 1.2.x) : pas de XISF.
\textbf{N.I.N.A.}, code (étiquette \texttt{Version-3.2}, 27/11/2025 : \texttt{XISF.cs} l.~304–328 : lz4, lz4hc,
zlib, avec ou sans \texttt{+sh} ; \texttt{Float32Converter.cs} l.~23 : $(\text{ushort})(v \times 65535)$,
\texttt{bounds} jamais lu) et décodage réel d'un de nos fichiers. \textbf{ASTAP}, page officielle (XISF « uncompressed »,
\texttt{.fits.fz} Rice depuis 2026.07.16, GZIP non). Fichiers d'essai pour les testeurs : \texttt{tests/donnees/}
(image de synthèse en XISF par défaut, XISF compatible et FITS, moins de 300~Ko chacun).

\paragraph{Écart à la révision 1.} Coupole écrit dans \texttt{Observation:Center:RA/Dec} le centre de la solution
astrométrique (choix du traitement de référence, antérieur à la révision) ; la révision 1 réserve ces propriétés au
point visé par le télescope et normalise la solution dans l'espace \texttt{AstrometricSolution:} (non écrit : la
solution reste dans les mots-clés WCS, que PixInsight lit). Gardé tel quel en 0.1.9 pour ne pas changer les fichiers
par rapport au traitement de référence.

% =====================================================================
\section{Qualité des images (facultatif)}\label{sec:qualite}
\capture{img/module_qualite_fr.png}{1.0}
Le module demande la bibliothèque SEP (Source Extractor en Python, LGPL, environ 1 Mo) ; il n'est jamais lancé
automatiquement (bouton \emph{Mesurer}, case \emph{Vérifier la qualité des images} du traitement, ou
\cmd{coupole qualite DOSSIER --ecrire}). Rapport par lot : \texttt{QUALITE.csv} et \texttt{QUALITE.txt}.

\paragraph{Poses de ciel seulement.} Par défaut, seules les poses de ciel (lights) sont mesurées. Sont exclus sans
être ouverts : les dossiers Flat(s), Dark(s), Bias, Offset(s), DarkFlat(s), Calibration, Master(s) (« FLAT L »,
« Dark-Bias » compris), cosmetized, calibrated, registered, debayered, \texttt{\_platesolve…} et leurs noms
français (Plats, Noirs, Biais) ; les fichiers \texttt{FLAT\_}, \texttt{DARK\_}, \texttt{BIAS\_},
\texttt{DARKFLAT\_} (N.I.N.A., ASIAIR), les masters (\texttt{masterDark}, \texttt{masterFlat},
\texttt{masterLight}, \texttt{superbias}) et les sorties de WBPP (\texttt{\_c}, \texttt{\_cc}, \texttt{\_r}…).
En dernier recours, une pose dont l'en-tête (\texttt{IMAGETYP}, \texttt{FRAME} ou la propriété XISF
\texttt{Observation:Image:Type}) dit flat, dark, bias, offset ou master est exclue à la mesure ; Light, Light Frame,
Science ou Object sont mesurées. Le résumé dit « 70 image(s) trouvée(s), 66 pose(s) de ciel mesurée(s), 0 en erreur,
4 fichier(s) de calibration exclu(s) » ; \emph{Fichiers exclus…} en donne la liste avec le motif. Case
\emph{Inclure aussi les poses de calibration} (ligne de commande : \cmd{--avec-calibration}) : tout est mesuré.

\paragraph{Fichiers des autres logiciels, erreurs.} Les XISF de N.I.N.A. (mots-clés sans commentaire), de l'ASIAIR
(sans \texttt{Metadata}), de PixInsight et de Siril sont lus par un lecteur tolérant : seul ce qui empêche
vraiment de lire les pixels est refusé ; nos propres fichiers restent contrôlés strictement. Un fichier non mesurable
apparaît dans le tableau, colonne \emph{état}, avec son motif (« erreur : format XISF non lu (…) ») ; le résumé
compte les erreurs et \texttt{coupole.log} garde la pile ; la ligne de commande sort avec le code 1 si aucune image
n'a pu être mesurée. Pose N.I.N.A. d'une caméra CMOS (9~576 × 6~388, 61~Mpx) : 4,2~s, 1~170~Mo de pointe (calcul en
float32), FWHM 3,5″ à 1,695″/px ; le nombre de processus est réduit d'après la taille des images (moitié de la
mémoire libre) et le nombre d'étoiles ajustées croît avec le capteur (120 à 400).

\paragraph{Vitesse, reprise, échantillon.} Les mesures tournent dans des processus séparés (nombre donné par le plan
machine : cœurs, mémoire, mode économe et bridage des Préférences ; au plus 3 quand le dossier est sur un partage
réseau, détecté et signalé, car la lecture limite alors la vitesse), la lecture d'une image recouvrant le calcul des
autres ; l'interface reste libre, \emph{Arrêter} est immédiat. Chaque image mesurée est gardée aussitôt dans le
cache \texttt{\_traitement/qualite.sqlite} par (chemin, taille, date) et affichée ; le \texttt{QUALITE.csv} de son
lot (écriture atomique) est réécrit au plus toutes les 5~s et à la fin du lot ou de l'arrêt : relancer ne refait pas
ce qui est fait, fermer puis rouvrir reprend.

\paragraph{Inventaire en flux.} La mesure commence dès le premier lot trouvé, pendant que l'inventaire continue
(« Inventaire du dossier : 1 234 fichiers trouvés dans 56 lots… », \emph{Arrêter} immédiat) ; l'inventaire n'ouvre
aucun fichier et lit taille et date dans la liste de chaque dossier (un \texttt{stat} au plus par image). Dans une
sortie de Coupole, seuls les dossiers des lots sont lus, d'après \texttt{INDEX\_LOTS.csv} ou la base d'état ; les
restes de traitement (\texttt{converties}, \texttt{telechargements}) ne sont jamais pris pour des lots. Sur un
partage, 32 fils lisent les dossiers. Le cache n'est \textbf{jamais} écrit sur un partage réseau : le client SMB de
Linux transmet les verrous de SQLite au serveur, qui les refuse (« database is locked » au bout de 30~s, trois fois
avant la première mesure en 0.1.8 : les 90~s d'« Inventaire du dossier… » observées sur un vrai partage) ; il va dans
le dossier de cache de l'utilisateur (un ancien cache du partage est relu). Mesuré sur un vrai partage Samba, 831
XISF de 09\_Galaxies : 0.1.8 92,7~s avant la première mesure (latence de 10~ms) ; 0.1.9 inventaire complet en
0,24–0,57~s (10~ms) et 0,57–1,05~s (20~ms), premier résultat en 0,2–0,6~s. Au-delà de 200 images, Coupole propose un \textbf{échantillon} de 5 images par lot (réparties dans
le temps : première, dernière, milieu…, réglable dans l'onglet et par \cmd{--echantillon N}), qui suffit à juger un
lot (mesuré en premier, pendant que vous choisissez), ou \emph{Tout mesurer} avec la durée estimée d'après les
images déjà mesurées ; la progression donne le temps
restant et le débit (images/min). Sur un processeur à 6 cœurs : 1,21 → 0,33~s par image T120
(modèle de Moffat vectorisé, dérivées analytiques, 120 étoiles par image), soit la banque entière en une dizaine de
minutes sur 5 processus.

\begin{center}\small
\begin{tabularx}{\linewidth}{@{}lXl@{}}
\toprule
\textbf{Mesure} & \textbf{Méthode} & \textbf{Validation synthétique}\\
\midrule
FWHM & Moffat elliptique à $\beta$ libre ajusté sur chaque étoile isolée et non saturée (modèle intégré sur le pixel) ; médiane & ≤ 0,4~\% (Gauss et Moffat, 2,5–8 px)\\
ellipticité & $1-b/a$ du même ajustement & ± 0,02\\
carte 3 × 3 & médianes par neuvième de champ & ordre retrouvé\\
fond, bruit & SEP, maille de 64 px & 0,5~\% ; 5~\%\\
gradient, résidu & plan ajusté sur la carte du fond ; amplitude du résidu & gradient et vignetage retrouvés\\
saturées & pic au-delà de SATURATE/DATAMAX (sinon 60~000 ADU) & 5 sur 5\\
traînées & objets de grand axe ≥ 12 px et $a/b \ge 6$ & 1 sur 1, 0 faux\\
RSN & seulement avec le gain (EGAIN/GAIN) & donné ou absent\\
échantillonnage & FWHM < 2 px, 2–5 px, > 5 px & —\\
\bottomrule
\end{tabularx}
\end{center}
Sur 14 images IRIS (12 nuits, FWHM de 2,4 à 11,3 px), la FWHM mesurée et celle écrite par PinPoint dans l'en-tête
sont corrélées (r = 0,974), rapport médian 1,046 (0,85 à 1,16) : les deux estimateurs diffèrent (méthode de PinPoint
non décrite). Non proposés faute de validation : mag\,arcsec$^{-2}$ (pas de point zéro), détection des « donuts » de
poussière. Une FWHM qui croît vers un bord peut venir du capteur, de la courbure de champ, de la collimation, mais
aussi de la turbulence, du suivi ou de la mise au point : la carte seule ne tranche pas.

% =====================================================================
\section{Spectres et séries}
\capture{img/module_donnees_fr.png}{1.0}
Formats lus : FITS (image 1D avec axe WCS linéaire \texttt{CRVAL/CDELT/CRPIX}, cubes radio 1 × 1 × N, tables
binaires) et CSV ou texte en colonnes (en-tête facultatif, unité entre parenthèses ou crochets). Vitesse radio :
$v = c\,(1 - f/f_0)$ (convention radio de la norme FITS), $f_0$ lue dans \texttt{RESTFRQ}, sinon raie H\,I à
1~420,405~751~768 MHz. Le référentiel est celui de l'axe des fréquences (\texttt{SPECSYS}) ; aucune correction LSR
ou barycentrique n'est faite implicitement. Les formats des radio-antennes d'enseignement s'ajoutent par un lecteur
(CONTRIBUTING.md, « Ajouter un format »).

% =====================================================================
\section{Sites et heures}
\capture{img/module_sites_fr.png}{0.95}
Le temps est toujours tenu en UTC (dates « aware ») ; l'heure locale du site et celle de l'utilisateur sont
affichées en plus. La « date du soir » d'une nuit est la date locale du site à midi précédent ; le plein jour se
calcule par la hauteur du Soleil au site (astropy, sans téléchargement IERS). Sites livrés : OHP (MPC 511), Meudon
(005), Paris (007), coordonnées géodésiques calculées depuis les constantes de parallaxe du Minor Planet Center.
Les sites partenaires s'ajoutent (fuseau IANA, code MPC s'il existe) quand leurs coordonnées sont connues de
source sûre. Carte : tuiles OpenStreetMap en cache disque, deux téléchargements au plus, User-Agent identifiant
l'application, attribution affichée ; hors ligne, fond simple et points. Sur un écran dense (mise à l'échelle
$\geq 150$~\%, Mac Retina), la carte dessine les tuiles du niveau de zoom supérieur à demi-taille (quatre par tuile
logique, l'équivalent « @2x ») : elle reste nette au lieu d'être agrandie par le système.

% =====================================================================
\section{Cosmologie}\label{sec:cosmo}
\capture{img/module_cosmo_fr.png}{1.0}
Le redshift $z$ d'un objet lointain ne donne pas « une » distance mais plusieurs, toutes justes, qui répondent à des
questions différentes. Le module les calcule dans le modèle $\Lambda$CDM avec
\[ E(z)^2 = \Omega_r(z)(1+z)^4 + \Omega_m(1+z)^3 + \Omega_k(1+z)^2 + \Omega_\Lambda, \qquad
   D_C = \frac{c}{H_0}\int_0^z\frac{dz'}{E(z')}, \]
rayonnement compris (photons du fond diffus et neutrinos, dont un massif de 0,06~eV).
\begin{center}
\begin{tabularx}{\linewidth}{@{}lX@{}}
\toprule
Grandeur & Usage \\
\midrule
$D_C$ comobile radiale & distance actuelle de l'objet ; sature à $\approx 46$~G\,al (horizon des particules)\\
$D_M$ comobile transverse & égale à $D_C$ si l'univers est plat ; entre dans $D_L$ et $D_A$\\
$D_L = (1+z)D_M$ & photométrie : $F = L/(4\pi D_L^2)$ ; module de distance $\mu = 5\log_{10}(D_L/10~\mathrm{pc})$\\
$D_A = D_M/(1+z)$ & tailles : $\theta = \ell/D_A$ ; échelle en kpc par seconde d'arc ; maximum vers $z\approx1{,}6$\\
$c\,t_L$ & chemin parcouru par la lumière (le plus intuitif, le moins physique)\\
$t_L$, âge à $z$, $t_0$ & temps de regard en arrière ; contrôle $t_L + t(z) = t_0$\\
$V_C$ & volume comobile jusqu'à $z$ (Hogg 1999, exact aussi en univers courbe)\\
vitesses & $cz$ (naïf), Doppler relativiste (inadapté), $H_0 D_C$ (la seule correcte)\\
\bottomrule
\end{tabularx}
\end{center}

\paragraph{Paramètres.} Planck 2018 par défaut ($H_0 = 67{,}66$, $\Omega_m = 0{,}3111$ dont $\Omega_\nu = 0{,}00144$,
$T_0 = 2{,}7255$~K) avec l'incertitude $1\sigma$ propagée depuis $H_0$ et $\Omega_m$ \emph{corrélés}
($\rho = -0{,}976$ : l'oublier surestimerait l'incertitude) et, au choix, la comparaison avec SH0ES
($H_0 = 73{,}04$, tension de Hubble). Pour comparer : Planck 2015, WMAP 9 ans, un $\Lambda$CDM « de manuel »
(70~/ 0,3~/ 0,7 sans rayonnement) et un jeu personnalisé ($H_0$, $\Omega_m$, $\Omega_k$, rayonnement de Planck 2018).
Plages : $H_0$ de 40 à 100, $\Omega_m$ de 0,1 à 1, $\Omega_k$ de $-0{,}3$ à $0{,}3$ ($\pm 0{,}05$ pour Planck 2018) ;
une valeur hors plage est \emph{refusée avec une explication}, jamais corrigée en silence.

\paragraph{Saisie.} Un $z$ (virgule ou point), un des exemples (M\,87, 3C\,273, $z = 1$, $z = 2{,}34$, ULAS J1120,
GN-z11, réionisation, fond diffus), ou un nom d'objet : son redshift est demandé à SIMBAD (nom tolérant :
« m87 », « ngc4486 » ; plusieurs candidats : vous choisissez). $z \le 0$ (objet qui s'approche, comme M\,31) ne donne
pas de distance ; sous $z = 0{,}03$ la vitesse propre de l'objet domine (avertissement) ; au-delà de $z \approx 1090$
l'univers est opaque. Tableau et courbes s'exportent en CSV (« ; », UTF-8 ; valeurs brutes en Mpc, Gyr, Gpc$^3$).

\paragraph{Curseur.} Sous le champ $z$, un curseur parcourt tout le domaine sur une échelle logarithmique
(de 0,001 à 1100, repères 0,01 / 0,1 / 1 / 10 / 100 / 1000). En glissant, le marqueur vertical des courbes et le
tableau des valeurs suivent en direct : les valeurs sont alors \emph{interpolées} sur la grille des courbes (300
points logarithmiques, écart typique $< 0{,}5$~\%, colonnes $\sigma$ et SH0ES suspendues) et la ligne d'état le dit ;
au relâchement, le calcul exact reprend. Flèches : pas fin (1~\% d'une décade) ; Page haut/bas : une décade. Le
champ $z$ et le curseur restent synchronisés dans les deux sens.

\capture{img/module_cosmo_cote_fr.png}{1.0}
\paragraph{Disposition.} Côte à côte, le tableau prend la largeur de ses colonnes (au contenu, jamais tronquées) et
les courbes tout le reste ; empilés, le tableau montre ses 18 lignes si la place le permet, sinon la place est
partagée en gardant au moins 300~px aux courbes. \emph{Affichage > Disposition de la Cosmologie} (ou la liste
\emph{Disposition} du module) : automatique (côte à côte à partir de 1~550~px de large, retour à l'empilement sous
1~450~px : pas de bascule pendant un redimensionnement), toujours côte à côte, ou toujours empilés ; le choix et la
position du séparateur de chaque disposition sont gardés.

\paragraph{Unités.} Notation à exposants négatifs : $H_0$ en km\,s$^{-1}$\,Mpc$^{-1}$, vitesses en km\,s$^{-1}$ ;
dans les CSV, en ASCII non ambigu (\texttt{km s\^{}-1 Mpc\^{}-1}), de même sur une console qui ne sait pas afficher
les exposants (Windows cp1252).

\paragraph{Origine et vérification.} Le noyau de calcul est repris du calculateur « cosmologie-redshift » du même
auteur (dépôt public \url{https://github.com/ARP273-ROSE/cosmologie-redshift}), avec astropy.cosmology. Avant
l'intégration, il a été comparé le 8 octobre 2026 à astropy et à une intégration indépendante sous SageMath
(mpmath, 18 chiffres, fonction de Fermi-Dirac exacte pour les neutrinos), de $z = 10^{-8}$ à 1100, pour
$\Omega_k = 0, \pm0{,}01, \pm0{,}05$ : écart relatif $\le 3\times10^{-6}$ sur les distances et
$\le 2\times10^{-5}$ sur les âges et $E(z)$ (ajustement de Komatsu des neutrinos dans astropy) ; aucun bug. Le
volume comobile d'astropy perd en revanche tous ses chiffres à très petit $z$ en univers courbe (jusqu'à un volume
négatif) : Coupole emploie là un développement limité, raccordé à la forme fermée. Scripts :
\texttt{outils/verif\_cosmo/}.

\begin{ccode}
coupole cosmo 0,158 --shoes\\
coupole cosmo --objet "3C 273"\\
coupole cosmo 1 2 3 --modele wmap9 --csv resultats.csv\\
coupole cosmo --courbes 0,001 1500 300 --csv courbes.csv
\end{ccode}

% =====================================================================
\section{Animations}\label{sec:anim}
\capture{img/module_anim_fr.png}{1.0}
Le module refait en Python les animations pédagogiques \emph{ClassAction} de l'université du Nebraska-Lincoln
(\emph{Astronomy Education}), écrites à l'origine en Flash et devenues illisibles avec la fin de ce lecteur. Elles
tournent désormais sous Windows, macOS et Linux, en français et en anglais, dans le thème de Coupole, et leurs calculs
reposent sur des formules de référence vérifiées. D'après ClassAction, University of Nebraska-Lincoln Astronomy
Education, licence Apache 2.0 (fichiers \texttt{NOTICE} et \texttt{LICENCES/Apache-2.0.txt}).

\paragraph{Catalogue.} À gauche, les 92 animations du catalogue d'origine, plus trois (la machine d'Anticythère, section~\ref{sec:anticythere}, Orion en 3D, section~\ref{sec:orion}, et Gravité et orbites, section~\ref{sec:gravite}), classées par thème (coordonnées et
mouvements, Lune, Renaissance et orbites, propriétés stellaires, lumière et spectres, exoplanètes…), avec une recherche
(française ou anglaise, sans accents), un filtre par thème et par état. Les thèmes suivent un ordre pédagogique, de
l'introduction et du ciel vu de la Terre jusqu'à la cosmologie. Chacune a son nom et sa description dans les
deux langues, et son titre d'origine. Depuis la version 0.4.0, les 92 sont disponibles : celles qui n'étaient que
des films (ombre d'un bâtiment, clou qui fond, lanceur de marteau, ballon, réactions de fusion, embouteillage…) sont
devenues des simulations, et celles qui reposent sur des données réelles (spectres, courbes de lumière, exoplanètes,
planètes) en citent la source dans le fichier \texttt{NOTICE}, mise à jour quand elle était périmée.

\paragraph{Manipuler.} Glisser sur une sphère la fait tourner ; glisser une étoile, la Lune, la Terre ou un
observateur les déplace. \emph{Lecture / Pause} et \emph{Vitesse} (0,1 à 10~fois, échelle logarithmique) animent le
temps ; \emph{Réinitialiser} revient à l'état de départ. Chaque réglage a un curseur et un champ numérique (la valeur
exacte est gardée : le champ n'en montre qu'un arrondi) ; les valeurs calculées s'affichent sous les réglages. Tous
les réglages, le thème et l'animation choisie sont retrouvés au lancement suivant. Une seule animation tourne à la
fois, et seulement quand elle est affichée ; une image coûte quelques millisecondes (rien de lourd dans le fil
graphique). Ce que fait la souris dans l'animation affichée est rappelé en petit sous la scène (plus d'info-bulle
par-dessus le dessin). Le bouton \emph{Version web} ouvre l'animation, avec ses réglages, dans Coupole Web
(\texttt{arp273-rose.github.io/coupole-web}) : les mêmes animations dans un navigateur, sur ordinateur, tablette ou
téléphone.

\subsection{Rendu, formules et cartes (0.4.1)}
\paragraph{Astres.} Le Soleil, la Lune, la Terre et Mars sont dessinés en sphères : dégradé radial dont le foyer est
décalé vers la lumière (limbe assombri) et, quand un Soleil est dessiné dans le même panneau, côté nuit ombré au-delà du
terminateur ; le Soleil a un limbe assombri et un halo discret. Tout reste vectoriel et lissé ; les brosses sont mises
en cache, si bien qu'une image coûte toujours quelques millisecondes (médiane 1,5~ms par image à 720 × 470, comme avant ces changements, mesurée sur les
92 animations). Les symboles (icônes de la table des éclipses, repères) restent plats.

\paragraph{Formules.} Les formules sont composées par un petit moteur de mise en page écrit pour Coupole (QPainter,
sans \LaTeX{} ni matplotlib) : fractions empilées, exposants et indices bien placés, radicaux, parenthèses extensibles,
$\times 10^n$ en vrai exposant, variables d'une lettre en italique mathématique, mots (sin, log, TSL) droits. La police
est STIX Two Math (licence SIL OFL 1.1, sous-ensemble embarqué de 139~ko, \texttt{LICENCES/OFL-STIX-Two.txt}). Les
formules dessinées dans les animations (calculateur de la loi de Newton, algèbre de la gravitation, module de distance,
petits angles, luminosité, flux, excentricité…) et celles des descriptions, des libellés et des valeurs affichées sont
toutes mises en forme ; les mises en page sont gardées en cache.

\paragraph{Étiquettes.} Toutes les étiquettes à cartouche passent par un placeur commun : si leur place est prise
(autre étiquette, petit astre, texte déjà écrit, bord de l'image), elles essaient 24 positions autour de leur objet et
prennent la première libre, reliée à l'objet par un trait de rappel. Un test vérifie, pour les 92 animations, quatre
tailles (420 à 2560 pixels de large) et trois instants, qu'aucune étiquette n'en recouvre une autre ni ne sort de
l'image.

\paragraph{Cartes.} Les côtes viennent de Natural Earth 1:50m (domaine public), simplifiées (Douglas-Peucker 0,07°,
65~ko), chargées au premier globe affiché. Le simulateur jour-nuit propose quatre projections. Équirectangulaire :
correspondance linéaire entre la souris et la latitude-longitude. Pôle Nord ou pôle Sud : azimutale équidistante
centrée sur le pôle (distance au pôle vraie), bord à l'équateur ou 30° au-delà pour voir les deux tropiques ;
méridiens tous les 15° (30° marqués), parallèles tous les 15°, tropiques et cercles polaires en pointillé ; par défaut
la Terre tourne sous un terminateur fixe (midi en bas), et \emph{Terminateur tournant} fixe la carte (Greenwich en
bas au Nord). La nuit et les crépuscules y sont tracés méridien par méridien, à la latitude exacte où le Soleil est à
0°, −6°, −12° ou −18° : une boucle autour du pôle, sans artefact au pôle (au solstice, elle est tangente au cercle
polaire). Globe 3D vu de l'espace (projection orthographique) : il se tourne dans tous les sens à la souris (rotation
de boule de commande : le point saisi suit la souris, avec une légère inertie), la molette zoome, le double-clic
revient au point subsolaire, le clic droit place l'observateur ; atmosphère fine, sphère éclairée du côté du Soleil,
lignes coupées exactement au limbe, nuit et crépuscules exacts (parties visibles des calottes de la sphère opposées
au Soleil), axe des pôles et flèche « vers le Soleil ». Au pôle Nord et au pôle Sud, le jour ou la nuit polaire
s'affiche avec les valeurs.

\paragraph{Voie lactée.} La courbe de rotation n'est plus faite de deux morceaux raccordés à 5~kpc : c'est celle d'un
modèle de masse publié, McMillan (2017), calculée avec galpy (bulbe, disques d'étoiles mince et épais, gaz HI et H2,
halo NFW de matière noire ; composantes affichables, $v^2$ étant la somme des $v^2$ des composantes), avec
$R_0 = 8{,}21$~kpc et $v_0 = 233{,}1$~km/s. Les mesures réelles s'y superposent avec leurs barres d'erreur : Eilers et
al. (2019) de 5 à 25~kpc et la compilation galkin (Pato \& Iocco 2017) de 2 à 5~kpc. La courbe képlérienne part de la
masse que le modèle contient dans la sphère de rayon $R$. La Galaxie est vue du pôle Nord galactique, à l'échelle :
barre longue (Wegg et al. 2015), bras spiraux log-périodiques (Reid et al. 2019), Soleil à 8,21~kpc avec son sens de
rotation, cercle du rayon choisi. SageMath retrouve la vitesse du halo et la masse du bulbe du modèle à
$1{,}3\times10^{-4}$ et $5\times10^{-5}$ près.

\paragraph{Éclipses.} Dans le simulateur des ombres d'éclipse, le Soleil est dessiné en entier, et \emph{Rejouer une
éclipse} montre chacune des 133 éclipses de la table 2015–2044 en proportions réelles : la Lune traverse l'ombre et la
pénombre de la Terre, ou l'ombre de la Lune balaie la Terre ; l'écart à l'axe vient de l'argument de latitude moyen de
la Lune ($y = D \sin 5{,}145^\circ \sin F$, Meeus), ramené dans le domaine compatible avec le type de l'éclipse, et un
curseur fait avancer l'instant.

\paragraph{Corrections.} Le globe du démonstrateur longitude/latitude (et toutes les sphères) tournait à l'envers du
glisser : le point saisi suit maintenant la souris. Les rayons du Soleil du simulateur des saisons s'arrêtent au limbe
et au bord du panneau ; les noms des étoiles de la Grande Ourse 3D et du lanceur de marteau ne se chevauchent plus ;
quelques textes trop longs pour les petites fenêtres ont été déplacés (légendes, horloges, niveaux de l'hydrogène,
ombre du bâtiment, rayons du Soleil, spectre électromagnétique).

\capture{img/anim_daylightsimulator_nord_fr.png}{0.95}
\capture{img/anim_daylightsimulator_globe_fr.png}{0.95}
\capture{img/anim_shadowsim_rejeu_fr.png}{0.95}
\capture{img/anim_gravcalc_fr.png}{0.95}
\capture{img/anim_hammerthrower_fr.png}{0.95}
\capture{img/anim_bigdipper_fr.png}{0.95}

\subsection{La machine d'Anticythère (0.4.2)}\label{sec:anticythere}
\capture{img/anim_antikythera_avant_fr.png}{0.95}
Une 93\textsuperscript{e} animation, hors ClassAction, dans « Astronomie antique » : le calculateur astronomique grec
du \textsc{ii}\textsuperscript{e} siècle av.~J.-C., en 3D, adapté d'\emph{Anticythere3D}, du même auteur (application
complète : \texttt{github.com/ARP273-ROSE/Anticythere3D}). Les nombres de dents (Freeth et al., \emph{Nature}~444,
2006 ; \emph{Nature}~454, 2008), les trains, les rapports exacts (fractions), l'implantation des 15 arbres
(optimisation sous contraintes, 17 entraxes exacts), le tenon et la fente, les cadrans et leurs graduations sont ceux
d'Anticythere3D 1.3.1, repris à l'identique, avec ses tests.

\paragraph{Manipuler.} On saisit la poignée de la manivelle (flanc droit de la face avant) et on la tourne : un tour de
poignée fait $48/223$ de tour de la grande roue b1, soit 78,6~jours. Ailleurs, glisser fait tourner la machine dans tous
les sens (orientation par quaternion, jamais invalide) ; \emph{Lecture} tourne la manivelle d'un tour par seconde
(fois la vitesse) ; \emph{Aujourd'hui}, \emph{+1~lunaison}, \emph{+1~an}. \emph{Enveloppe} : machine fermée (coffret,
plaques de bronze gravées en grec) ou ouverte (les 33 roues colorées par fonction, légende) ; \emph{Éclatement des
étages} écarte les 16 plans d'engrènement le long des arbres ; \emph{Vue} : face avant, face arrière ou trois-quarts.
Survoler une roue (ou la toucher) donne son nom, ses dents, son rapport exact en tours par an et sa fonction.

\paragraph{Ce que montrent les cadrans.} Face avant : zodiaque (360 divisions, noms grecs, lettres du parapegma),
anneau calendaire égyptien de 365~jours, tourné pour que l'aiguille du Soleil indique la date (ère de Nabonassar,
1\textsuperscript{er} Thoth de l'an~1 = JD~1~448~638, vérifié avec SageMath), aiguilles du Soleil et de la Lune, sphère
de phase demi-argentée, anneaux du Cosmos (Freeth et al., \emph{Scientific Reports}~11:5821, 2021) dont la petite sphère
indique la longitude géocentrique de la planète (relation de période et épicycle). Face arrière : spirale métonique
(5~tours, 235~cases, mois corinthiens, années en chiffres grecs), spirale du Saros (4~tours, 223~cases, glyphes
$\mathrm{H}$ et $\Sigma$ calculés), cadrans callippique, de l'exeligmos (0, 8, 16~h) et des Jeux, chacun sur l'arbre qui porte son
aiguille. Les valeurs : date (grégorienne ou julienne), date égyptienne, Soleil, Lune, phase, correction d'anomalie,
nœud, cases des spirales, éclipse possible, planètes, écart au ciel réel (Meeus).

\paragraph{Rendu sans OpenGL.} Projection orthographique et algorithme du peintre : toutes les pièces sont planes et
parallèles, chacune est dessinée dans son plan par une transformation affine exacte, les plans sont triés le long de
la visée, les arbres coupés en tronçons aux plans des roues qui peuvent les recouvrir ; le coffret fermé est convexe,
seules ses faces visibles sont dessinées. Les cadrans gravés sont dessinés en vectoriel (mêmes tracés
qu'Anticythere3D), mis en cache dans une image dont la résolution suit le zoom ; une denture de moins de 3~pixels par
dent est remplacée par son cercle primitif. Une image coûte 1 à 4~ms fermée, 5 à 9~ms ouverte (QPainter, 420 à
720~pixels, NAS chargé) ; dans Coupole Web, 0,4 à 1,2~ms (Chromium, téléphone émulé) et 1,5 à 5~ms avec un processeur
ralenti 4~fois.

\paragraph{Vérifications.} SageMath (\texttt{outils/verif\_anim/anticythere\_verif.sage}, nombres de dents ressaisis,
arithmétique rationnelle) : Lune sidérale $254/19$, d'où $254-19=235$ lunaisons en 19~ans ($235/19$, réduite de la
fraction continue de l'année en lunaisons), métonique $5/19$, callippique $1/76$, Saros $940/4237=(4/223)(235/19)$,
exeligmos $235/12711$ (54,09~ans), porte-satellite $477/4237$ (8,8826~ans), mois anomalistique de la machine
27,5533~j, amplitude du tenon-fente $\arcsin(1{,}1/9{,}6)=6{,}58^\circ$, 17 entraxes. Les tests d'Anticythere3D
(rapports, porte-satellite, tenon-fente, retour des cadrans, planètes, calage, entraxes) sont repris, avec le rendu
(vues, glissers aberrants pendant la lecture, roue survolée, manivelle) ; les valeurs affichées sont identiques au
caractère près à celles de Coupole Web.

\capture{img/anim_antikythera_arriere_fr.png}{0.95}
\capture{img/anim_antikythera_ouverte_fr.png}{0.95}

\paragraph{Autres corrections (0.4.2).} Les titres des panneaux ne débordent plus sur le panneau voisin dans une
fenêtre étroite (Phases de Vénus, chez Copernic comme chez Ptolémée, et 20 autres animations) : coupés sur deux lignes,
raccourcis ou terminés par « … », le titre complet restant en info-bulle et dans le nom accessible de la scène. Trois
vues d'un spectre : l'ampoule, opaque, cache le nuage placé derrière elle (spectre continu) ; l'absorption n'apparaît
que si le gaz est entre le télescope et l'ampoule. Le simulateur jour-nuit s'ouvre désormais sur le globe 3D.

\subsection{Spectres physiques, Orion en 3D, fiches au survol (0.5.0)}\label{sec:orion}
\paragraph{Spectres sur bases physiques.} L'Explorateur de spectres et la Parallaxe spectroscopique ne reprennent plus
les segments de droite et les 37 raies de l'original : les intensités des six familles de raies sont calculées par
Saha–Boltzmann sur des atmosphères modèles (grille ATLAS9 de Castelli \& Kurucz), à la profondeur optique $2/3$, pour
chaque type spectral et chaque classe de luminosité (V, III, I) ; 133 raies du NIST (longueurs d'onde dans l'air, forces
d'oscillateur), bandes de TiO et de la bande G du CH dégradées vers le rouge, continu de Planck coloré par les fonctions
CIE 1931. Courbes lisses, avec le type spectral en abscisse et la température de la classe en haut ; noms des raies les
plus fortes au-dessus du spectre. G2~V = 5\,770~K (Mamajek 2022). Les cas extrapolés (atmosphère hors de la grille,
étalonnage borné ou interpolé) sont signalés sous le spectre et dans les valeurs ; la Parallaxe spectroscopique prend
la magnitude absolue publiée (classe V partout, III pour les étoiles O et de G à M7, I pour les O) et, ailleurs,
l'ancien étalonnage, signalé lui aussi (courbe en pointillé). Le moteur et ses données (copiés tels quels, mêmes tests
de référence que Coupole Web, vérifiés par SageMath) sont décrits dans \texttt{NOTICE}.
\capture{img/anim_spectrum010_m5_fr.png}{0.95}

\paragraph{Orion en 3D.} Une 94\textsuperscript{e} animation, propre à Coupole, dans un nouveau thème « Milieu
interstellaire et formation des étoiles ». Orion telle qu'on la voit (vue du Soleil : la perspective est la projection
gnomonique du ciel, nord en haut, est à gauche), puis \emph{Dépliement} : chaque étoile glisse sur sa ligne de visée
jusqu'à sa distance mesurée (SIMBAD, Gaia~DR3, Hipparcos~2007, mesures dédiées : Bételgeuse, $\zeta$~Ori,
$\sigma$~Ori, Trapèze en VLBA) et la vue tourne de côté : Tabit est à 8~pc, Bellatrix à 77, Bételgeuse à 168, Rigel
à 265, la Ceinture vers 380–400~pc. Barres d'erreur le long de la visée ; Alnilam, Mintaka, Meissa et Hatysa, dont
la mesure propre est peu fiable, sont placées à la distance publiée de leur groupe (option \emph{Distances mesurées
brutes}). Nuages moléculaires Orion~A et B (squelettes 3D de Zucker et al.~2021), anneau de $\lambda$~Orionis, boucle
de Barnard sur la coquille d'Orion (Foley et al.~2023), nébuleuse d'Orion en « blister » (front d'ionisation en bol,
Voile, Trapèze ; Wen \& O'Dell~1995, O'Dell et al.~2009) et Tête de Cheval. \emph{Regarder} : plongée vers M42 ou
vers la Tête de Cheval, avec les images de l'ESO (VST, VLT) ou de Hubble posées sur le ciel à leur place exacte ;
leur crédit, mot pour mot, est affiché sur la scène et en lien sous la scène. Couleur vraie des étoiles (corps noir à
leur température, CIE 1931), option \emph{Couleurs accentuées}. Tout est à l'échelle, sauf la taille des étoiles.
Glisser fait tourner la scène autour du pivot ; double-clic : retour ; \emph{Lecture} : la scène tourne lentement.
Une image coûte 5 à 8~ms (QPainter, 720~pixels, NAS chargé), 0,4 à 1,1~ms dans Coupole Web.

\paragraph{Fiches au survol.} Dans Orion en 3D, la Grande Ourse 3D, l'horloge de la Grande Ourse, le lever
héliaque, la comparaison des repères et le diagramme HR, survoler une étoile ou une nébuleuse (ou la toucher, ou y
venir au clavier : Tab de la plus brillante à la plus faible, flèches, Échap) ouvre une petite fiche, posée à côté de
l'objet, jamais dessus : noms et désignations, type, magnitudes, distance en parsecs et en années-lumière avec son
incertitude (« mesure incertaine » quand il le faut, autres mesures), température, rayon et luminosité publiés, âge et
étoiles excitatrices des nébuleuses, et la source de chaque valeur. Le résumé est lu par les lecteurs d'écran. Sur une
petite toile, la fiche est compacte. Grande Ourse 3D : distances mises à jour (parallaxes adoptées par SIMBAD) ;
diagramme HR : case \emph{Couleur = température}.
\capture{img/anim_orion3d_fiche_fr.png}{0.95}

\paragraph{Typographie.} Unités à la manière SI (km$\cdot$s$^{-1}$, W$\cdot$m$^{-2}$…), vrai signe moins, espaces
fines, virgule décimale, puissances de dix en exposant, formules aux parenthèses et radicaux extensibles, indices de
plusieurs lettres en romain ($\lambda_{\mathrm{max}}$) : dans toutes les animations, avec un formateur commun au
bureau, au web et aux ressources, et un test qui passe en revue tous les textes affichés.

\paragraph{Vérifications.} SageMath : conversion parsec $\to$ année-lumière depuis les définitions de l'UAI, distances
et barres des 300 fiches, parallaxe du Trapèze (Gaia~DR3, 4 étoiles : 402 $\pm$ 4~pc, montrée dans la fiche à côté des
388 $\pm$ 5~pc du VLBA), matrice ICRS $\to$ galactique contre sa définition ; couleurs (21 étoiles) et modèle des spectres
(46 comparaisons) vérifiés par les jeux de données. Les nombres d'Orion en 3D, le contenu des fiches et leur placement
sont identiques à ceux de Coupole Web (références communes).

\subsection{Modules, questions et ressources (0.6.0)}\label{sec:classe}
\paragraph{Navigation par modules.} Le module \emph{Animations} s'ouvre sur l'accueil des \textbf{22 modules} de
ClassAction (et deux modules bonus) : boutons \emph{Modules}, \emph{Bibliothèque}, \emph{Animations} (le catalogue), \emph{Mes favoris},
\emph{Ma progression}, une recherche et un fil d'Ariane (avec retour à la page précédente). Une page de module a des
onglets \emph{Questions} (Échauffement, Générales, Défis, Discussion, numérotées W1, G1, C1, D1, avec leur état
$\circ$, $\bullet$, $\checkmark$, $\times$), \emph{Animations}, \emph{Fiches}, \emph{Tables}, \emph{Images}. Double-clic ou Entrée ouvre un élément ; une
fiche, une table ou une image s'ouvre dans sa propre fenêtre (plusieurs à la fois, redimensionnables, Échap les
ferme) ; survoler un élément montre un aperçu ; clic droit : \emph{Ajouter à\dots{}} un module personnel.
\capture{img/classe_accueil_fr.png}{0.95}
\capture{img/classe_module_fr.png}{0.95}

\paragraph{Questions.} Les 517 questions d'instruction par les pairs de ClassAction (Apache~2.0) sont traduites et
réécrites en données (les mêmes que dans Coupole Web, même moteur : les empreintes de plus de 8\,000 instances sont
identiques). Menus \emph{Adapter la question} (variantes), \emph{Nouveau tirage}, figures vectorielles, choix A, B, C\dots{}
(touches A à H), \emph{Valider} (Entrée) enregistre la tentative, \emph{Montrer la réponse} (R) fait clignoter la bonne
réponse en orange puis la laisse orange, \emph{Indice} (I), explication, éléments de réponse des questions de
discussion, volet « Ce qui a changé par rapport à l'original », ressources liées $\star$ à droite, question précédente et
suivante (flèches, ordre W $\to$ G $\to$ C $\to$ D en boucle), plein écran (F).
\capture{img/classe_question_fr.png}{0.95}

\paragraph{Fiches, tables, images.} 110 fiches, 14 tables et 109 schémas écrits pour Coupole (valeurs recalculées et
sourcées) : texte et formules par le moteur de formules de Coupole, schémas dessinés par un moteur QPainter qui lit leur
description (largeur réelle des textes par \texttt{QFontMetricsF} pour les marges, contrôle de non-chevauchement),
tables à colonnes numériques alignées à droite (source en info-bulle), photographies libres avec leur crédit et leur
licence (liste complète dans \texttt{NOTICE}).
\capture{img/classe_fiche_fr.png}{0.7}

\paragraph{Mes modules et sauvegarde.} \emph{Mes modules} : créer, copier (y compris un module de ClassAction),
renommer, supprimer (avec \emph{Annuler}), ajouter ou retirer n'importe quelle question ou ressource, réordonner
(onglet \emph{Organiser} : glisser, ou Alt+$\uparrow$ et Alt+$\downarrow$). Tout est gardé dans \texttt{sauvegarde.json}, dans le dossier
de configuration (écriture atomique, copie de la version précédente) : réglages de chaque animation, réponses, scores
et tentatives, progression, Mes modules, préférences, au même format que Coupole Web. \emph{Ma progression} : bilan par
module, questions à revoir, \emph{Exporter} et \emph{Importer} (fusionner ou remplacer ; fichiers de Coupole Web
acceptés), \emph{Réinitialiser} tout ou partie (une animation, un module, les réglages, les réponses, Mes modules, les
préférences, tout), avec confirmation, copie gardée dans \texttt{sauvegardes/} et \emph{Annuler}.
\capture{img/classe_progression_fr.png}{0.95}

\paragraph{Orion en 3D, refaite.} Sur le modèle de la Grande Ourse 3D : la Terre, les étoiles principales
(Bételgeuse, Bellatrix, Meissa, la Ceinture, Saiph, Rigel, l'épée avec Hatysa et le Trapèze) à leur distance, les
lignes de visée et la figure sur la sphère céleste ; on fait tourner la scène au glisser, boutons \emph{De face (vue de
la Terre)} et \emph{De côté}. M42 (cloque creusée dans le nuage OMC-1, couche émissive, Trapèze, Voile, Barre) et
IC~434 (front rouge éclairé par $\sigma$~Ori, pilier de la Tête de Cheval, bord d'Orion~B, NGC~2023) sont dessinées en
volume d'après leur géométrie publiée, à leur position exacte, la taille agrandie $\times$12 (c'est écrit dans la scène).
Scène fixe : rien n'est redessiné tant que rien ne bouge.
\capture{img/anim_orion3d_cote_fr.png}{0.95}
\capture{img/anim_orion3d_face_fr.png}{0.95}

\paragraph{Spectres.} Les intensités des familles de raies sont calculées tous les huitièmes de type (même physique que
le modèle, étalonnage interpolé de façon monotone et lisse, PCHIP) : plus de paliers ni de cassures. La bande de
spectre a un \emph{contraste pédagogique} par défaut (raies élargies et approfondies, dit sous la bande) ou une
\emph{échelle réaliste}.
\capture{img/anim_spectrum010_g2_fr.png}{0.95}

\subsection{Présentation et vérifications (0.7.0)}\label{sec:finale}
\paragraph{Présentation.} L'accueil montre les modules en cartes de même largeur (numéro, titre, contenu, barre de
progression fine) ; la vue courante est mise en évidence dans la barre \emph{Modules}, \emph{Animations},
\emph{Mes modules}, \emph{Ma progression}. Sur une question, les ressources liées sont à droite en fenêtre large et
dans le menu « $\star$~n » de l'en-tête en fenêtre étroite ; la figure ne dépasse pas 45~\% de la hauteur, si bien que les
choix restent visibles en 1024~$\times$~600. Les textes des figures et des schémas restent lisibles dans les deux thèmes
et ne se recouvrent jamais, quelle que soit la police du système. L'interface s'adapte de 1024~$\times$~600 à la 4K et
de 100 à 200~\% de mise à l'échelle ; en petite fenêtre, la page d'une animation défile au lieu de tasser la scène.
\capture{img/classe_question_etroite_fr.png}{0.8}

\paragraph{Vérifications automatiques.} Les 517 questions sont notées une à une (bonne et mauvaise réponse), les 95
animations sont manipulées (chaque réglage aux bornes, chaque bouton, glissers aléatoires) sans exception et sans gel
du fil graphique, la sauvegarde fait l'aller-retour export $\to$ effacement $\to$ import à l'identique (et accepte un
export de Coupole Web), la langue et le thème changent à chaud, les réglages reviennent à la réouverture. Chaque paquet
construit (Windows, macOS, Linux) passe un essai « fumée » sans écran : \texttt{python -m coupole.essai\_fumee}.

\paragraph{Kit commun.} Coupole suit désormais le standard des applications de bureau de l'auteur : la mise à jour et
les rapports d'incident sont les modules communs \texttt{updater.py} et \texttt{reporting.py}, identiques d'une
application à l'autre et réglés par \texttt{kit.json} ; réglages, rapports et données vivent hors du dossier
d'installation (l'accord donné aux versions précédentes est repris) ; aucun nom de machine n'est envoyé (identifiant
aléatoire). \texttt{coupole -{}-check} vérifie une installation sans écran (fenêtre, modules, animations, données,
manuels, logo, style) et rend le code 0 si tout va bien. L'installeur Windows est bilingue ; le dernier installeur a
un lien permanent (\texttt{releases/latest/download/Coupole-Setup.exe}) ; sous Linux, l'archive
\texttt{Coupole-X-linux-x86\_64.tar.gz} (ou \texttt{-arm64}) s'installe dans le profil par \texttt{installer.sh}
(commande \texttt{coupole}), et la boîte de fichiers est par défaut celle, complète, de Qt.

\subsection{Recherche, favoris, bibliothèque, « Maintenant » (0.7.1)}\label{sec:favoris}
\paragraph{Recherche bilingue.} En français ou en anglais, quelle que soit la langue de l'interface, la recherche
trouve le même contenu (modules, animations, questions, fiches, tables, images) et l'affiche dans la langue de
l'interface : « seasons » donne « Simulateur des saisons (NAAP) », « Kepler's third law » la troisième loi de Kepler.
L'index réunit les titres, les mots-clés, les textes et les titres d'origine de ClassAction dans les deux langues ; il
ignore les accents et la casse, ramène les pluriels simples (saisons, galaxies, planètes) et tolère une faute (deux
pour les mots longs). Les résultats sont groupés par type, les mots trouvés en gras ; quand seul le titre de l'autre
langue correspond, il est rappelé (« En anglais : \dots{} »). Le moteur est le même, ligne pour ligne, que celui de
Coupole Web, avec les mêmes cas de référence.
\capture{img/classe_recherche_fr.png}{0.95}

\paragraph{Favoris et collections.} L'étoile $\star$ d'une page (module, question, fiche, table, image, animation),
le clic droit ou Ctrl+D dans une liste marquent un favori ; « Mes favoris » (barre du haut et accueil) les regroupe
par type, avec un filtre et un tri. Les collections (les « Mes modules » d'avant : vos propres séquences dans l'ordre
choisi) passent au second plan : bouton « Organiser en collections\dots{} ». Favoris et collections sont dans la
sauvegarde commune (format version 2), échangée avec Coupole Web.
\capture{img/classe_favoris_fr.png}{0.95}

\paragraph{Bibliothèque.} La présentation d'origine de ClassAction : un onglet par type (Questions, Animations, Fiches,
Tables, Images) montre tout le contenu de ce type, rangé par module dans l'ordre d'origine, avec des vignettes (les
animations dessinent elles-mêmes la leur) ; filtres par module et « favoris seulement ».
\capture{img/classe_bibliotheque_fr.png}{0.95}

\paragraph{Maintenant et en direct.} Les 19 animations qui dépendent de la date ou de l'heure ont les boutons
\emph{Maintenant} (instant présent de l'ordinateur, en temps universel, affiché dans votre fuseau : « 09:20 heure de
Paris (UTC+2) · 07:20 TU », sans erreur d'heure d'été), \emph{En direct} (l'heure suit l'horloge, sans redessiner
sans cesse ; s'arrête dès qu'on change la date ou qu'on relance la lecture) et \emph{Ma position} (un des sites
connus). Une animation réglée en heure solaire reçoit l'heure solaire moyenne au méridien du fuseau (hors heure
d'été), ou à la longitude du site choisi.
\capture{img/anim_maintenant_fr.png}{0.95}

\paragraph{Ombres d'éclipse en vraie géométrie.} Le type d'éclipse est calculé avec les tailles réelles et la distance
de la Lune telle qu'elle est placée : en l'approchant, bien alignée, on passe d'aucune éclipse à une partielle
(pénombre seule), une annulaire (antombre), puis une totale (ombre), avec la largeur de la bande ; la distance est
comparée au périgée et à l'apogée extrêmes (356\,375 et 406\,712~km, éphéméride DE440, 1900--2100). La vue de côté
trace les cônes de la Terre et de la Lune (tailles et écarts agrandis $\times$6, indiqué) ; la vue 3D montre l'orbite
inclinée de 5,1° et la ligne des nœuds, et pourquoi il n'y a pas d'éclipse à chaque nouvelle lune. « Rejouer une
éclipse » replace la Lune selon la géométrie réelle (JPL DE440 : distance, gamma, trajectoire) ; les 133 éclipses de
la table retrouvent leur type.
\capture{img/anim_shadowsim_fr.png}{0.95}

\subsection{Gravité et orbites (0.7.2)}\label{sec:gravite}
Une 95\textsuperscript{e} animation, dans le thème « Renaissance » et le module « L'astronomie de la Renaissance »,
d'après la simulation PhET « Gravity and Orbits » du projet PhET Interactive Simulations de l'University of Colorado
Boulder, montrée en cours. Seules l'idée et la liste des fonctions sont reprises : le code, les dessins et les textes
sont propres à Coupole, aucun fichier ni logo de PhET n'est utilisé (licences de PhET vérifiées le 15~octobre~2026 et
citées dans \texttt{NOTICE}).

\paragraph{Scénarios.} Soleil et Terre ; Soleil, Terre et Lune ; Terre et Lune ; Terre et station spatiale (ISS) ;
libre (Soleil, Terre, Lune et une comète d'aphélie 2,4~ua). Les couples partent sur des orbites circulaires ; le
Soleil, la Terre et la Lune ensemble partent de leurs positions et vitesses \emph{réelles} du 1\textsuperscript{er}
janvier 2026 (éphéméride JPL DE440, chaque orbite ramenée dans l'écliptique par une rotation autour de sa ligne des
nœuds) : une Lune lancée sur un cercle de Kepler prendrait, dans le champ de marée du Soleil, une orbite déformée dont
la période dépend de la phase de départ (27,05 à 27,81~j).

\paragraph{Deux vues.} « Modèle » : astres agrandis et distance Terre--Lune agrandie 78 fois, pour lire (les rapports
sont écrits sur la toile). « À l'échelle » : tailles et distances vraies ; vus du Soleil, la Terre et la Lune ne font
qu'un point (entouré d'un anneau discret), le Soleil n'est qu'un petit disque, et l'ISS rase la Terre. Le passage de
l'une à l'autre se fait en douceur (rayons et écart interpolés en logarithme, 0,9~s), puis plus rien n'est redessiné.
\capture{img/anim_gravite_modele_fr.png}{0.95}
\capture{img/anim_gravite_echelle_fr.png}{0.95}

\paragraph{Calcul.} N corps en gravitation newtonienne, schéma symplectique de Yoshida d'ordre~4 (trois pas de Verlet
de vitesse composés), pas adaptatif ($0{,}01\sqrt{r^3/G(m_1+m_2)}$ pour la paire la plus serrée) ; constantes : $G$
(CODATA 2018), $GM_\odot$ et $GM_\oplus$ nominaux (UAI 2015), $GM_{\mathrm{Lune}}$ (DE440), rayons (UAI 2015, NASA), unité
astronomique (UAI 2012), ISS (NASA : 419\,725~kg à 415~km). Collisions aux rayons réels : le corps le plus léger
s'écrase, la quantité de mouvement est gardée. Les nombres sont les mêmes, au bit près, que dans Coupole Web.

\paragraph{Réglages et souris.} Cases : force de gravité (flèches bleues, longueur relative à celle du départ pour
chaque corps), vecteur vitesse (flèches vertes, même échelle pour tous), trajectoires, grille (taille du carreau
écrite), ruban de mesure (vraie distance), gravité (sans elle, chaque corps file en ligne droite), Lune vue du Soleil ;
masses du Soleil, de la Terre et de la Lune (0,5 à 2 fois les vraies) ; boutons \emph{Pas}, \emph{Retour} (positions et
vitesses du dernier départ), \emph{Effacer les trajectoires}. Glisser un astre le déplace (ses satellites le suivent ;
le temps est suspendu pendant qu'on le tient), tirer le bout d'une flèche verte change la vitesse, glisser les bouts du
ruban mesure.
\capture{img/anim_gravite_iss_fr.png}{0.95}

\paragraph{Lune vue du Soleil.} Un panneau, à côté de la scène (fenêtre large) ou dessous (fenêtre haute), jamais par-dessus le mouvement, montre la vraie trajectoire de la Lune autour du Soleil sur les 60 derniers
jours, à la même échelle en largeur et en hauteur : elle se confond avec celle de la Terre et ne fait jamais de boucle
(les boucles de la vue « Modèle » viennent de l'agrandissement de la distance Terre--Lune). Dessous, le rapport des
courbures Lune~/ Terre reste toujours positif (0,49 à 1,41) : la trajectoire de la Lune est toujours tournée vers le
Soleil, qui la tire 2,2 fois plus fort que la Terre ($GM_\odot/d_\odot^2$ contre $GM_\oplus/d_{\mathrm{Lune}}^2$).

\paragraph{Valeurs et vérifications.} Distances, vitesses relatives (km·s$^{-1}$), périodes mesurées (passages
successifs par la direction de départ vue du corps central) comparées aux vraies ou à Kepler, énergie totale et moment
cinétique avec leur variation relative (de l'ordre de $10^{-14}$). Mesures : Terre 365,26~j (année sidérale
365,256~j), Lune 27,34~j (mois sidéral 27,322~j), ISS 92,87~min, comète 671,0~j ; une intégration indépendante sous
SageMath (DOP853, tolérance $10^{-13}$, mêmes conditions initiales) donne les mêmes périodes à $2\times10^{-6}$ près
(\texttt{outils/verif\_anim/gravite\_orbites\_verif.sage}).
\capture{img/anim_gravite_libre_fr.png}{0.95}

\subsection{Les animations disponibles}
\begin{center}
\small
\begin{tabularx}{\linewidth}{@{}>{\raggedright}p{0.30\linewidth}X@{}}
\toprule
Animation & Ce qu'elle montre, ce qu'elle calcule\\
\midrule
Longitude / latitude & globe (côtes, villes, équateur, méridien de Greenwich, ligne de changement de date) ; un point à déplacer, coordonnées en décimal ou sexagésimal\\
Mouvements du Soleil & sphère locale : course diurne, culmination, ombre, mois, analemme ; hauteur, azimut, $\alpha$, $\delta$, angle horaire, temps sidéral, équation du temps, lever et coucher ($h_0=-0{,}833^\circ$), durée du jour\\
Saisons (NAAP) & orbite (axe fixe dans l'espace), Terre de côté ou depuis le Soleil, point subsolaire, hauteur à midi $90^\circ-|\varphi-\delta|$, étalement d'un faisceau $1/\sin h$\\
Phases de la Lune (NAAP) & orbite vue du dessus, phase vue de la Terre, ciel de l'observateur ; fraction éclairée $(1-\cos\phi)/2$, âge (lunaison 29,53~j), pas de $\pm$1~j, 1~h, 1~min\\
Rayons du Soleil & rayons parallèles, latitude du rayon le plus direct ($\delta_\odot$), tropiques et cercles polaires, équinoxes et solstices 2026\\
Azimut / hauteur & repère horizontal : arcs de l'azimut (depuis le nord vers l'est) et de la hauteur ; étoile à glisser\\
Coordonnées équatoriales & sphère céleste, Terre au centre, arcs de $\alpha$ (depuis le cercle horaire 0~h) et de $\delta$, écliptique\\
Comparaison des systèmes & sphère céleste et diagramme d'horizon côte à côte ; régions circumpolaire, de lever et coucher, jamais levée ; traînées ; Grande Ourse, Orion, Croix du Sud\\
Temps sidéral et angle horaire & $H = \mathrm{TSL}-\alpha$ ramené entre $-12$ et $+12$~h, arc de l'angle horaire, le ciel tourne avec le TSL\\
Écliptique & carte $(\alpha,\delta)$ vue de l'intérieur ou de l'extérieur, écliptique $\tan\delta=\tan\varepsilon\sin\alpha$, trajet du Soleil entre deux dates\\
Obliquité & définition de l'angle ; Terre (IAU 2006, J2000 et 2026), Mars, Uranus\\
Écliptique (zodiaque) & Soleil devant les 13 constellations traversées (limites officielles de l'UAI, Ophiuchus compris) ; constellation de minuit ; signes de 30$^\circ$\\
Temps sidéral et solaire & orbite et deux horloges de 24~h ; jour sidéral $86\,400\,N/(N+1)$~s pour une année de $N$ jours (365, 365,25 ou 365,2422)\\
Plages de déclinaison & coupe Terre et sphère céleste ; limites $\delta=\pm(90^\circ-|\varphi|)$, hauteur du pôle $=|\varphi|$\\
Ombres d'éclipse & Terre et Lune à glisser, ombre et pénombre par les tangentes communes, type d'éclipse ; encart aux proportions réelles (méthode d'Aristarque)\\
Mars selon Ptolémée & déférent excentrique, point équant, épicycle parcouru en un an, Soleil ; boucles de la trajectoire\\
Épicycles & trois corps, l'un gardé fixe : héliocentrique et géocentrique donnent les mêmes positions relatives\\
Mouvement rétrograde & Terre et Mars sur leurs orbites, lignes de visée, Mars devant les étoiles (boucle avec l'inclinaison de 1,85$^\circ$) ; durée et amplitude de la rétrogradation\\
Orbite planétaire (NAAP) & lois de Kepler : ellipse et foyers, aires balayées, $P^2=a^3$, vitesse par la relation de vis-viva ; orbites des planètes\\
Parallaxe & $d=1/p$ ; conversions seconde d'arc, parsec, année-lumière, ua ; Proxima, Sirius, Véga\\
Module de distance & $m-M=5\log_{10}d-5$, une grandeur verrouillée, l'étoile sur une échelle de 1 à 1~000~pc\\
\midrule
Durée du jour & courbe de la durée du jour sur l'année pour une latitude, arcs de jour et de nuit du parallèle, moyenne annuelle ; horizon géométrique ou réfraction\\
Hauteur méridienne & coupe du méridien : $h = 90^\circ-\varphi+\delta$ (ou $180^\circ-h$ au nord du zénith), plages du Soleil ($\pm23{,}44^\circ$) et de la Lune ($\pm28{,}6^\circ$)\\
Trajectoires du Soleil & trajets diurnes aux solstices, aux équinoxes et le jour choisi, points de lever et de coucher, durées du jour\\
Vue d'ensemble du Soleil & construction en quatre étapes : pôles, équateur céleste, trajet des équinoxes, trajets des solstices\\
Position sur l'horizon & point de lever (ou de coucher) qui glisse sur l'horizon au fil de l'année, écart entre les solstices\\
Jour et nuit & carte plane : terminateur, point subsolaire (latitude $\delta_\odot$, midi solaire vrai), crépuscules civil, nautique, astronomique\\
Antipodes & un point et son antipode $(-\varphi, \lambda\pm180^\circ)$ sur le globe et la carte, déplacement contraint, multiples de 5$^\circ$\\
Grande Ourse en 3D & les sept étoiles à leur distance (81 à 123~al, parallaxes adoptées par SIMBAD) en couleur vraie, lignes de visée, projection sur la sphère céleste ; fiche au survol\\
Repères céleste et horizontal & passage d'un repère à l'autre par une rotation de $90^\circ-\varphi$ autour de l'axe est-ouest\\
\bottomrule
\end{tabularx}
\end{center}

\begin{center}
\small
\begin{tabularx}{\linewidth}{@{}>{\raggedright}p{0.30\linewidth}X@{}}
\toprule
Animation & Ce qu'elle montre, ce qu'elle calcule\\
\midrule
Vocabulaire des phases & la Lune à chacune des huit phases nommées, lunaison animée (29,53~s à l'écran : une seconde par jour), âge et jours écoulés\\
Phases avec un ballon & lampe, ballon et œil vus de dessus ; ce que voit l'œil, angle de phase $i$, fraction $(1+\cos i)/2$, phase lunaire équivalente\\
Trois points de vue & la Lune vue de la Terre, du Soleil (toujours pleine) et d'un œil qu'on place librement\\
Positions et phases & deux corps à placer autour du Soleil, chacun vu de l'autre ; tout près, l'un devient la lune de l'autre\\
Phases et diagramme d'horizon & Soleil et Lune sur huit positions de l'équateur céleste : l'heure, la phase et la position dans le ciel\\
Questionnaire des phases & on cache la phase, la Lune ou le Soleil ; réponse au curseur ou à la souris, verdict (même phase nommée)\\
Plans bissecteurs & plans perpendiculaires au Soleil et à la ligne Terre--Lune : la phase est la partie éclairée de la moitié tournée vers nous\\
Déclinaison de la pleine lune & opposée au Soleil à $\pm5{,}145^\circ$ près ; courbe annuelle, hauteur à minuit $90^\circ-\varphi+\delta$\\
Retard synodique & jour sidéral et jour solaire : $1/S = 1/P_\text{rot} - 1/\text{année}$ (rotation réelle ou ralentie à 240~h ou 1~200~h)\\
Marées & marée d'équilibre de la Lune et du Soleil $\zeta=(M/M_\oplus)(R/d)^3R\,(3\cos^2\psi-1)/2$, force de marée, vives et mortes eaux, niveau chez un observateur\\
Inclinaison de l'orbite lunaire & orbite inclinée de 5,145$^\circ$, mois draconitique 27,212~j, année draconitique 346,62~j ; vue du Soleil à l'échelle, ombre et pénombre\\
Table des éclipses & les 133 éclipses de 2015 à 2044 en bandes obliques ; saisons d'éclipses (Soleil à un nœud), vérification par les phases moyennes de Meeus\\
\bottomrule
\end{tabularx}
\end{center}

\begin{center}
\small
\begin{tabularx}{\linewidth}{@{}>{\raggedright}p{0.30\linewidth}X@{}}
\toprule
Animation & Ce qu'elle montre, ce qu'elle calcule\\
\midrule
Phases de Vénus & système de Copernic : élongation, phase, diamètre apparent (de 10 à 60$''$), plus grande élongation $\arcsin a$\\
Vénus selon Ptolémée & épicycle sur la ligne Terre--Soleil : jamais plus qu'un croissant\\
Configurations planétaires & conjonctions, oppositions, quadratures, plus grandes élongations ; période synodique, arrêt ou pause à chaque configuration\\
Période synodique & $1/S=1/E-1/P$ (supérieure) ou $1/P-1/E$ (inférieure), $E$ année sidérale (365,256~j)\\
3$^\text{e}$ loi de Kepler & $P^2=a^3/M$, graphique logarithmique, les planètes sur la droite\\
Excentricité & $e=c/a$, $b=\sqrt{a^2-c^2}$, somme des distances aux foyers $=2a$\\
Algèbre de la gravitation & facteurs $\frac13$ à 3 sur $M_1$, $M_2$, $d$ : force multipliée par $ab/c^2$ (fraction exacte)\\
Loi de la gravitation & $F=GM_1M_2/R^2$ de l'électron à la Voie lactée, accélérations, préréglages, mémoire des résultats\\
Orbites autour de la Terre & $P=2\pi\sqrt{r^3/GM}$ : Station spatiale, GPS, géostationnaire (42~164~km), Lune\\
Gravité et orbites & d'après PhET : Soleil, Terre, Lune, ISS, comète en gravitation (Yoshida d'ordre 4, état réel DE440 du 1/01/2026) ; vues modèle et à l'échelle, forces, vitesses, trajectoires, ruban, gravité coupable, périodes mesurées, Lune vue du Soleil\\
\midrule
Petits angles & $\alpha\approx206\,265''\,d/D$ contre la valeur exacte $2\arcsin(d/2D)$\\
Horloge de la Grande Ourse & ciel du nord en projection stéréographique ; un tour par jour sidéral, environ 1$^\circ$ par jour à heure fixe\\
Lever héliaque & premier matin où l'étoile est levée à l'aube (Soleil à $-7^\circ$) ; année entière, Sirius, Véga…\\
Ombre d'un bâtiment & au lieu du film d'origine : l'ombre portée jour après jour, à heure fixe ou au midi vrai\\
\bottomrule
\end{tabularx}
\end{center}

\begin{center}
\small
\begin{tabularx}{\linewidth}{@{}>{\raggedright}p{0.30\linewidth}X@{}}
\toprule
Animation & Ce qu'elle montre, ce qu'elle calcule\\
\midrule
Luminosité d'une étoile & $L/L_\odot=(R/R_\odot)^2(T/T_\odot)^4$, $L=4\pi R^2\sigma T^4$\\
Flux & une ampoule et un détecteur : $F=P/4\pi d^2$\\
Vitesse d'une étoile & $v_r=c\,\Delta\lambda/\lambda_0$ (H$\alpha$), $v_t=4{,}74\,\mu d$, $v=\sqrt{v_r^2+v_t^2}$\\
Parallaxe spectroscopique & type et classe → $M_V$ publiée (ou ancien étalonnage, signalé) ; $d=10^{(m-M+5)/5}$~pc ; spectre calculé\\
Diagramme HR (NAAP) & étoiles proches et brillantes, séquences, droites d'égal rayon, $R/R_\odot=\sqrt{L/L_\odot}/(T/T_\odot)^2$\\
\bottomrule
\end{tabularx}
\end{center}

\begin{center}
\small
\begin{tabularx}{\linewidth}{@{}>{\raggedright}p{0.30\linewidth}X@{}}
\toprule
Animation & Ce qu'elle montre, ce qu'elle calcule\\
\midrule
Clou chauffé jusqu'à la fusion & au lieu du film : couleur et éclat du clou par la loi de Planck et les fonctions colorimétriques CIE (invisible sous $\approx800$~K) ; courbe de Planck, $\lambda_\text{max}=b/T$, part rayonnée dans le visible\\
Courbes du corps noir (NAAP) & jusqu'à cinq courbes, maximum (Wien), aire $\sigma T^4/\pi$, couleur perçue ; magnitudes U, B, V, R et indices de couleur\\
Explorateur de spectres & continu, émission, absorption ; six familles de raies (He~\textsc{ii}, He~\textsc{i}, H, métaux ionisés et neutres, molécules) calculées par Saha–Boltzmann selon le type spectral et la classe ; 133 raies NIST, bandes TiO et CH\\
Spectre électromagnétique & $\nu=c/\lambda$, $E=h\nu$, corps noir dont c'est le maximum, domaine ; l'onde dessinée à côté d'un objet de même taille\\
Effet Doppler & crêtes émises par une source en mouvement, signal reçu par le temps retardé $|\vec r_\text{obs}-\vec r_s(\tau)|=c(t-\tau)$ : $f=f_0/(1-v_r/c)$ ; observateur à glisser\\
Filtres (NAAP) & source (corps noir, gaussienne, Soleil) × filtres (U, B, V, R, atmosphère, gaussien) = spectre reçu ; couleur perçue (CIE, sRGB)\\
Atome d'hydrogène (NAAP) & photon absorbé si $E=13{,}598\,(1/n^2-1/m^2)$~eV, ionisation au-delà, réémission ; niveaux d'énergie, journal\\
Les trois spectres (Kirchhoff) & ampoule, nuage de gaz, télescope à glisser : continu, émission ou absorption aux mêmes longueurs d'onde\\
\midrule
Lanceur de marteau & au lieu du film : athlète et marteau (7,26~kg), étoile et planète, autour du centre de masse ; $r_1/r_2=m_2/m_1$\\
Courbe de vitesse radiale & Soleil et Jupiter ($K=12{,}5$~m/s, 11,86~ans), Soleil et Terre, 51~Pegasi~b ; $K=(2\pi G/P)^{1/3}M_p\sin i/(M_\ast+M_p)^{2/3}$\\
Le Soleil et les planètes & centre du Soleil autour du barycentre (positions vraies des planètes, éléments de Standish, masses DE440, 1800--2050) ; jusqu'à 2,2~$R_\odot$\\
Période d'un pulsar & impulsions émises le long d'une orbite : intervalles d'arrivée $P(1\pm v/c)$ (temps de trajet exact)\\
Décalage des raies & étoile et planète vues de dessus ; raies du spectre décalées de $\Delta\lambda/\lambda=v/c$ (amplifié), décalage réel de H$\alpha$\\
Transits (NAAP) & profondeur $(R_p/R_\ast)^2$ (aire exacte de recouvrement), durée, paramètre d'impact ; 9 planètes réelles (NASA Exoplanet Archive) ; mesures bruitées\\
Vitesses radiales (NAAP) & $v_r=K[\cos(\nu+\omega)+e\cos\omega]$, orbites excentriques et inclinées, $M_p\sin i$ ; planètes réelles, mesures bruitées\\
\bottomrule
\end{tabularx}
\end{center}

\begin{center}
\small
\begin{tabularx}{\linewidth}{@{}>{\raggedright}p{0.30\linewidth}X@{}}
\toprule
Animation & Ce qu'elle montre, ce qu'elle calcule\\
\midrule
Binaires à éclipses (NAAP) & deux étoiles et la courbe de lumière (filtre V, brillance par la loi de Planck) ; 46 binaires réelles et leurs mesures (CALEB), dont 7 « à compléter » ; écart modèle--mesures\\
Courbe de lumière & un couple de période 3,65~j observé mesure après mesure, chaque point reporté à sa phase\\
Centre de masse & $r_1=d\,m_2/(m_1+m_2)$, $m_1r_1=m_2r_2$ ; le couple tourne autour\\
Repliement (roue de charrette) & roue vue image par image : rotation apparente ramenée dans $]-180^\circ/n,180^\circ/n]$, roue arrêtée ou à l'envers ; même piège pour une étoile variable\\
\midrule
Propriétés des planètes & histogrammes : distance, période, masse, rayon, satellites (Jupiter 95, Saturne 274), rotation, densité (NASA, 2026)\\
Températures de formation & $T=600~\text{K}\,(d/1~\text{ua})^{-0{,}747}$ ; état (gaz ou solide) de neuf substances, ligne des glaces\\
La neige vue d'une voiture & les flocons semblent jaillir d'un point : radiant à $\arctan(v_\text{neige}/v_\text{voiture})$, comme une pluie d'étoiles filantes\\
Chaîne proton-proton, cycle CNO & au lieu des films : chaque réaction, énergie libérée (AME 2020), bilan 26,731~MeV, part des neutrinos\\
Orion en 3D & la constellation dépliée à ses vraies distances, barres d'erreur, nuages Orion A et B, M42 en « blister », Tête de Cheval, $\lambda$~Ori, boucle de Barnard ; plongée avec les images ESO et Hubble ; fiches au survol\\
Rotation de la Voie lactée & modèle de McMillan (2017) et ses composantes, mesures d'Eilers et al. (2019) et galkin, $M=v^2R/G$, courbe képlérienne ; Galaxie vue du pôle (barre, bras de Reid et al. 2019)\\
Ondes de densité & tracteur et voitures : $n_\text{bouchon}/n=(v-v_t)/(v_\text{lente}-v_t)$, analogie des bras spiraux\\
Univers-ballon & au lieu du film : galaxies sur un ballon qui gonfle, $v=Hd$, $H=\dot a/a$ ; le ballon se tourne à la souris\\
Décalage vers le rouge & spectre de M32 décalé de $1+z$, raies, lumière reçue dans U, B, V, R ; vitesse par l'effet Doppler relativiste\\
L'atmosphère terrestre & masse d'air (sec $z$, Kasten \& Young, Pickering), extinction $\Delta m = kX$ et rougissement, fenêtres atmosphériques (ESO SkyCalc), réfraction, Soleil aplati, lever avancé, dispersion\\
Télescope & grossissement $F/f$, champ réel, pouvoir séparateur $1{,}22\lambda/D$, pouvoir collecteur, mise au point ; Lune, Saturne, amas globulaire\\
Capteur CCD & pose (photons au hasard, loi de Poisson, saturation) puis lecture ligne par ligne et case par case\\
Loi de Snell-Descartes & $n_1\sin\theta_1=n_2\sin\theta_2$, réflexion totale, part réfléchie (Fresnel), angle de Brewster\\
Regard vers le passé & l'explosion d'une étoile à $d$ années-lumière vue $d$ ans plus tard ; supernovæ de 1054 et de 1006\\
\bottomrule
\end{tabularx}
\end{center}

\capture{img/anim_eclipticsimulator_fr.png}{0.95}
\capture{img/anim_lunar_applet_fr.png}{0.95}
\capture{img/anim_shadowsim_fr.png}{0.95}
\capture{img/anim_celhorcomp_fr.png}{0.95}
\capture{img/anim_kepler_fr.png}{0.95}
\capture{img/anim_marsorbit_fr.png}{0.95}
\capture{img/anim_longlat_fr.png}{0.95}
\capture{img/anim_eclipsetable_fr.png}{0.95}
\capture{img/anim_mooninc_fr.png}{0.95}
\capture{img/anim_tidesim_fr.png}{0.95}
\capture{img/anim_venusphases_fr.png}{0.95}
\capture{img/anim_configurationssimulator_fr.png}{0.95}
\capture{img/anim_heliacalrisingsim_fr.png}{0.95}
\capture{img/anim_transitmovie_fr.png}{0.95}
\capture{img/anim_hrexplorer_fr.png}{0.95}
\capture{img/anim_bbexplorer_fr.png}{0.95}
\capture{img/anim_hydrogenatom_fr.png}{0.95}
\capture{img/anim_threeviewsspectra_fr.png}{0.95}
\capture{img/anim_transitsimulator_fr.png}{0.95}
\capture{img/anim_ca_extrasolarplanets_starwobble_fr.png}{0.95}
\capture{img/anim_eclipsingbinarysim_fr.png}{0.95}
\capture{img/anim_galacticredshift_fr.png}{0.95}
\capture{img/anim_milkywayrotationalvelocity_fr.png}{0.95}
\capture{img/anim_telescope10_fr.png}{0.95}
\capture{img/anim_fusion02_fr.png}{0.95}

\subsection{Calculs et vérification}
Les animations d'origine employaient des approximations commodes (obliquité de 23,5$^\circ$, orbite circulaire,
ajustement de Fourier pour le Soleil, 3,26~al par parsec). Coupole les remplace par :
\begin{itemize}
\item jour julien et temps sidéral moyen de Greenwich (IAU 1982, Meeus éq.~12.4) ; obliquité moyenne P03 / IAU 2006
      ($\varepsilon_{J2000}=84\,381{,}406''$) ;
\item position apparente du Soleil et équation du temps (Meeus chap.~25 et 28, précision $\approx 0{,}01^\circ$) ;
\item changements de repère par trigonométrie sphérique (azimut compté depuis le nord vers l'est) et matrices de
      rotation pour le dessin ;
\item équation de Kepler (Newton), relation de vis-viva ($GM_\odot = 1{,}327\,124\,400\,18\times10^{20}$~m$^3$\,s$^{-2}$),
      3$^\text{e}$ loi $P^2=a^3$ ;
\item parsec et année-lumière exacts de l'UAI ($1~\mathrm{pc}=3{,}261\,564$~al) ; rayons nominaux de l'UAI.
\item phases moyennes de la Lune et argument de latitude $F$ (Meeus chap.~49 et 54), nœud de l'orbite lunaire
      (chap.~47) ; précession rigoureuse IAU 1976 (chap.~21) ; marée d'équilibre ($M_\oplus/M_{\mathrm{Lune}}=81{,}300\,568$) ;
\item périodes synodiques, $G=6{,}674\,30\times10^{-11}$ (CODATA 2018), $GM_\oplus=3{,}986\,004\,418\times10^{14}$ ;
      Stefan-Boltzmann ($\sigma$ exacte) et constantes solaires nominales de l'UAI ($R_\odot$, $L_\odot$, d'où
      $T_\odot=5\,772$~K ; $M_{\text{bol},\odot}=4{,}74$ et facteur 2,5 de Pogson, l'original écrivait 2,51189). Les
      étalonnages empiriques de l'original (type spectral, corrections bolométriques, séquences du diagramme HR,
      intensités des raies) sont repris tels quels ;
\item (lot~3) constantes exactes du SI ; loi de Planck, constante de Wien tirée de $x=5(1-e^{-x})$
      ($b=2{,}897\,772$~mm\,K ; l'original : 2,897~768), $\sigma$ ; hydrogène avec la masse réduite
      ($13{,}598$~eV, $hc=1\,239{,}84$~eV\,nm ; l'original : 13,6 et 1~240) ; Snell-Descartes et Fresnel ; effet
      Doppler par le temps retardé ; couleur perçue par les fonctions CIE 1931 (Judd-Vos) et la transformation sRGB
      normalisée ; magnitudes dans les filtres de Landolt ;
\item (lot~3) deux corps : $P=2\pi\sqrt{a^3/G(M_1+M_2)}$ ($0{,}1158$~j pour $a=1~R_\odot$, l'original : 0,1155),
      équation de Kepler, $K$ et $v_r$, aire exacte de recouvrement de deux disques (transits, éclipses) ; masses par
      les $GM$ de l'UAI et de DE440, barycentre du Système solaire par les éléments de Standish ; bilans des réactions
      de fusion par les excès de masse AME 2020 ; $G$ en kpc\,(km/s)$^2/M_\odot$.
\end{itemize}
Le script \texttt{outils/verif\_anim/verif\_anim.sage} recalcule chaque grandeur sous SageMath en 100~bits, par une
autre voie que le code (rotations au lieu des formules fermées, racines numériques, triangles semblables) ; ses
valeurs sont figées dans \texttt{tests/anim\_reference\_sage.json} et \texttt{tests/test\_anim\_calculs.py} y compare
Coupole (écarts $\le 10^{-9}$), ainsi qu'aux exemples publiés de Meeus (temps sidéral du 10 avril 1987,
Soleil du 13 octobre 1992) et à astropy (Soleil apparent à mieux que $0{,}01^\circ$, limites des constellations à
$0{,}05^\circ$). Quelques valeurs obtenues : éclipse centrale de Lune, ombre de 4~607~km de rayon à la distance de
la Lune, soit 2,65 rayons lunaires, totalité de 1~h~41~min ; rétrogradation de Mars (orbites circulaires) de 73~jours
sur 15,9$^\circ$ ; jour sidéral de 23~h~56~min~4,091~s.

Lot~2 : les 133 éclipses de la table tombent à moins de 1,01~jour d'une phase moyenne, avec $|\sin F|\le0{,}28$, et
toutes dans une saison d'éclipses ; en 2026, le Soleil passe aux nœuds le 27~février et le 22~août ; lever héliaque
de Sirius à 40,8$^\circ$~N le 11~août (Véga le 14~novembre, Aldébaran le 19~juin) ; orbite géostationnaire de
42~164~km ; marée d'équilibre de la Lune de 53,8~cm d'amplitude totale, Soleil 0,459 fois la Lune ; précession à
$10^{-4}$~degré d'astropy.

Lot~3 : constante de Wien, $\sigma$, loi de Planck à $10^{-11}$ près de SageMath, et $\sigma T^4$ égal à
l'intégrale numérique de $\pi B_\lambda$ ; H$\alpha$ à 656,47~nm dans le vide ; Fresnel identique aux lois des sinus et
des tangentes, $r_p$ nul à l'angle de Brewster ; fréquence Doppler égale à $f_0/(1-v_r/c)$ à $10^{-5}$ ; Soleil autour
du barycentre à 1,65~$R_\odot$ le 1\textsuperscript{er}~janvier 2000, à mieux que $5\times10^{-5}$~ua d'ERFA ; Jupiter
fait bouger le Soleil de 12,5~m/s ; HD~209458~b : transit de 1,44~\% et 3~h~05 ; les 39 binaires complètes suivent leurs
mesures à 0,037 près (médiane, flux relatif), et DI~Her « à compléter » retrouve $e=0{,}49$ par ajustement ; fusion
26,731~MeV ; masse de la Galaxie dans l'orbite du Soleil $9{,}9\times10^{10}~M_\odot$ ; le bouchon des ondes de densité
atteint la densité prévue ($\times5$).  Les courbes de lumière des binaires, dessinées en images dans l'original, ont été
numérisées par \texttt{outils/verif\_anim/donnees\_lot3.py}.

\begin{ccode}
coupole anim liste\\
coupole anim liste --faites --json\\
coupole anim info sunmotions
\end{ccode}


\paragraph{L'atmosphère terrestre (0.8.0).} Animation propre à Coupole (thème Télescopes), quatre vues : épaisseur
traversée (masse d'air : $\sec z$, Kasten \& Young 1989, Pickering 2002, coquille sphérique de 8~km, $\approx 38$ à
l'horizon), extinction et rougissement ($\Delta m = kX$, Rayleigh $\propto \lambda^{-4}$, aérosols, ozone), fenêtres
atmosphériques (modèle de ciel de l'ESO de 0,3 à 30~µm, schéma des rayons $\gamma$ à la radio) et réfraction (Bennett,
Sæmundsson, Soleil aplati à l'échelle, lever avancé, dispersion d'après Ciddor 1996, pression et température).
Formules vérifiées par SageMath (\texttt{outils/verif\_atmosphere/}).
\capture{img/anim_earthatmosphere_fenetres_fr.png}{0.9}
\capture{img/anim_earthatmosphere_refraction_fr.png}{0.9}

\paragraph{Temps sidéral et angle horaire : l'hémisphère Sud.} L'arc de l'angle horaire est affiché par défaut (un
choix déjà fait est gardé). La fiche « Ce qu'il faut voir » porte un encadré \emph{Dans l'hémisphère Sud} :
$H = \mathrm{TSL} - \alpha$ ne change pas (compté vers l'ouest depuis le méridien, $H < 0$ avant la culmination), mais
la plupart des astres culminent vers le nord ($\delta > \varphi$) ; face au nord, le ciel défile de droite à gauche ;
le pôle Sud céleste n'a pas d'étoile brillante ($\sigma$~Octantis, magnitude 5,4, à environ 1°) ; une étoile est
circumpolaire si $\delta < -(90^\circ - |\varphi|)$. Ces affirmations sont vérifiées par SageMath
(\texttt{outils/verif\_anim/temps\_sideral\_sud\_verif.sage}).

% =====================================================================
\section{Calculateurs de TD}\label{sec:td}
Des calculateurs qui suivent \emph{pas à pas} la démarche d'un TD ou d'un TP du cours, avec vos nombres, puis montrent
le calcul exact et ce qui change. Le premier : \textbf{Parallaxe (TP du cours)}, qui reprend le sujet officiel
(altitude de l'ISS, distance Terre-Lune).

\capture{img/td_parallaxe_iss_fr.png}{1.0}
\paragraph{La démarche du sujet.}
\begin{enumerate}
\item relever dans Stellarium les élévations $h_1$ et $h_2$ de l'objet vu des deux lieux, au même instant (et, si
possible, les azimuts) — pour la Lune, la distance angulaire Lune–étoile $\varphi$ du sujet convient aussi ;
\item parallaxe $2\alpha = h_1 - h_2$ (objet vu du même côté par les deux) ; si l'objet est \emph{entre} les deux lieux,
les visées se croisent au-dessus de la base et $2\alpha = 180^\circ - h_1 - h_2$ ;
\item distance $R$ entre les lieux « par exemple avec Google Earth » (à la surface) ou avec \texttt{EarthLocation}
d'astropy (en ligne droite : la corde) ;
\item $2\alpha = R/D$, $2\alpha$ en radians, d'où $D$.
\end{enumerate}
Les positions s'écrivent \emph{comme on veut} : décimal avec point ou virgule, signé ou avec N/S/E/O/W, degrés-minutes-%
secondes avec ° ′ ″ ' " ou des espaces (« 48°47′6″N », « 48 47 06 N », « 2,45° E », « -2.35 », « 48d47m6s »),
latitude et longitude dans n'importe quel ordre si les lettres le disent (sans lettre : latitude puis longitude). Une
latitude impossible ($> 90^\circ$) est corrigée et signalée ; une base énorme qui deviendrait raisonnable en échangeant
latitude et longitude d'un des lieux est signalée. Minutes ou secondes $\ge 60$, signe et lettre contradictoires :
erreur dite, rien n'est deviné.

\paragraph{Corde, surface et flèche.} La corde $R = \lvert\vec r_1 - \vec r_2\rvert$ relie les positions géocentriques
WGS84 (comme \texttt{EarthLocation.from\_geodetic}) ; la distance à la surface est la géodésique de l'ellipsoïde
(Vincenty). La corde passe \emph{sous} le sol : sa \textbf{flèche} (écart au milieu entre la corde et la surface) vaut
$\approx R^2/(8R_\oplus)$ — 3~km pour 400~km — et le calculateur la donne exactement sur l'ellipsoïde, avec la
profondeur de la corde sous l'horizon de chaque lieu ; une loupe la rend visible sur le dessin.

\capture{img/td_parallaxe_lune_fr.png}{1.0}
\paragraph{Le calcul exact.} Les deux horizons ne sont pas parallèles : ils font entre eux l'angle au centre $\theta$
($5{,}9^\circ$ pour 660~km). Avec les deux azimuts, les deux visées sont triangulées dans l'espace (point le plus proche
des deux droites) : distances aux deux lieux, distance au centre de la Terre, altitude au-dessus de l'ellipsoïde ;
sans azimut, l'objet est supposé dans le plan vertical des deux lieux. L'exemple Lune le montre : $h_1 - h_2 = 4{,}45^\circ$
est presque entièrement dû à $\theta$, la formule du sujet appliquée aux élévations donne $\approx 8\,500$~km, le calcul
exact $357\,917$~km — c'est pourquoi le TP de la Lune mesure la distance Lune–étoile.

\paragraph{Éphémérides et incertitudes.} Avec l'heure (bouton \emph{Maintenant}) : Lune par JPL DE440 (32~Mo,
téléchargée une fois, puis hors ligne) ou Meeus à défaut — distances topocentriques et géocentrique, élévations
prédites sans et avec réfraction ; ISS par le TLE de CelesTrak (le plus proche de l'instant parmi ceux gardés en
cache, date affichée, avertissement au-delà de 3 jours ; l'historique des TLE n'est pas librement disponible sans
compte Space-Track : pour un passage ancien, collez un TLE de l'époque) et SGP4. Incertitudes propagées sur $D$
(élévations, positions, heure multipliée par la vitesse de l'objet en élévation). \emph{Exemple} remplit des données
fictives cohérentes ; \emph{Exporter le détail} écrit un CSV pour la feuille « Calculs » du partage.

\paragraph{Vérification.} Lecteur de coordonnées : une batterie de cas commune avec Coupole Web
(\texttt{tests/td\_coordonnees\_cas.json}). Géométrie et calculs contre SageMath à 200 bits
(\texttt{outils/verif\_td/verif\_td.sage}) : corde, surface, flèche, profondeurs, angle au centre à 1~mm et
$10^{-9}$~degré près, triangulation, formule du sujet.

\begin{ccode}
coupole td parallaxe --p1 "48.8566, 2.3522" --h1 62,483 --az1 14,184 --alt1 35\\
\hspace*{2em}--p2 "45°45′50″N 4°50′09″E" --h2 33.603 --az2 346.928 --alt2 170
\end{ccode}

% =====================================================================
\section{Machine et performance}
\capture{img/module_machine_fr.png}{0.9}
Détection sans dépendance obligatoire : système, architecture, processeur, cœurs physiques et logiques (bornés par
l'affinité et le quota du conteneur), mémoire totale et disponible, place libre, cartes graphiques (nvidia-smi,
\texttt{/sys/class/drm}, system\_profiler, CIM). Parallélisme : téléchargements 3 (4 au plus), conversions = cœurs
physiques − 1 bornées par la mémoire disponible à raison de 500 Mo par conversion (pic mesuré sur une image IRIS
4096², la moitié de la mémoire libre au plus) ; mode économe (1 conversion, 2 téléchargements) automatique sous 4 Go
ou 2 cœurs ; bridage manuel, lui aussi borné par la mémoire. Exemples : 2 cœurs et 2 Go → 2 téléchargements,
1 conversion ; 6 cœurs et 20 Go libres → 3 et 5 ; 32 cœurs et 60 Go libres → 3 et 16 (plafond). La carte
graphique n'est pas utilisée par la Banque OHP (travail limité par le réseau, le disque et Zstandard) ;
\texttt{coupole.core.calcul} offre CuPy aux modules futurs.

Rien de lourd ne s'exécute dans le fil graphique : astropy, scipy et SEP ne sont chargés qu'au premier besoin, les
sondes de la machine et la détection d'ASTAP tournent en fond, et la fermeture de la fenêtre arrête puis attend tous
les fils, tâches et processus (voir \texttt{docs/AUDIT\_2026-10.md}). Les grandes listes restent fluides à 80~000
lignes : les cellules ne sont calculées que pour les lignes affichées, le tri se fait sur une clé par ligne (et non par
comparaisons deux à deux), la sélection ne fait plus parcourir toutes les lignes à chaque dessin, et le changement de
thème ne recalcule que les couleurs (voir \texttt{docs/AUDIT2\_2026-10.md}).

\begin{center}\small
\begin{tabular}{@{}lr@{}}
\toprule
\textbf{Mesure (processeur à 6 cœurs, NVMe ; conteneur Debian)} & \textbf{Valeur}\\
\midrule
\texttt{coupole --version} & 0,1–0,2 s\\
fenêtre affichée (astropy, scipy, SEP non chargés) & 0,42 s (1,5 s avant l'audit)\\
inventaire chargé (fil de fond) & 3,3 s au premier lancement, 0,8 s ensuite\\
chaîne complète, 24 images T120 de 8 Mo, plafond 8 Mo/s & 24,5 s = durée du téléchargement seul (24,2 s)\\
même chaîne sans plafond (serveur local) & 2,3 s, 86 Mo/s de FITS, 10 images/s\\
conversion d'une image T120 (un processus) & 0,13 s (zstd 14 ms, lecture 18 ms, écriture et relecture 29 ms)\\
pic mémoire d'une conversion T120 / IRIS 4096² (processus dédié) & 163 Mo / 365 Mo (170 / 573 avant)\\
interface pendant 3 téléchargements (28 Mo/s) & écart max. du minuteur 20 ms : 20,8 ms\\
\midrule
\multicolumn{2}{@{}l}{\emph{Second audit (0.1.4) : banque réelle de 7~625 XISF ; partage réseau simulé à 2 ms par accès}}\\
rangement final sans rien à déplacer, sur le partage & 1,5 s (60,7 s avant)\\
48 images vers une destination sur le partage & 4,9 s (20,3 s avant)\\
qualité, relance de la banque déjà mesurée (partage) & 1,9 s (72 s avant)\\
réorganiser la banque (lecture des 7~625 en-têtes) & 57 s (243 s avant)\\
tout sélectionner, 80~000 images (10 × la banque) & 0,2 s (3,1 s avant ; 0,5 s → 0,17 s sur la banque)\\
trier 80~000 images & 0,1–0,2 s (3,7–4,1 s avant)\\
\bottomrule
\end{tabular}
\end{center}

% =====================================================================
\section{Ligne de commande}
Tout ce que fait l'interface se fait en ligne de commande ; \cmd{--json} pour les scripts. Sans argument,
\cmd{coupole} ouvre l'interface. Les commandes ont un nom français et un alias anglais.
\begin{ccode}
coupole ohp estimer "NGC 6888" --nuit 2025-07-16\\
coupole ohp traiter "NGC 6888" --nuit 2025-07-16 --filtre B --filtre V --filtre R --dest \textasciitilde/OHP --oui\\
coupole ohp telecharger Pluton --sans-dates-douteuses --dest \textasciitilde/FITS\\
coupole qualite \textasciitilde/OHP --ecrire
\end{ccode}
{\small% Généré par outils/tables_cli.py — ne pas modifier à la main.
\begin{tabularx}{\linewidth}{@{}>{\ttfamily\small\raggedright}p{0.38\linewidth}X@{}}
\toprule
\multicolumn{2}{@{}p{\linewidth}@{}}{\sffamily\bfseries coupole}\\
\midrule
--version & affiche la version et quitte\\
--lang,\allowbreak{} --langue \{auto,\allowbreak{}fr,\allowbreak{}en\} & langue des messages (auto : celle du système)\\
-v,\allowbreak{} --verbeux,\allowbreak{} --verbose & journal détaillé\\
--reinitialiser-interface,\allowbreak{} --reset-interface & efface la disposition gardée (fenêtre, colonnes, filtres, dossiers des dialogues) puis quitte ; les réglages ne changent pas\\
\bottomrule
\end{tabularx}
\par\medskip
\begin{tabularx}{\linewidth}{@{}>{\ttfamily\small\raggedright}p{0.38\linewidth}X@{}}
\toprule
\multicolumn{2}{@{}p{\linewidth}@{}}{\sffamily\bfseries coupole astap\newline\normalfont\small Cherche ASTAP et son catalogue d'étoiles, indique quoi installer pour ce système, ou enregistre leur emplacement. Coupole fonctionne sans ASTAP.}\\
\midrule
--guide & affiche le guide d'installation pour ce système\\
--definir,\allowbreak{} --set CHEMIN & chemin de l'exécutable ASTAP (ou de son dossier)\\
--catalogue,\allowbreak{} --catalog DOSSIER & dossier du catalogue d'étoiles (D80, D50...)\\
--oublier,\allowbreak{} --forget & oublie les chemins enregistrés (retour à la détection)\\
--json & sortie JSON (pour les scripts)\\
\bottomrule
\end{tabularx}
\par\medskip
\begin{tabularx}{\linewidth}{@{}>{\ttfamily\small\raggedright}p{0.38\linewidth}X@{}}
\toprule
\multicolumn{2}{@{}p{\linewidth}@{}}{\sffamily\bfseries coupole rapports (reports)\newline\normalfont\small Rien n'est envoyé sans accord. Sans accord, les rapports restent dans un dossier local.}\\
\midrule
--activer,\allowbreak{} --enable & autorise l'envoi des rapports (anonymes)\\
--desactiver,\allowbreak{} --disable & interdit l'envoi des rapports\\
\bottomrule
\end{tabularx}
\par\medskip
\begin{tabularx}{\linewidth}{@{}>{\ttfamily\small\raggedright}p{0.38\linewidth}X@{}}
\toprule
\multicolumn{2}{@{}p{\linewidth}@{}}{\sffamily\bfseries coupole maj (update)\newline\normalfont\small Interroge les versions publiées sur GitHub. Le paquet autonome se met à jour lui-même ; une installation par pip/pipx indique la commande à lancer.}\\
\midrule
--appliquer,\allowbreak{} --apply & télécharge et installe la mise à jour (paquet autonome)\\
\bottomrule
\end{tabularx}
\par\medskip
\begin{tabularx}{\linewidth}{@{}>{\ttfamily\small\raggedright}p{0.38\linewidth}X@{}}
\toprule
\multicolumn{2}{@{}p{\linewidth}@{}}{\sffamily\bfseries coupole sources\newline\normalfont\small Sans option : liste les adresses et leur origine (défaut livré, fichier distant, valeur forcée).}\\
\midrule
--forcer,\allowbreak{} --set CLÉ VALEUR & force une valeur (jamais écrasée par une mise à jour)\\
--defaut,\allowbreak{} --default CLÉ & revient à la valeur par défaut pour cette clé\\
--tout-reinitialiser,\allowbreak{} --reset-all & revient aux valeurs par défaut pour toutes les clés\\
--tester,\allowbreak{} --test CLÉ & teste la connexion à cette adresse\\
--distant,\allowbreak{} --remote & récupère maintenant le fichier de sources publié dans le dépôt\\
\bottomrule
\end{tabularx}
\par\medskip
\begin{tabularx}{\linewidth}{@{}>{\ttfamily\small\raggedright}p{0.38\linewidth}X@{}}
\toprule
\multicolumn{2}{@{}p{\linewidth}@{}}{\sffamily\bfseries coupole ohp inventaire (inventory)\newline\normalfont\small Inventaire de la banque par le service TAP public (http://tap-ufe.obspm.fr/tap, table ivoa.obscore), gardé en cache local.}\\
\midrule
--rafraichir,\allowbreak{} --refresh & réinterroge le service TAP de l'Observatoire (une requête)\\
--manquantes,\allowbreak{} --missing & liste ce qui reste à télécharger par rapport au dossier de sortie (--dest)\\
--dest DOSSIER & dossier de destination (défaut : \textasciitilde{}/Coupole/OHP\_DU\_ECU ou celui des réglages)\\
--json & sortie JSON (pour les scripts)\\
\bottomrule
\end{tabularx}
\par\medskip
\begin{tabularx}{\linewidth}{@{}>{\ttfamily\small\raggedright}p{0.38\linewidth}X@{}}
\toprule
\multicolumn{2}{@{}p{\linewidth}@{}}{\sffamily\bfseries coupole ohp catalogue (catalog)\newline\normalfont\small Les 194 noms de la base regroupés en objets, classés par type.}\\
\midrule
--type TYPE & catégorie : ast, neocp, com, pla, tno, pn, neb, amas, gal, eto (répétable)\\
--telescope \{T120,\allowbreak{}IRIS\} & télescope (répétable)\\
--chercher,\allowbreak{} --search TEXTE & filtre les objets dont le nom contient ce texte\\
--json & sortie JSON (pour les scripts)\\
\bottomrule
\end{tabularx}
\par\medskip
\begin{tabularx}{\linewidth}{@{}>{\ttfamily\small\raggedright}p{0.38\linewidth}X@{}}
\toprule
\multicolumn{2}{@{}p{\linewidth}@{}}{\sffamily\bfseries coupole ohp images (list)\newline\normalfont\small Une ligne par image : date, télescope, filtre, pose, drapeaux (doublon, date partagée, plein jour).}\\
\midrule
OBJET & objet(s) : « NGC 6888 », « Pluton », « 1998 AG6 », ou un nom de la base\\
--type TYPE & catégorie : ast, neocp, com, pla, tno, pn, neb, amas, gal, eto (répétable)\\
--telescope \{T120,\allowbreak{}IRIS\} & télescope (répétable)\\
--filtre,\allowbreak{} --filter F & filtre : B, V, R (Johnson et Cousins), Bc, Vc, Rc, Ha, OIII, g, r, i, z, ou le nom exact (« Johnson R ») ; répétable ou séparé par des virgules\\
--nuit,\allowbreak{} --night AAAA-MM-JJ & nuit (date du soir) ; répétable\\
--sans-dates-douteuses,\allowbreak{} --no-doubtful-dates & écarte les poses à DATE-OBS partagée ou datées en plein jour\\
--tout,\allowbreak{} --all & toute la banque (≈ 78 Go)\\
--csv FICHIER & écrit aussi la liste dans ce fichier CSV\\
--dest DOSSIER & dossier de destination (défaut : \textasciitilde{}/Coupole/OHP\_DU\_ECU ou celui des réglages)\\
\bottomrule
\end{tabularx}
\par\medskip
\begin{tabularx}{\linewidth}{@{}>{\ttfamily\small\raggedright}p{0.38\linewidth}X@{}}
\toprule
\multicolumn{2}{@{}p{\linewidth}@{}}{\sffamily\bfseries coupole ohp estimer (estimate)\newline\normalfont\small Estimation avant téléchargement, avec la place libre à destination.}\\
\midrule
OBJET & objet(s) : « NGC 6888 », « Pluton », « 1998 AG6 », ou un nom de la base\\
--type TYPE & catégorie : ast, neocp, com, pla, tno, pn, neb, amas, gal, eto (répétable)\\
--telescope \{T120,\allowbreak{}IRIS\} & télescope (répétable)\\
--filtre,\allowbreak{} --filter F & filtre : B, V, R (Johnson et Cousins), Bc, Vc, Rc, Ha, OIII, g, r, i, z, ou le nom exact (« Johnson R ») ; répétable ou séparé par des virgules\\
--nuit,\allowbreak{} --night AAAA-MM-JJ & nuit (date du soir) ; répétable\\
--sans-dates-douteuses,\allowbreak{} --no-doubtful-dates & écarte les poses à DATE-OBS partagée ou datées en plein jour\\
--tout,\allowbreak{} --all & toute la banque (≈ 78 Go)\\
--format \{xisf,\allowbreak{}xisf16,\allowbreak{}fz,\allowbreak{}fits\} & format de sortie : xisf (PixInsight, défaut), xisf16 (XISF compatible N.I.N.A. et Siril : entiers 16 bits, zlib), fz (FITS compressé sans perte), fits (float32)\\
--dest DOSSIER & dossier de destination (défaut : \textasciitilde{}/Coupole/OHP\_DU\_ECU ou celui des réglages)\\
\bottomrule
\end{tabularx}
\par\medskip
\begin{tabularx}{\linewidth}{@{}>{\ttfamily\small\raggedright}p{0.38\linewidth}X@{}}
\toprule
\multicolumn{2}{@{}p{\linewidth}@{}}{\sffamily\bfseries coupole ohp traiter (process)\newline\normalfont\small Traitement complet d'une sélection, reprenable : relancer la même commande reprend là où elle s'est arrêtée.}\\
\midrule
OBJET & objet(s) : « NGC 6888 », « Pluton », « 1998 AG6 », ou un nom de la base\\
--type TYPE & catégorie : ast, neocp, com, pla, tno, pn, neb, amas, gal, eto (répétable)\\
--telescope \{T120,\allowbreak{}IRIS\} & télescope (répétable)\\
--filtre,\allowbreak{} --filter F & filtre : B, V, R (Johnson et Cousins), Bc, Vc, Rc, Ha, OIII, g, r, i, z, ou le nom exact (« Johnson R ») ; répétable ou séparé par des virgules\\
--nuit,\allowbreak{} --night AAAA-MM-JJ & nuit (date du soir) ; répétable\\
--sans-dates-douteuses,\allowbreak{} --no-doubtful-dates & écarte les poses à DATE-OBS partagée ou datées en plein jour\\
--tout,\allowbreak{} --all & toute la banque (≈ 78 Go)\\
--dest DOSSIER & dossier de destination (défaut : \textasciitilde{}/Coupole/OHP\_DU\_ECU ou celui des réglages)\\
--format \{xisf,\allowbreak{}xisf16,\allowbreak{}fz,\allowbreak{}fits\} & format de sortie : xisf (PixInsight, défaut), xisf16 (XISF compatible N.I.N.A. et Siril : entiers 16 bits, zlib), fz (FITS compressé sans perte), fits (float32)\\
--noms,\allowbreak{} --names \{fr,\allowbreak{}en\} & langue des noms de dossiers, d'objets et des en-têtes (défaut : langue de l'interface)\\
--astap \{auto,\allowbreak{}tous,\allowbreak{}suspectes,\allowbreak{}jamais,\allowbreak{}all,\allowbreak{}suspicious,\allowbreak{}never\} & vérification ASTAP : auto/tous (chaque image si ASTAP est là), suspectes, jamais\\
--telechargements,\allowbreak{} --downloads N & téléchargements simultanés (défaut automatique ; 4 au plus)\\
--conversions N & conversions simultanées (défaut : selon cœurs et mémoire)\\
--econome,\allowbreak{} --eco & mode économe : 1 conversion, 2 téléchargements\\
--debit,\allowbreak{} --rate Mo/s & plafond de débit vers le serveur en Mo/s (défaut : réglage ; aucun plafond, 0 = illimité)\\
--garder-fits,\allowbreak{} --keep-fits & garde aussi les FITS d'origine (dans \_traitement/fits\_origine)\\
--garder-doublons,\allowbreak{} --keep-duplicates & traite aussi les doublons (inventaire et pixels identiques) au lieu de les écarter\\
--oui,\allowbreak{} --yes & accepte une sélection de plus de 5 Go sans confirmation\\
\bottomrule
\end{tabularx}
\par\medskip
\begin{tabularx}{\linewidth}{@{}>{\ttfamily\small\raggedright}p{0.38\linewidth}X@{}}
\toprule
\multicolumn{2}{@{}p{\linewidth}@{}}{\sffamily\bfseries coupole ohp tout (all)\newline\normalfont\small Toute la banque (≈ 8 000 images, ≈ 78 Go de FITS → ≈ 30 Go en XISF), avec estimation du volume et du temps avant confirmation ; relancer la même commande reprend là où elle s'était arrêtée.}\\
\midrule
--dest DOSSIER & dossier de destination (défaut : \textasciitilde{}/Coupole/OHP\_DU\_ECU ou celui des réglages)\\
--format \{xisf,\allowbreak{}xisf16,\allowbreak{}fz,\allowbreak{}fits\} & format de sortie : xisf (PixInsight, défaut), xisf16 (XISF compatible N.I.N.A. et Siril : entiers 16 bits, zlib), fz (FITS compressé sans perte), fits (float32)\\
--noms,\allowbreak{} --names \{fr,\allowbreak{}en\} & langue des noms de dossiers, d'objets et des en-têtes (défaut : langue de l'interface)\\
--astap \{auto,\allowbreak{}tous,\allowbreak{}suspectes,\allowbreak{}jamais,\allowbreak{}all,\allowbreak{}suspicious,\allowbreak{}never\} & vérification ASTAP : auto/tous (chaque image si ASTAP est là), suspectes, jamais\\
--telechargements,\allowbreak{} --downloads N & téléchargements simultanés (défaut automatique ; 4 au plus)\\
--conversions N & conversions simultanées (défaut : selon cœurs et mémoire)\\
--econome,\allowbreak{} --eco & mode économe : 1 conversion, 2 téléchargements\\
--debit,\allowbreak{} --rate Mo/s & plafond de débit vers le serveur en Mo/s (défaut : réglage ; aucun plafond, 0 = illimité)\\
--garder-fits,\allowbreak{} --keep-fits & garde aussi les FITS d'origine (dans \_traitement/fits\_origine)\\
--garder-doublons,\allowbreak{} --keep-duplicates & traite aussi les doublons (inventaire et pixels identiques) au lieu de les écarter\\
--oui,\allowbreak{} --yes & accepte une sélection de plus de 5 Go sans confirmation\\
\bottomrule
\end{tabularx}
\par\medskip
\begin{tabularx}{\linewidth}{@{}>{\ttfamily\small\raggedright}p{0.38\linewidth}X@{}}
\toprule
\multicolumn{2}{@{}p{\linewidth}@{}}{\sffamily\bfseries coupole ohp nouveautes (new)\newline\normalfont\small Réinterroge le service TAP, compare à la copie locale (dossier de sortie) et liste les nouveautés ; --telecharger les traite dans la même arborescence (jamais sans cette option).}\\
\midrule
--telecharger,\allowbreak{} --download & télécharge et convertit les nouveautés listées\\
--dest DOSSIER & dossier de destination (défaut : \textasciitilde{}/Coupole/OHP\_DU\_ECU ou celui des réglages)\\
--format \{xisf,\allowbreak{}xisf16,\allowbreak{}fz,\allowbreak{}fits\} & format de sortie : xisf (PixInsight, défaut), xisf16 (XISF compatible N.I.N.A. et Siril : entiers 16 bits, zlib), fz (FITS compressé sans perte), fits (float32)\\
--noms,\allowbreak{} --names \{fr,\allowbreak{}en\} & langue des noms de dossiers, d'objets et des en-têtes (défaut : langue de l'interface)\\
--astap \{auto,\allowbreak{}tous,\allowbreak{}suspectes,\allowbreak{}jamais,\allowbreak{}all,\allowbreak{}suspicious,\allowbreak{}never\} & vérification ASTAP : auto/tous (chaque image si ASTAP est là), suspectes, jamais\\
--telechargements,\allowbreak{} --downloads N & téléchargements simultanés (défaut automatique ; 4 au plus)\\
--conversions N & conversions simultanées (défaut : selon cœurs et mémoire)\\
--econome,\allowbreak{} --eco & mode économe : 1 conversion, 2 téléchargements\\
--debit,\allowbreak{} --rate Mo/s & plafond de débit vers le serveur en Mo/s (défaut : réglage ; aucun plafond, 0 = illimité)\\
--garder-fits,\allowbreak{} --keep-fits & garde aussi les FITS d'origine (dans \_traitement/fits\_origine)\\
--garder-doublons,\allowbreak{} --keep-duplicates & traite aussi les doublons (inventaire et pixels identiques) au lieu de les écarter\\
--oui,\allowbreak{} --yes & accepte une sélection de plus de 5 Go sans confirmation\\
\bottomrule
\end{tabularx}
\par\medskip
\begin{tabularx}{\linewidth}{@{}>{\ttfamily\small\raggedright}p{0.38\linewidth}X@{}}
\toprule
\multicolumn{2}{@{}p{\linewidth}@{}}{\sffamily\bfseries coupole ohp reorganiser (reorganise, reorganize)\newline\normalfont\small Parcourt un dossier, reconnaît les fichiers produits par Coupole (en-tête) et les déplace dans l'arborescence des lots de --dest ; conflit de nom → suffixe et ligne de journal ; rien n'est écrasé.}\\
\midrule
DOSSIER & dossier où se trouvent les fichiers à ranger\\
--dest DOSSIER & dossier de destination (défaut : \textasciitilde{}/Coupole/OHP\_DU\_ECU ou celui des réglages)\\
\bottomrule
\end{tabularx}
\par\medskip
\begin{tabularx}{\linewidth}{@{}>{\ttfamily\small\raggedright}p{0.38\linewidth}X@{}}
\toprule
\multicolumn{2}{@{}p{\linewidth}@{}}{\sffamily\bfseries coupole ohp telecharger (download)\newline\normalfont\small Rangement : <dossier>/<objet>/<nuit>\_<télescope>/<fichier d'origine>.}\\
\midrule
OBJET & objet(s) : « NGC 6888 », « Pluton », « 1998 AG6 », ou un nom de la base\\
--type TYPE & catégorie : ast, neocp, com, pla, tno, pn, neb, amas, gal, eto (répétable)\\
--telescope \{T120,\allowbreak{}IRIS\} & télescope (répétable)\\
--filtre,\allowbreak{} --filter F & filtre : B, V, R (Johnson et Cousins), Bc, Vc, Rc, Ha, OIII, g, r, i, z, ou le nom exact (« Johnson R ») ; répétable ou séparé par des virgules\\
--nuit,\allowbreak{} --night AAAA-MM-JJ & nuit (date du soir) ; répétable\\
--sans-dates-douteuses,\allowbreak{} --no-doubtful-dates & écarte les poses à DATE-OBS partagée ou datées en plein jour\\
--tout,\allowbreak{} --all & toute la banque (≈ 78 Go)\\
--dest DOSSIER & dossier de destination (défaut : \textasciitilde{}/Coupole/OHP\_DU\_ECU ou celui des réglages)\\
--telechargements,\allowbreak{} --downloads N & téléchargements simultanés (défaut automatique ; 4 au plus)\\
--debit,\allowbreak{} --rate Mo/s & plafond de débit vers le serveur en Mo/s (défaut : réglage ; aucun plafond, 0 = illimité)\\
--oui,\allowbreak{} --yes & accepte une sélection de plus de 5 Go sans confirmation\\
\bottomrule
\end{tabularx}
\par\medskip
\begin{tabularx}{\linewidth}{@{}>{\ttfamily\small\raggedright}p{0.38\linewidth}X@{}}
\toprule
\multicolumn{2}{@{}p{\linewidth}@{}}{\sffamily\bfseries coupole ohp anomalies (anomalies-report)\newline\normalfont\small Doublons (même fichier, copie diurne, pixels identiques), dates partagées ou diurnes, champ incohérent, classement à vérifier, instrument inconnu. Rien n'est supprimé.}\\
\midrule
--dest DOSSIER & ajoute les doublons de pixels trouvés dans ce dossier de traitement\\
--csv FICHIER & écrit aussi la liste dans ce fichier CSV\\
--nouveaux,\allowbreak{} --new & seulement les images nouvelles depuis le dernier rafraîchissement\\
\bottomrule
\end{tabularx}
\par\medskip
\begin{tabularx}{\linewidth}{@{}>{\ttfamily\small\raggedright}p{0.38\linewidth}X@{}}
\toprule
\multicolumn{2}{@{}p{\linewidth}@{}}{\sffamily\bfseries coupole ohp classer (classify)\newline\normalfont\small Sans argument : liste les noms à vérifier. Avec un nom de la base et --objet : mémorise la correction (fusion avec un objet existant si --type est omis). --en-ligne : interroge JPL SBDB et CDS Sesame.}\\
\midrule
NOM\_BASE & nom de cible tel qu'il figure dans la base\\
--objet,\allowbreak{} --object OBJET & objet canonique à associer (existant : fusion)\\
--type \{ast,\allowbreak{}neocp,\allowbreak{}com,\allowbreak{}pla,\allowbreak{}tno,\allowbreak{}pn,\allowbreak{}neb,\allowbreak{}amas,\allowbreak{}gal,\allowbreak{}eto,\allowbreak{}autre\} & catégorie : ast, neocp, com, pla, tno, pn, neb, amas, gal, eto (répétable)\\
--oublier,\allowbreak{} --forget & oublie la correction de ce nom\\
--en-ligne,\allowbreak{} --online & résout en ligne les noms inconnus (JPL SBDB, CDS Sesame), avec cache\\
\bottomrule
\end{tabularx}
\par\medskip
\begin{tabularx}{\linewidth}{@{}>{\ttfamily\small\raggedright}p{0.38\linewidth}X@{}}
\toprule
\multicolumn{2}{@{}p{\linewidth}@{}}{\sffamily\bfseries coupole ohp ranger (sort)\newline\normalfont\small Utile après une interruption, ou pour regrouper plusieurs traitements dans un même dossier.}\\
\midrule
--dest DOSSIER & dossier de destination (défaut : \textasciitilde{}/Coupole/OHP\_DU\_ECU ou celui des réglages)\\
\bottomrule
\end{tabularx}
\par\medskip
\begin{tabularx}{\linewidth}{@{}>{\ttfamily\small\raggedright}p{0.38\linewidth}X@{}}
\toprule
\multicolumn{2}{@{}p{\linewidth}@{}}{\sffamily\bfseries coupole ohp bilan (status)\newline\normalfont\small état d'avancement d'un dossier de traitement}\\
\midrule
--dest DOSSIER & dossier de destination (défaut : \textasciitilde{}/Coupole/OHP\_DU\_ECU ou celui des réglages)\\
--json & sortie JSON (pour les scripts)\\
\bottomrule
\end{tabularx}
\par\medskip
\begin{tabularx}{\linewidth}{@{}>{\ttfamily\small\raggedright}p{0.38\linewidth}X@{}}
\toprule
\multicolumn{2}{@{}p{\linewidth}@{}}{\sffamily\bfseries coupole ohp fusionner (merge)\newline\normalfont\small fusionner la base d'état du partage et la base de travail locale (après une divergence)}\\
\midrule
--dest DOSSIER & dossier de destination (défaut : \textasciitilde{}/Coupole/OHP\_DU\_ECU ou celui des réglages)\\
\bottomrule
\end{tabularx}
\par\medskip
\begin{tabularx}{\linewidth}{@{}>{\ttfamily\small\raggedright}p{0.38\linewidth}X@{}}
\toprule
\multicolumn{2}{@{}p{\linewidth}@{}}{\sffamily\bfseries coupole ohp reconnaitre (recognise, recognize)\newline\normalfont\small Lit l'en-tête de chaque image convertie du dossier de sortie (OHP:Source:URL ou HISTORY) et l'inscrit comme possédée dans la base de suivi.  Ne télécharge rien, ne déplace rien, ne réécrit aucun fichier.}\\
\midrule
--dest DOSSIER & dossier de destination (défaut : \textasciitilde{}/Coupole/OHP\_DU\_ECU ou celui des réglages)\\
\bottomrule
\end{tabularx}
\par\medskip
\begin{tabularx}{\linewidth}{@{}>{\ttfamily\small\raggedright}p{0.38\linewidth}X@{}}
\toprule
\multicolumn{2}{@{}p{\linewidth}@{}}{\sffamily\bfseries coupole ohp metadonnees (metadata)\newline\normalfont\small Pour chaque XISF/FITS du dossier : FOCALLEN accordé à l'échelle mesurée (ancienne valeur en HISTORY) et propriétés XISF Instrument:Sensor:XPixelSize, Instrument:Camera:XBinning, Instrument:Telescope:FocalLength… Les pixels ne sont pas touchés (relus et comparés), écriture atomique, journal dans \_traitement/metadonnees.csv. Sans --reecrire : simulation.}\\
\midrule
DOSSIER & dossier de sortie (ou un lot)\\
--reecrire,\allowbreak{} --rewrite & écrit vraiment (sinon : dit seulement ce qui changerait) ; note aussi dans la base d'état l'emplacement de chaque fichier, relatif au dossier de sortie (copie d'un ancien traitement, autre machine)\\
--noms,\allowbreak{} --names \{fr,\allowbreak{}en\} & langue des noms de dossiers, d'objets et des en-têtes (défaut : langue de l'interface)\\
\bottomrule
\end{tabularx}
\par\medskip
\begin{tabularx}{\linewidth}{@{}>{\ttfamily\small\raggedright}p{0.38\linewidth}X@{}}
\toprule
\multicolumn{2}{@{}p{\linewidth}@{}}{\sffamily\bfseries coupole archives liste (list)\newline\normalfont\small liste les archives connues (étape, missions, état)}\\
\midrule
--json & sortie JSON (pour les scripts)\\
\bottomrule
\end{tabularx}
\par\medskip
\begin{tabularx}{\linewidth}{@{}>{\ttfamily\small\raggedright}p{0.38\linewidth}X@{}}
\toprule
\multicolumn{2}{@{}p{\linewidth}@{}}{\sffamily\bfseries coupole archives chercher (search)\newline\normalfont\small Cherche dans MAST, l'ESO et IRSA (par défaut) les produits finaux publics autour d'un objet ; --archive opus ou pds avec --corps pour les sondes planétaires.}\\
\midrule
OBJET & objet (nom ou coordonnées)\\
--rayon,\allowbreak{} --radius ARCMIN & rayon de recherche en minutes d'arc (défaut : 3)\\
--archive A & archive à interroger (répétable) : mast, eso, irsa, noirlab, koa, sdss, goa, smoka, opus, pds\\
--mission M & mission (répétable) : JWST, HST, GALEX, ESO, Spitzer, WISE, 2MASS, NOIRLab, Keck, SDSS, Voyager, Cassini, Juno\\
--instrument I & instrument (morceau du nom, répétable)\\
--filtre,\allowbreak{} --filter F & filtre (nom exact, répétable)\\
--debut,\allowbreak{} --from AAAA-MM-JJ & observations à partir de cette date\\
--fin,\allowbreak{} --to AAAA-MM-JJ & observations jusqu'à cette date\\
--tous-niveaux,\allowbreak{} --all-levels & aussi les poses calibrées une à une et les données brutes\\
--non-publiques,\allowbreak{} --proprietary & aussi les données encore réservées (pour information : non téléchargeables)\\
--corps,\allowbreak{} --body CORPS & corps du Système solaire (Voyager, Cassini, Juno)\\
--limite,\allowbreak{} --limit N & nombre maximal de lignes par archive (défaut : 2000)\\
--id ID & garder seulement cet identifiant (répétable)\\
--variantes,\allowbreak{} --variants & montrer aussi les variantes (autres produits des mêmes poses) et les images de détection ; par défaut, un produit par filtre\\
--programme,\allowbreak{} --program N & ne garder que les produits de ce programme (répétable), par exemple 9978\\
--json & sortie JSON (pour les scripts)\\
--csv FICHIER & écrit aussi les résultats dans ce fichier CSV\\
\bottomrule
\end{tabularx}
\par\medskip
\begin{tabularx}{\linewidth}{@{}>{\ttfamily\small\raggedright}p{0.38\linewidth}X@{}}
\toprule
\multicolumn{2}{@{}p{\linewidth}@{}}{\sffamily\bfseries coupole archives estimer (estimate)\newline\normalfont\small Comme « chercher », puis additionne les tailles (annoncées ou mesurées) sans rien télécharger.}\\
\midrule
OBJET & objet (nom ou coordonnées)\\
--rayon,\allowbreak{} --radius ARCMIN & rayon de recherche en minutes d'arc (défaut : 3)\\
--archive A & archive à interroger (répétable) : mast, eso, irsa, noirlab, koa, sdss, goa, smoka, opus, pds\\
--mission M & mission (répétable) : JWST, HST, GALEX, ESO, Spitzer, WISE, 2MASS, NOIRLab, Keck, SDSS, Voyager, Cassini, Juno\\
--instrument I & instrument (morceau du nom, répétable)\\
--filtre,\allowbreak{} --filter F & filtre (nom exact, répétable)\\
--debut,\allowbreak{} --from AAAA-MM-JJ & observations à partir de cette date\\
--fin,\allowbreak{} --to AAAA-MM-JJ & observations jusqu'à cette date\\
--tous-niveaux,\allowbreak{} --all-levels & aussi les poses calibrées une à une et les données brutes\\
--non-publiques,\allowbreak{} --proprietary & aussi les données encore réservées (pour information : non téléchargeables)\\
--corps,\allowbreak{} --body CORPS & corps du Système solaire (Voyager, Cassini, Juno)\\
--limite,\allowbreak{} --limit N & nombre maximal de lignes par archive (défaut : 2000)\\
--id ID & garder seulement cet identifiant (répétable)\\
--variantes,\allowbreak{} --variants & montrer aussi les variantes (autres produits des mêmes poses) et les images de détection ; par défaut, un produit par filtre\\
--programme,\allowbreak{} --program N & ne garder que les produits de ce programme (répétable), par exemple 9978\\
\bottomrule
\end{tabularx}
\par\medskip
\begin{tabularx}{\linewidth}{@{}>{\ttfamily\small\raggedright}p{0.38\linewidth}X@{}}
\toprule
\multicolumn{2}{@{}p{\linewidth}@{}}{\sffamily\bfseries coupole archives telecharger (download)\newline\normalfont\small Comme « chercher », puis télécharge (confirmation au-delà du seuil : --oui), extrait l'image scientifique et la range dans <dest>/Archives.}\\
\midrule
OBJET & objet (nom ou coordonnées)\\
--rayon,\allowbreak{} --radius ARCMIN & rayon de recherche en minutes d'arc (défaut : 3)\\
--archive A & archive à interroger (répétable) : mast, eso, irsa, noirlab, koa, sdss, goa, smoka, opus, pds\\
--mission M & mission (répétable) : JWST, HST, GALEX, ESO, Spitzer, WISE, 2MASS, NOIRLab, Keck, SDSS, Voyager, Cassini, Juno\\
--instrument I & instrument (morceau du nom, répétable)\\
--filtre,\allowbreak{} --filter F & filtre (nom exact, répétable)\\
--debut,\allowbreak{} --from AAAA-MM-JJ & observations à partir de cette date\\
--fin,\allowbreak{} --to AAAA-MM-JJ & observations jusqu'à cette date\\
--tous-niveaux,\allowbreak{} --all-levels & aussi les poses calibrées une à une et les données brutes\\
--non-publiques,\allowbreak{} --proprietary & aussi les données encore réservées (pour information : non téléchargeables)\\
--corps,\allowbreak{} --body CORPS & corps du Système solaire (Voyager, Cassini, Juno)\\
--limite,\allowbreak{} --limit N & nombre maximal de lignes par archive (défaut : 2000)\\
--id ID & garder seulement cet identifiant (répétable)\\
--variantes,\allowbreak{} --variants & montrer aussi les variantes (autres produits des mêmes poses) et les images de détection ; par défaut, un produit par filtre\\
--programme,\allowbreak{} --program N & ne garder que les produits de ce programme (répétable), par exemple 9978\\
--dest DOSSIER & dossier de sortie (défaut : celui des Préférences)\\
--format \{xisf,\allowbreak{}fits\} & format des images préparées (défaut : xisf)\\
--sans-original,\allowbreak{} --no-original & (comportement par défaut depuis 0.2.2) n'efface que le fichier d'origine une fois l'image préparée\\
--garder-original,\allowbreak{} --keep-original & garder aussi le fichier d'origine (dans <champ>/\_originaux) ; par défaut, seul le produit préparé reste\\
--masque,\allowbreak{} --mask & écrire aussi le masque des pixels sans donnée (compressé, après recadrage)\\
--debit,\allowbreak{} --rate MO/S & plafond de débit en Mo/s (défaut : réglage ; aucun plafond, 0 = illimité)\\
--paralleles,\allowbreak{} --parallel N & téléchargements simultanés (défaut : 2, au plus 6)\\
--max N & au plus N fichiers (les plus proches)\\
--oui,\allowbreak{} --yes & confirme un téléchargement au-delà du seuil\\
\bottomrule
\end{tabularx}
\par\medskip
\begin{tabularx}{\linewidth}{@{}>{\ttfamily\small\raggedright}p{0.38\linewidth}X@{}}
\toprule
\multicolumn{2}{@{}p{\linewidth}@{}}{\sffamily\bfseries coupole archives preparer (prepare)\newline\normalfont\small Extrait l'image scientifique d'un FITS à extensions, ou convertit un produit planétaire, en XISF ou FITS Float32 à côté du fichier.}\\
\midrule
FICHIER & fichiers à préparer\\
--format \{xisf,\allowbreak{}fits\} & format des images préparées (défaut : xisf)\\
--etiquette,\allowbreak{} --label FICHIER & label PDS3 détaché (.LBL) ou PDS4 (.xml)\\
\bottomrule
\end{tabularx}
\par\medskip
\begin{tabularx}{\linewidth}{@{}>{\ttfamily\small\raggedright}p{0.38\linewidth}X@{}}
\toprule
\multicolumn{2}{@{}p{\linewidth}@{}}{\sffamily\bfseries coupole archives aligner (align)\newline\normalfont\small Rééchantillonne les images (WCS) sur la grille de l'une d'elles ou une grille optimale ; écrit les fichiers alignés, le masque commun, l'aperçu couleur et composition.txt.}\\
\midrule
FICHIER & images préparées (FITS ou XISF avec WCS)\\
--dest DOSSIER & dossier des fichiers alignés\\
--reference N & indice de l'image de référence (défaut : 0, la première)\\
--optimale,\allowbreak{} --optimal & grille optimale couvrant toutes les images\\
--methode,\allowbreak{} --method \{bilineaire,\allowbreak{}adaptative,\allowbreak{}exacte,\allowbreak{}proche\} & rééchantillonnage : bilineaire, adaptative, exacte (reproject), proche\\
--echelle,\allowbreak{} --scale X & facteur de taille des pixels de sortie (défaut : 1)\\
--sans-recadrage,\allowbreak{} --no-crop & ne recadre pas sur la zone commune\\
--format \{xisf,\allowbreak{}fits\} & format des images préparées (défaut : xisf)\\
\bottomrule
\end{tabularx}
\par\medskip
\begin{tabularx}{\linewidth}{@{}>{\ttfamily\small\raggedright}p{0.38\linewidth}X@{}}
\toprule
\multicolumn{2}{@{}p{\linewidth}@{}}{\sffamily\bfseries coupole archives catalogue (catalog)\newline\normalfont\small Liste les cibles d'une archive, les plus observées d'abord (requête agrégée, cache de 30 jours), ou les programmes célèbres.  Listes : celebres, hst, jwst, eso, voyager, cassini, juno.}\\
\midrule
liste \{celebres,\allowbreak{}famous,\allowbreak{}hst,\allowbreak{}jwst,\allowbreak{}eso,\allowbreak{}voyager,\allowbreak{}cassini,\allowbreak{}juno\} & liste à afficher (celebres par défaut)\\
--type T & type d'objet : galaxie, amas\_galaxies, amas\_stellaire, nebuleuse, etoile, systeme\_solaire, champ, autre\\
--filtre,\allowbreak{} --filter OBJET & ne garder que les cibles dont le nom contient ce texte\\
--page N & page à afficher (100 cibles par page)\\
--actualiser,\allowbreak{} --refresh & relire la liste en ligne (ignorer le cache)\\
--json & sortie JSON (pour les scripts)\\
\bottomrule
\end{tabularx}
\par\medskip
\begin{tabularx}{\linewidth}{@{}>{\ttfamily\small\raggedright}p{0.38\linewidth}X@{}}
\toprule
\multicolumn{2}{@{}p{\linewidth}@{}}{\sffamily\bfseries coupole archives reorganiser (reorganise)\newline\normalfont\small Renomme les fichiers déjà téléchargés (<cible>\_<instrument>\_<filtre>\_<date>\_<type>\_<n>x<pose>s), les range par champ, les recadre sur leur zone utile, supprime les masques pleine taille et met variantes et images de détection dans « \_variantes ».  Jamais d'écrasement ; journal dans \_etat/JOURNAL.txt.}\\
\midrule
--dest DOSSIER & dossier de sortie (défaut : celui des Préférences)\\
--simuler,\allowbreak{} --dry-run & montrer ce qui serait fait, sans rien changer\\
--garder-masques,\allowbreak{} --keep-masks & récrire les masques (compressés) au lieu de les supprimer\\
--sans-recadrage,\allowbreak{} --no-crop & ne pas recadrer les images\\
\bottomrule
\end{tabularx}
\par\medskip
\begin{tabularx}{\linewidth}{@{}>{\ttfamily\small\raggedright}p{0.38\linewidth}X@{}}
\toprule
\multicolumn{2}{@{}p{\linewidth}@{}}{\sffamily\bfseries coupole archives instruments\newline\normalfont\small topo des instruments et des filtres (ou la fiche d’un filtre)}\\
\midrule
FILTRE & nom d'un filtre (F656N, F444W, Ks…) : longueur d'onde, largeur, usage, équivalent amateur\\
--instrument I & instrument (morceau du nom, répétable)\\
\bottomrule
\end{tabularx}
\par\medskip
\begin{tabularx}{\linewidth}{@{}>{\ttfamily\small\raggedright}p{0.38\linewidth}X@{}}
\toprule
\multicolumn{2}{@{}p{\linewidth}@{}}{\sffamily\bfseries coupole archives hudf\newline\normalfont\small Liste les produits de haut niveau du HUDF (XDF, HUDF09, HUDF12) ; --telecharger récupère le préréglage « couleur en 3 filtres » (XDF F435W, F775W, F160W, 3 × 110 Mo) ou les filtres donnés par --filtre.}\\
\midrule
--telecharger,\allowbreak{} --download & télécharger et préparer (sinon : liste seulement)\\
--filtre,\allowbreak{} --filter F & filtre XDF à télécharger (répétable)\\
--dest DOSSIER & dossier de sortie (défaut : celui des Préférences)\\
--format \{xisf,\allowbreak{}fits\} & format des images préparées (défaut : xisf)\\
--oui,\allowbreak{} --yes & confirme un téléchargement au-delà du seuil\\
\bottomrule
\end{tabularx}
\par\medskip
\begin{tabularx}{\linewidth}{@{}>{\ttfamily\small\raggedright}p{0.38\linewidth}X@{}}
\toprule
\multicolumn{2}{@{}p{\linewidth}@{}}{\sffamily\bfseries coupole archives bilan (status)\newline\normalfont\small ce qui est déjà téléchargé et préparé}\\
\midrule
--dest DOSSIER & dossier de sortie (défaut : celui des Préférences)\\
--json & sortie JSON (pour les scripts)\\
\bottomrule
\end{tabularx}
\par\medskip
\begin{tabularx}{\linewidth}{@{}>{\ttfamily\small\raggedright}p{0.38\linewidth}X@{}}
\toprule
\multicolumn{2}{@{}p{\linewidth}@{}}{\sffamily\bfseries coupole qualite\newline\normalfont\small Facultatif : fond, bruit, FWHM et ellipticité (carte 3×3), gradient du fond, saturation, traînées, échantillonnage}\\
\midrule
DOSSIER & dossier de lots (ou un fichier image)\\
--ecrire,\allowbreak{} --write & écrit QUALITE.csv et QUALITE.txt dans chaque lot\\
--echantillon,\allowbreak{} --sample N & mesure N images par lot (réparties dans le temps) au lieu de toutes\\
--tout,\allowbreak{} --all & tout mesurer même au-delà de 200 images (sinon : échantillon de 5 par lot, annoncé)\\
--processus,\allowbreak{} --processes N & nombre de processus de mesure (défaut : plan machine)\\
--avec-calibration,\allowbreak{} --with-calibration & mesurer aussi flats, darks, bias, masters et fichiers intermédiaires (exclus par défaut)\\
\bottomrule
\end{tabularx}
\par\medskip
\begin{tabularx}{\linewidth}{@{}>{\ttfamily\small\raggedright}p{0.38\linewidth}X@{}}
\toprule
\multicolumn{2}{@{}p{\linewidth}@{}}{\sffamily\bfseries coupole donnees lire (read)\newline\normalfont\small décrit le contenu d'un fichier (FITS, CSV)}\\
\midrule
FICHIER & fichier à lire\\
--json & sortie JSON (pour les scripts)\\
\bottomrule
\end{tabularx}
\par\medskip
\begin{tabularx}{\linewidth}{@{}>{\ttfamily\small\raggedright}p{0.38\linewidth}X@{}}
\toprule
\multicolumn{2}{@{}p{\linewidth}@{}}{\sffamily\bfseries coupole donnees exporter (export)\newline\normalfont\small exporte un spectre ou une série en CSV}\\
\midrule
FICHIER & fichier à lire\\
--csv FICHIER & exporte un spectre ou une série en CSV\\
--hdu HDU & numéro de la donnée dans le fichier (défaut : la première 1D)\\
--vitesse,\allowbreak{} --velocity & ajoute la vitesse radio (spectre en fréquence)\\
--f0 MHz & fréquence de repos en MHz (défaut : RESTFRQ, sinon H I)\\
\bottomrule
\end{tabularx}
\par\medskip
\begin{tabularx}{\linewidth}{@{}>{\ttfamily\small\raggedright}p{0.38\linewidth}X@{}}
\toprule
\multicolumn{2}{@{}p{\linewidth}@{}}{\sffamily\bfseries coupole cosmo\newline\normalfont\small Du redshift aux distances : comobile, luminosité, angulaire, temps de regard en arrière, âge, volume, module de distance, échelle (Planck 2018 et autres jeux de paramètres)}\\
\midrule
Z & redshift(s) à calculer\\
--objet,\allowbreak{} --object NOM & nom d'objet : redshift demandé à SIMBAD\\
--modele,\allowbreak{} --model \{planck18,\allowbreak{}planck15,\allowbreak{}wmap9,\allowbreak{}simple,\allowbreak{}perso\} & jeu de paramètres (défaut : planck18)\\
--h0 H0 & H₀ en km s⁻¹ Mpc⁻¹ (modèle perso)\\
--om OM & Ωm total (modèle perso)\\
--ok OK & Ωk, courbure (planck18 ou perso)\\
--shoes & ajoute la comparaison SH0ES (planck18)\\
--sans-incertitudes,\allowbreak{} --no-uncertainties & sans les incertitudes (plus rapide)\\
--csv FICHIER & écrit les résultats dans ce fichier CSV\\
--courbes,\allowbreak{} --curves ZMIN ZMAX N & écrit les courbes (z de ZMIN à ZMAX, N points) dans le fichier CSV donné par --csv\\
--json & sortie JSON (pour les scripts)\\
\bottomrule
\end{tabularx}
\par\medskip
\begin{tabularx}{\linewidth}{@{}>{\ttfamily\small\raggedright}p{0.38\linewidth}X@{}}
\toprule
\multicolumn{2}{@{}p{\linewidth}@{}}{\sffamily\bfseries coupole anim liste}\\
\midrule
--faites & seulement les animations disponibles\\
--json & sortie JSON (pour les scripts)\\
\bottomrule
\end{tabularx}
\par\medskip
\begin{tabularx}{\linewidth}{@{}>{\ttfamily\small\raggedright}p{0.38\linewidth}X@{}}
\toprule
\multicolumn{2}{@{}p{\linewidth}@{}}{\sffamily\bfseries coupole anim info}\\
\midrule
ID & identifiant ClassAction (ex. sunmotions)\\
--json & sortie JSON (pour les scripts)\\
\bottomrule
\end{tabularx}
\par\medskip
\begin{tabularx}{\linewidth}{@{}>{\ttfamily\small\raggedright}p{0.38\linewidth}X@{}}
\toprule
\multicolumn{2}{@{}p{\linewidth}@{}}{\sffamily\bfseries coupole sites\newline\normalfont\small Sites d'observation (carte OpenStreetMap), fuseaux horaires, heure UTC et heure locale}\\
\midrule
--ajouter,\allowbreak{} --add X & ajoute un site : ID NOM LAT LON ALT FUSEAU [MPC]\\
--supprimer,\allowbreak{} --delete ID & supprime un site ajouté\\
--heure,\allowbreak{} --time DATE\_ISO & affiche une date UTC (ISO) en heure du site et date du soir\\
--site ID & identifiant du site (défaut : ohp)\\
\bottomrule
\end{tabularx}
\par\medskip
\begin{tabularx}{\linewidth}{@{}>{\ttfamily\small\raggedright}p{0.38\linewidth}X@{}}
\toprule
\multicolumn{2}{@{}p{\linewidth}@{}}{\sffamily\bfseries coupole machine\newline\normalfont\small Diagnostic matériel, parallélisme retenu, carte graphique, ASTAP}\\
\midrule
--json & sortie JSON (pour les scripts)\\
\bottomrule
\end{tabularx}
\par\medskip
}

% =====================================================================
\section{Évolutivité}\label{chap:evolutif}
\paragraph{Sources.} Les adresses (TAP, collection, portail, réécriture des URL de téléchargement, résolveurs,
point de collecte, dépôt, tuiles de carte) sont dans \texttt{coupole/donnees/sources.json}. Trois couches : valeurs
livrées ; fichier publié dans le dépôt, récupéré au démarrage s'il est plus récent et passe le contrôle de forme
(clés connues, types, domaines d'une liste blanche ; pas de signature cryptographique) ; valeurs forcées par
l'utilisateur, jamais écrasées. Préférences > Sources : modifier, revenir au défaut, tester. Si le serveur des
images déménage, \texttt{ohp.telechargement.de}/\texttt{vers} réécrit les URL.
\paragraph{Classement.} Table livrée \texttt{data/classement.json} (194 noms vérifiés) ; alias publiés ; corrections
de l'utilisateur (\emph{Corriger le classement}, \cmd{coupole ohp classer}) ; règles sur la forme du nom ; résolution
en ligne facultative (JPL SBDB, CDS Sesame), en cache ; regroupement par position sur un objet fixe connu ; sinon
« non classé ». Tout classement automatique est marqué « à vérifier ».
\paragraph{Base.} Colonnes ObsCore en plus ignorées, en moins remplacées par une valeur par défaut, instruments et
filtres inconnus rangés à part, données non images jamais signalées « de jour » ; testé sur inventaires simulés.
\paragraph{Sécurité.} Aucun secret dans le dépôt ; fichier de sources distant et point de collecte en HTTPS
obligatoire, taille bornée, liste blanche de domaines ; archive de mise à jour téléchargée seulement depuis GitHub,
contrôlée avant extraction (chemins absolus ou « .. » refusés, liens symboliques refusés, taille dépliée et nombre de
membres bornés), remplacement avec retour arrière en cas d'échec ; XISF et FITS lus avec des bornes de taille ; noms
de dossiers construits depuis la base rendus sûrs (ASCII, caractères interdits, noms réservés de Windows, longueur) ;
ASTAP lancé sans interpréteur de commandes ; aucun \texttt{eval} ni \texttt{pickle} sur des données externes.

% =====================================================================
\section{Rapports d'incident et confidentialité}
Au premier lancement, Coupole demande s'il peut envoyer des rapports anonymes : installation, preuve de vie (une
fois par jour), plantage, crash natif, interface figée, signalement écrit. Contenu : version, système, processeur,
mémoire, carte graphique, identifiant d'installation tiré au hasard, trace d'erreur avec le dossier personnel
remplacé par « \textasciitilde », le nom d'utilisateur et le nom de la machine masqués et les chemins profonds
tronqués ; un rapport ne dépasse jamais 64 ko. Jamais de nom de machine, d'utilisateur, ni d'image. Sans accord, rien ne part : les rapports restent dans \texttt{\_rapports} du
dossier de réglages. Choix modifiable (menu Aide, \cmd{coupole rapports --activer|--desactiver}).

% =====================================================================
\section{Écrire un module}
Un module est un paquet Python avec un \texttt{manifest.json} (identifiant, version, nom en français et en anglais,
icône, textes, points d'entrée graphique et ligne de commande), découvert dans \texttt{coupole/modules/}, par point
d'entrée \texttt{coupole.modules}, ou dans le dossier \texttt{modules/} des réglages. Textes par \texttt{tr()},
info-bulles obligatoires, travail long hors du fil graphique : voir CONTRIBUTING.md.

% =====================================================================
\section{Dépannage}
\begin{center}\small
\begin{tabularx}{\linewidth}{@{}>{\raggedright}p{0.34\linewidth}X@{}}
\toprule
\textbf{Symptôme} & \textbf{Que faire}\\
\midrule
Rien ne démarre (Windows) & menu Démarrer > Coupole > Diagnostic : l'erreur reste affichée.\\
« Interface graphique indisponible » & bibliothèques Qt manquantes (Linux : \texttt{libxcb-cursor0}, \texttt{libegl1}) ; la ligne de commande reste utilisable.\\
Téléchargement en échec & relancer (reprise) ; vérifier la connexion ; Préférences > Sources > Tester ; serveur déménagé : réécriture d'URL.\\
« ASTAP non trouvé » & assistant ASTAP ; ou \texttt{coupole astap --definir CHEMIN --catalogue DOSSIER}.\\
Catalogue incomplet & réinstaller le catalogue (1~476 tuiles attendues).\\
Place insuffisante & choisir un autre dossier de sortie ; la place nécessaire inclut les FITS en attente.\\
Machine lente & mode économe, ou moins de conversions simultanées.\\
Nom de cible inconnu & onglet Catalogue, case « À vérifier », \emph{Corriger le classement}.\\
Qualité indisponible & \texttt{pip install sep} (inclus dans les paquets).\\
Réglages revenus par défaut & le fichier \texttt{reglages.json} était illisible : il est gardé à côté avec le suffixe \texttt{.corrompu-…}.\\
Dossier de sortie non inscriptible & message clair au lancement ; choisir un autre dossier (Préférences ou champ Dossier).\\
Dialogue de fichiers sans le NAS, ou générique & Linux : installer \texttt{xdg-desktop-portal} et le moteur du bureau (\texttt{-kde}, \texttt{-gnome} ou \texttt{-gtk}), relancer ; KDE : \texttt{kio-fuse} ; sinon taper le chemin du partage monté dans le champ (section~\ref{sec:dialogues}).\\
Dialogue du système qui ne s'ouvre pas & 0.1.10 et avant, Spectres et séries sous KDE : mettre à jour (filtre sans nom refusé par le portail) ; sinon Préférences > Boîtes de dialogue de fichiers : Qt ; chaque ouverture est notée dans \texttt{coupole.log}.\\
Images possédées mais rien à ouvrir (copie d'un ancien traitement) & mettre à jour (0.1.11) : les emplacements sont retrouvés par \texttt{journal.csv} (section~\ref{sec:ouvrir}).\\
Qualité : colonnes vides sur ses propres poses & mettre à jour (0.1.11) : lecteur XISF tolérant ; la colonne \emph{état} donne le motif.\\
\texttt{\textasciitilde/.local/bin/coupole} introuvable & ajouter \texttt{\textasciitilde/.local/bin} au \texttt{PATH} (ligne donnée par l'installeur).\\
Fichier \texttt{.part} ou \texttt{.tmp} dans \texttt{\_traitement} & reste d'une coupure : la reprise le réutilise (\texttt{.part}) ou le remplace ; rien à faire.\\
WBPP ne trouve pas focale et pixel & saisir les valeurs de l'encadré \emph{Pour PixInsight / N.I.N.A.} (onglet Lots) dans le panneau \emph{Astrometric solution} ; fichiers anciens : \cmd{coupole ohp metadonnees DOSSIER --reecrire}.\\
Image blanche dans N.I.N.A. ou Siril & format XISF compatible (options du téléchargement), section~\ref{sec:compat}.\\
« database is locked » vers un NAS (0.1.9 et avant) & mettre à jour : la base de travail locale règle le problème (section~\ref{sec:partage}).\\
« Deux versions de la base d'état » & deux ordinateurs (ou le NAS) ont traité vers le même dossier : accepter la fusion, ou \cmd{coupole ohp fusionner --dest DOSSIER} ; les deux bases d'origine sont gardées.\\
Tout « à télécharger » sous Windows alors que le dossier (lecteur réseau) est complet & 0.2.0 et avant : mettre à jour (section~\ref{sec:partage}) ; la ligne « Dossier de suivi » dit ce qui est lu.\\
Archives : image noire sauf un bout, masque de 300~Mo, 17 fichiers pour 3 filtres (0.2.1) & mettre à jour puis \emph{Réorganiser les téléchargements} : recadrage, masques supprimés, variantes à part (section~\ref{sec:archives}).\\
Archives : « jupiter », « saturne » introuvables, rien pour Cassini, pas d'aperçu (0.2.1) & mettre à jour : corps cherchés par nom de cible, limite partagée, vignettes lues en entier.\\
Archives : image noire à l'ouverture & normal : électrons/s ou MJy/sr ; étirer (STF, auto-ajustement).\\
« La base de suivi de ce dossier n'a pas pu être lue » & la cause suit le message (et \texttt{coupole.log}) : base abîmée, partage débranché, droits ; ne pas relancer un téléchargement complet.\\
Pas de \texttt{\_traitement/} mais des images rangées & \emph{Reconnaître les fichiers existants} (bandeau), ou \cmd{coupole ohp reconnaitre --dest DOSSIER}.\\
« Recopie de la base d'état sur le partage impossible » & partage débranché ou base ouverte ailleurs : rien n'est perdu, nouvel essai au passage suivant et au prochain lancement.\\
\bottomrule
\end{tabularx}
\end{center}

\section*{Sources}
\addcontentsline{toc}{section}{Sources}
Banque « OHP student observations », Observatoire de Paris / PADC (Paris Astronomical Data Centre), service TAP
public, licence Etalab 2.0 ; guide officiel de la base (PADC-BDD-DU-ECU). Spécification XISF 1.0 (Pleiades
Astrophoto). ASTAP : \url{https://www.hnsky.org/astap.htm}. Codes d'observatoire : Minor Planet Center. SEP :
Barbary (2016), Bertin \& Arnouts (1996). Moments adaptatifs : Bernstein \& Jarvis (2002). Fond de carte :
© OpenStreetMap contributors. SIMBAD et Sesame : CDS, Strasbourg (Wenger et al. 2000) ; Aladin Lite (CDS) ;
NED (NASA/IPAC) ; JPL Small-Body Database (NASA/JPL). Cosmologie : Planck Collaboration (2020, A\&A 641, A6) ;
Hinshaw et al. (2013, WMAP 9) ; Riess et al. (2022, SH0ES) ; Hogg (1999, astro-ph/9905116) ; Komatsu et al. (2011)
pour la densité des neutrinos massifs ; astropy.cosmology ; calculateur « cosmologie-redshift » (même auteur). Archives : MAST (STScI), ESO Science Archive Facility, IRSA (NASA/IPAC), Astro Data Archive de NOIRLab, Keck Observatory Archive (NExScI), SDSS, Planetary Data System (OPUS du nœud Anneaux et Lunes, Imaging Node) ; reproject (Astropy) ; Rector et al. (2007, AJ 133, 598) pour l'ordre chromatique ; pratiques des amateurs : voir \texttt{docs/archives\_methode.md}. Animations : ClassAction, University of Nebraska-Lincoln Astronomy Education (Apache License 2.0) ; Meeus (1998, \emph{Astronomical Algorithms}, 2\textsuperscript{e} éd.) ; obliquité P03 / IAU 2006 ; limites des constellations de l'UAI (données livrées avec astropy) ; côtes Natural Earth 1:50m (domaine public) ; police STIX Two Math (SIL OFL 1.1) ; NASA Exoplanet Archive, NASA Planetary Fact Sheet, éphéméride JPL DE440 et éléments de Standish, AME 2020, filtres de Landolt (1992), spectre de M32 (Santos et al. 2002), Catalog and Atlas of Eclipsing Binaries, courbe de rotation : modèle de McMillan (2017) par galpy, mesures d'Eilers et al. (2019) et compilation galkin (Pato \& Iocco 2017), bras de Reid et al. (2019), barre de Wegg et al. (2015) ; atmosphère : Kasten \& Young (1989), Pickering (2002), Hansen \& Travis (1974), Ångström (1929), Ciddor (1996), modèle de ciel de l'ESO (Noll et al. 2012 ; Jones et al. 2013) ; Partages du DU et Calculateurs de TD : sujets « Parallaxe de l'ISS » et « Parallaxe lunaire » du DU (Observatoire de Paris) ; JPL DE440 (Park et al. 2021) lue par jplephem (MIT) ; ERFA (BSD, dérivé de SOFA, UAI) ; SGP4 (Vallado et al. 2006, paquet sgp4, MIT) et TLE de CelesTrak ; WGS84 (NIMA TR8350.2) ; Vincenty (1975) ; Kåsa (1976) ; Bennett (1982) et Saemundsson (1986) pour la réfraction ; skyfield (Rhodes, MIT) pour les références ; astrometry.net et nova.astrometry.net (Lang et al. 2010, AJ 139, 1782) ; WebDAV (RFC 4918) et partages publics Nextcloud ; Office Open XML (ECMA-376) ; machine d'Anticythère : adaptée d'Anticythere3D (même auteur), d'après Freeth et al. (\emph{Nature} 444, 2006 ; \emph{Nature} 454, 2008 ; \emph{Scientific Reports} 11:5821, 2021) ; spectres physiques : CODATA 2022, NIST ASD, Mamajek (2022), Malkov et al. (2020), van Belle et al. (2021), Martins et al. (2005), Castelli \& Kurucz (2003, ATLAS9), NIST-JANAF, Allard et al. (2000), CIE 1931 ; Gravité et orbites : d'après la simulation PhET « Gravity and Orbits » (PhET Interactive Simulations, University of Colorado Boulder ; rien n'en est repris), CODATA 2018, UAI 2015 (B3) et 2012 (B2), DE440 (Park et al. 2021), NASA (fiches Terre, Lune, ISS), Yoshida (1990) ; Orion en 3D et fiches au survol : SIMBAD (CDS, Strasbourg), Gaia DR3 (ESA/DPAC), Hipparcos (van Leeuwen 2007), Zucker et al. (2021, CC0), Großschedl et al. (2018), Kounkel et al. (2017, 2018), Foley et al. (2023), O'Dell et al. (2009), Stellarium (figure, CC BY-SA 4.0) ; images ESO/G. Beccari, ESO, et NASA, ESA, M. Robberto (STScI/ESA) and the Hubble Space Telescope Orion Treasury Project Team, NASA, ESA, and the Hubble Heritage Team (AURA/STScI), sous licence CC BY 4.0 (détail dans \texttt{NOTICE}).
\end{document}
